% Leçon 33 — Expression orale : économie — Chinois 1ère A — Cours signé OnBuch+
% Programme officiel MINESEC. Compiler avec : tectonic ZHA33-oral-economie.tex
\newcommand{\DOCMATIERE}{Chinois}
\newcommand{\DOCNIVEAU}{1ère A}
\newcommand{\DOCMODULE}{Module 4 : Activités économiques et monde du travail}
\newcommand{\DOCLECON}{ZHA33}
\newcommand{\DOCTITRE}{Leçon 33 — Expression orale : économie}
% OnBuch+ — préambule « cours signés » ENRICHI (compilé avec tectonic / XeLaTeX)
% Système visuel « Pop » d'OnBuch+ : fond crème (jamais blanc), bordures encre
% épaisses, ombres « dures » décalées, titres Archivo Black en capitales.
% Version MATHÉMATIQUES : graphes (pgfplots/tikz), boîtes pédagogiques
% (définition, propriété, méthode, attention, le savais-tu, exercice résolu,
% exercice, TP, bilan), page de garde et signature OnBuch+ ancrée en bas.
%
% Macros attendues dans main.tex AVANT \input{preamble} :
%   \DOCMATIERE  \DOCNIVEAU  \DOCMODULE  \DOCLECON (numéro)  \DOCTITRE
\documentclass[11pt]{article}

\usepackage{fontspec}
\usepackage[french]{babel} % français = langue principale ; le chinois via \zh{..} / textebox (xeCJK)
\usepackage{xeCJK}
\usepackage{pifont} % \ding{51} \ding{55} utilisés par les modèles
\usepackage[a4paper,margin=2cm,top=3.4cm,bottom=2.8cm,headheight=1.4cm,footskip=1.3cm]{geometry}
\usepackage{amsmath,amssymb,amsfonts}
\usepackage{mathtools} % \xmapsto, \coloneqq... souvent utilisés par les modèles
\usepackage{cancel} % simplifications barrées (bilans de forces)
% \abs{z} (module, valeur absolue) et \norm{u} (norme), utilisés par les modèles
\providecommand{\abs}{}\let\abs\relax\DeclarePairedDelimiter{\abs}{\lvert}{\rvert}
\providecommand{\norm}{}\let\norm\relax\DeclarePairedDelimiter{\norm}{\lVert}{\rVert}
\usepackage[table]{xcolor} % [table] : fournit aussi \rowcolors (alternance de lignes)
\usepackage{pagecolor}
\usepackage{tikz}
\usetikzlibrary{arrows.meta,positioning,calc,shapes.geometric,shapes.misc,shapes.arrows,decorations.pathmorphing,decorations.pathreplacing,decorations.markings,patterns,shadows,mindmap,trees,fit,backgrounds,matrix,intersections,angles,quotes}
\usepackage{pgfplots}
\usepackage[version=4]{mhchem} % équations nucléaires \ce{^{235}_{92}U}
\usepackage[european,siunitx]{circuitikz} % schémas électriques
\pgfplotsset{compat=1.18}
\usepgfplotslibrary{fillbetween,groupplots} % aires entre courbes (calcul intégral)
\usepackage{enumitem}
\usepackage{fancyhdr}
% [most] charge toutes les bibliothèques tcolorbox (listings, minted, raster,
% documentation...) et gonfle inutilement la mémoire TeX sur un document long
% (~300 boîtes) : on ne charge que ce qui est réellement utilisé.
\usepackage[skins,breakable]{tcolorbox}
\usepackage{graphicx}
\usepackage{colortbl}
\usepackage{booktabs}
\usepackage{tabularx}
\usepackage{diagbox} % case d'en-tête diagonale des tableaux à double entrée
% Colonnes extensibles centrée / gauche / droite (C, L, R), souvent utilisées par les modèles
\newcolumntype{C}{>{\centering\arraybackslash}X}
\newcolumntype{L}{>{\raggedright\arraybackslash}X}
\newcolumntype{R}{>{\raggedleft\arraybackslash}X}
\usepackage{array}
\usepackage{multicol}
\usepackage{multirow}
\usepackage{siunitx}
% parse-numbers=false : accepte aussi \SI{2,5 \times 10^{-3}}{...}
\sisetup{locale=FR,output-decimal-marker={,},group-minimum-digits=4,parse-numbers=false}
\DeclareSIUnit\ohm{\ensuremath{\symup{\Omega}}} % U+2126 absent de la police
\DeclareSIUnit\division{div} % réglages d'oscilloscope (ms/div, V/div)
\DeclareSIUnit\tr{tr} % tours (vitesse de rotation, ex. \SI{1500}{\tr\per\minute})
\DeclareSIUnit\molar{M} % concentration molaire (ex. \SI{3}{\molar}, \SI{1}{\micro\molar})
\DeclareSIUnit\an{an} % année (le modèle confond parfois avec \year, un compteur TeX) — ex. \SI{5}{\kilo\meter\per\an}
\DeclareSIUnit\FCFA{FCFA} % franc CFA (ex. \SI{500}{\FCFA\per\kilo\gram})
\DeclareSIUnit\degreeBrix{\textdegree Brix} % teneur en sucre d'un sirop/jus (réfractométrie)
\DeclareSIUnit\month{mois} % le modèle utilise parfois \month, un compteur TeX, comme unité de durée
\usepackage{qrcode}
% \degree : commande fréquemment utilisée par les modèles pour un angle ou une
% température en dehors de \SI{}{} (non fournie par les paquets chargés).
\providecommand{\degree}{^{\circ}}
% \qed (fin de démonstration), utilisé par les modèles sans amsthm
\providecommand{\qed}{\hfill$\square$}
% \barre : double barre des valeurs interdites dans un tableau de variations (tkz-tab)
\providecommand{\barre}{\ensuremath{\|}}
% environnement proof (amsthm non chargé) utilisé par les modèles
\ifdefined\proof\else\newenvironment{proof}[1][Démonstration]{\par\noindent\textit{#1.}\ }{\qed\par}\fi
\usepackage{lastpage}
\usepackage{hyperref}
\hypersetup{hidelinks}

% ============================================================
% Palette « Pop » OnBuch+ — mêmes hex que la classe P (plus_ui.dart)
% ============================================================
\definecolor{poporange}{HTML}{F59321}
\definecolor{popdark}{HTML}{E07A0C}
\definecolor{popcream}{HTML}{FFF6E8}
\definecolor{popink}{HTML}{1C1714}
\definecolor{poppurple}{HTML}{7A5AE0}
\definecolor{popgreen}{HTML}{2E9E6A}
\definecolor{poppink}{HTML}{E0457B}
\definecolor{popblue}{HTML}{2F7DE1}
\definecolor{popgold}{HTML}{C98A12}
\definecolor{popmuted}{HTML}{8A7F76}
\definecolor{popline}{HTML}{EADFD2}
\definecolor{popgray}{HTML}{8A7F76}
\colorlet{popgrey}{popgray}
% Alias tolérants (le modèle nomme parfois les couleurs différemment)
\colorlet{popwhite}{white}
\colorlet{popblack}{popink}
\colorlet{popred}{poppink}
\colorlet{popyellow}{popgold}
% Teintes claires (fonds de tableaux/encadrés) — le modèle nomme parfois
% les couleurs avec un suffixe L (ex. popblueL) sans les avoir définies.
\colorlet{poporangeL}{poporange!20}
\colorlet{popdarkL}{popdark!20}
\colorlet{poppurpleL}{poppurple!20}
\colorlet{popgreenL}{popgreen!20}
\colorlet{poppinkL}{poppink!20}
\colorlet{popblueL}{popblue!20}
\colorlet{popgoldL}{popgold!20}
\colorlet{popredL}{popred!20}
\colorlet{popgrayL}{popgray!20}
% Teintes claires pour fonds de tableaux / zones
\colorlet{popblueL}{popblue!10!white}
\colorlet{poporangeL}{poporange!14!white}
\colorlet{popgreenL}{popgreen!12!white}
\colorlet{poppurpleL}{poppurple!12!white}
\colorlet{poppinkL}{poppink!12!white}
\colorlet{popgoldL}{popgold!14!white}

\pagecolor{popcream}

% ============================================================
% Typographie
% ============================================================
\defaultfontfeatures{Ligatures=TeX}
\setmainfont{PlusJakartaSans-Medium.ttf}[Path=../../../fonts/,BoldFont=PlusJakartaSans-Bold.ttf]
% Symboles, API et emojis absents de la police principale : repli sur DejaVu Sans (caractères actifs)
\newfontface\symfont{DejaVuSans.ttf}[Path=../../../fonts/]
\makeatletter
\newcommand\symactive[1]{\catcode#1=13 \begingroup\lccode`\~=#1 \lowercase{\endgroup\def~}{{\symfont\char#1}}}
\newcommand\symblank[1]{\catcode#1=13 \begingroup\lccode`\~=#1 \lowercase{\endgroup\def~}{}}
\newcommand\symhyph[1]{\catcode#1=13 \begingroup\lccode`\~=#1 \lowercase{\endgroup\def~}{\mbox{-}}}
\newcount\symc
\newcommand\symrange[2]{\symc=#1 \@whilenum\symc<#2 \do{\expandafter\symactive\expandafter{\the\symc}\advance\symc by 1 }}
\newcommand\symblankrange[2]{\symc=#1 \@whilenum\symc<#2 \do{\expandafter\symblank\expandafter{\the\symc}\advance\symc by 1 }}
\makeatother
\symrange{"0250}{"02FF}\symactive{"2713}\symactive{"2714}\symactive{"2717}\symactive{"2718}\symactive{"2610}\symactive{"2611}\symactive{"2612}
\symrange{"2600}{"27BF}
\catcode"2705=13 \begingroup\lccode`\~="2705 \lowercase{\endgroup\def~}{{\symfont\char"2713}}\symblank{"26BD}\symblank{"26A1}
\symrange{"2460}{"257F}\symactive{"223C}\symactive{"2248}\symactive{"2260}
\symactive{"01F8}\symactive{"01F9}\symactive{"1E3E}\symactive{"1E3F}
\symhyph{"2011}\symhyph{"2E1A}\symblankrange{"1F300}{"1FAFF}\symblank{"FE0F}
% Caractères chinois : WenQuanYi Zen Hei (sous-ensemble GB2312 + pinyin), gras/italique simulés
\setCJKmainfont{WenQuanYiZenHei-subset.ttf}[Path=../../../fonts/,AutoFakeBold=3,AutoFakeSlant=0.2]
\xeCJKsetup{CJKspace=false}
% Pinyin : voyelles à caron (ǎ ǐ ǒ ǔ ǖ ǘ ǚ ǜ) absentes de la police principale -> caractères actifs
\newfontface\pyfont{WenQuanYiZenHei-subset.ttf}[Path=../../../fonts/]
\newcommand\pyactive[1]{\catcode#1=13 \begingroup\lccode`\~=#1 \lowercase{\endgroup\def~}{{\pyfont\char#1}}}
\pyactive{"01CD}\pyactive{"01CE}\pyactive{"01CF}\pyactive{"01D0}\pyactive{"01D1}\pyactive{"01D2}\pyactive{"01D3}\pyactive{"01D4}\pyactive{"01D5}\pyactive{"01D6}\pyactive{"01D7}\pyactive{"01D8}\pyactive{"01D9}\pyactive{"01DA}\pyactive{"01DB}\pyactive{"01DC}
% Maths en sans-serif (Fira Math) avec chiffres et lettres droites de Plus
% Jakarta, pour que les formules se fondent dans le texte.
\usepackage[math-style=ISO,bold-style=ISO]{unicode-math}
\setmathfont{FiraMath-Regular.otf}[Path=../../../fonts/]
\setmathfont{PlusJakartaSans-Medium.ttf}[Path=../../../fonts/,range={up/num,up/{Latin,latin},\mathup}]
\usepackage{icomma} % virgule décimale sans espace en mode math
% Caractères mathématiques Unicode absents des polices de texte (Plus Jakarta,
% Archivo Black) : rendus par leur équivalent mathématique, en texte comme en
% formule — sinon ils s'affichent en carré vide (ex. « Arithmétique dans ℤ »).
\usepackage{newunicodechar}
\newunicodechar{ℝ}{\ensuremath{\mathbb{R}}}
\newunicodechar{ℤ}{\ensuremath{\mathbb{Z}}}
\newunicodechar{ℂ}{\ensuremath{\mathbb{C}}}
\newunicodechar{ℕ}{\ensuremath{\mathbb{N}}}
\newunicodechar{ℚ}{\ensuremath{\mathbb{Q}}}
\newunicodechar{𝔻}{\ensuremath{\mathbb{D}}}
\newunicodechar{α}{\ensuremath{\alpha}}
\newunicodechar{β}{\ensuremath{\beta}}
\newunicodechar{γ}{\ensuremath{\gamma}}
\newunicodechar{δ}{\ensuremath{\delta}}
\newunicodechar{ε}{\ensuremath{\varepsilon}}
\newunicodechar{θ}{\ensuremath{\theta}}
\newunicodechar{λ}{\ensuremath{\lambda}}
\newunicodechar{μ}{\ensuremath{\mu}}
\newunicodechar{σ}{\ensuremath{\sigma}}
\newunicodechar{τ}{\ensuremath{\tau}}
\newunicodechar{φ}{\ensuremath{\varphi}}
\newunicodechar{ψ}{\ensuremath{\psi}}
\newunicodechar{ω}{\ensuremath{\omega}}
\newunicodechar{Σ}{\ensuremath{\Sigma}}
% majuscules produites par \MakeUppercase dans les titres (α-aminés -> Α)
\newunicodechar{Α}{\ensuremath{\alpha}}
\newunicodechar{Β}{\ensuremath{\beta}}
\newunicodechar{Γ}{\ensuremath{\Gamma}}
\newunicodechar{Δ}{\ensuremath{\Delta}}
\newunicodechar{Ε}{\ensuremath{\varepsilon}}
\newunicodechar{Θ}{\ensuremath{\Theta}}
\newunicodechar{Λ}{\ensuremath{\Lambda}}
\newunicodechar{Μ}{\ensuremath{\mu}}
\newunicodechar{Τ}{\ensuremath{\tau}}
\newunicodechar{Φ}{\ensuremath{\Phi}}
\newunicodechar{Ψ}{\ensuremath{\Psi}}
\newunicodechar{Ω}{\ensuremath{\Omega}}
\newunicodechar{↦}{\ensuremath{\mapsto}}
\newunicodechar{∈}{\ensuremath{\in}}
\newunicodechar{∉}{\ensuremath{\notin}}
\newunicodechar{∧}{\ensuremath{\wedge}}
\newunicodechar{⇔}{\ensuremath{\Leftrightarrow}}
\newunicodechar{⟺}{\ensuremath{\Longleftrightarrow}}
\newunicodechar{⇒}{\ensuremath{\Rightarrow}}
\newunicodechar{⟹}{\ensuremath{\Longrightarrow}}
\newunicodechar{∘}{\ensuremath{\circ}}
\newunicodechar{∪}{\ensuremath{\cup}}
\newunicodechar{∩}{\ensuremath{\cap}}
\newunicodechar{≡}{\ensuremath{\equiv}}
\newunicodechar{⊂}{\ensuremath{\subset}}
\newunicodechar{⊥}{\ensuremath{\perp}}
\newunicodechar{∃}{\ensuremath{\exists}}
\newunicodechar{∀}{\ensuremath{\forall}}
\newunicodechar{∛}{\ensuremath{\sqrt[3]{\,}}}
\newunicodechar{∜}{\ensuremath{\sqrt[4]{\,}}}
\newunicodechar{⃗}{\ensuremath{{}^{\to}}}
\newunicodechar{ⁿ}{\ifmmode^{n}\else\textsuperscript{\ensuremath{n}}\fi}
\newunicodechar{ˣ}{\ifmmode^{x}\else\textsuperscript{\ensuremath{x}}\fi}
\newunicodechar{ᵇ}{\ifmmode^{b}\else\textsuperscript{\ensuremath{b}}\fi}
\newunicodechar{ʳ}{\ifmmode^{r}\else\textsuperscript{\ensuremath{r}}\fi}
\newunicodechar{ᵘ}{\ifmmode^{u}\else\textsuperscript{\ensuremath{u}}\fi}
\newunicodechar{ʸ}{\ifmmode^{y}\else\textsuperscript{\ensuremath{y}}\fi}
\newunicodechar{ᵃ}{\ifmmode^{a}\else\textsuperscript{\ensuremath{a}}\fi}
\newunicodechar{ᵉ}{\ifmmode^{e}\else\textsuperscript{\ensuremath{e}}\fi}
\newunicodechar{ᵏ}{\ifmmode^{k}\else\textsuperscript{\ensuremath{k}}\fi}
\newunicodechar{ᵖ}{\ifmmode^{p}\else\textsuperscript{\ensuremath{p}}\fi}
\newunicodechar{ᴺ}{\ifmmode^{N}\else\textsuperscript{\ensuremath{N}}\fi}
\newunicodechar{ᶜ}{\ifmmode^{c}\else\textsuperscript{\ensuremath{c}}\fi}
\newunicodechar{ʰ}{\ifmmode^{h}\else\textsuperscript{\ensuremath{h}}\fi}
\newunicodechar{ᵠ}{\ifmmode^{q}\else\textsuperscript{\ensuremath{q}}\fi}
\newunicodechar{ₙ}{\ifmmode_{n}\else\textsubscript{\ensuremath{n}}\fi}
\newunicodechar{ₐ}{\ifmmode_{a}\else\textsubscript{\ensuremath{a}}\fi}
\newunicodechar{ᵢ}{\ifmmode_{i}\else\textsubscript{\ensuremath{i}}\fi}
\newunicodechar{ᵣ}{\ifmmode_{r}\else\textsubscript{\ensuremath{r}}\fi}
\newunicodechar{ₖ}{\ifmmode_{k}\else\textsubscript{\ensuremath{k}}\fi}
\newunicodechar{ᵦ}{\ifmmode_{\beta}\else\textsubscript{\ensuremath{\beta}}\fi}
\newunicodechar{ₓ}{\ifmmode_{x}\else\textsubscript{\ensuremath{x}}\fi}
\newfontfamily\popdisplay{ArchivoBlack-Regular.ttf}[Path=../../../fonts/]
\newfontfamily\poplogo{SpaceGrotesk-Bold.ttf}[Path=../../../fonts/]
\newfontfamily\popsemi{PlusJakartaSans-SemiBold.ttf}[Path=../../../fonts/]
\newfontfamily\popxbold{PlusJakartaSans-ExtraBold.ttf}[Path=../../../fonts/]
\newfontfamily\popxboldls{PlusJakartaSans-ExtraBold.ttf}[Path=../../../fonts/,LetterSpace=3.0]

\setlist[enumerate]{leftmargin=1.7em,itemsep=5pt,topsep=4pt}
\setlist[itemize]{leftmargin=1.7em,itemsep=4pt,topsep=4pt,label={\color{poporange}\textbullet}}
\setlength{\parindent}{0pt}
\setlength{\parskip}{5pt}
\renewcommand{\arraystretch}{1.25}

% pgfplots : style Pop par défaut
\pgfplotsset{
  every axis/.append style={axis line style={popink,line width=1pt},tick style={popink},
    grid style={popline,line width=0.5pt},label style={font=\small},tick label style={font=\footnotesize}},
  popcurve/.style={poporange,line width=1.8pt,no markers,smooth},
}
% Même style disponible dans un \draw TikZ simple (hors pgfplots)
\tikzset{popcurve/.style={poporange,line width=1.8pt}}
\tikzset{popgrid/.style={popline,very thin}}
\colorlet{popcurve}{poporange} % le nom du style est parfois utilisé comme couleur

% ============================================================
% Pill / squiggle
% ============================================================
\newcommand{\poppill}[3]{%
  \tikz[baseline=-0.55ex]\node[draw=popink,line width=1.2pt,rounded corners=2.6pt,%
    fill=#2,inner xsep=6.5pt,inner ysep=3pt]{%
    \popxboldls\fontsize{7}{7}\selectfont\color{#3}\MakeUppercase{#1}};%
}
\newcommand{\popsquiggle}[1]{%
  \tikz[baseline=-0.4ex]{
    \draw[#1,line width=2.6pt,line cap=round]
      plot [smooth] coordinates {(0,0.09) (0.34,0.01) (0.68,0.21) (1.02,0.01) (1.36,0.10)};
  }%
}
% Mot-clé surligné (marqueur orange sous le texte)
\newcommand{\cle}[1]{{\popxbold\color{popdark}#1}}

% ============================================================
% En-tête / pied
% ============================================================
\pagestyle{fancy}
\fancyhf{}
\renewcommand{\headrulewidth}{0pt}
\renewcommand{\footrulewidth}{0pt}
\lhead{\raisebox{-0.15ex}{\poplogo\fontsize{13}{13}\selectfont\color{popink}OnBuch\color{poporange}+}}
\rhead{\poppill{\DOCMATIERE\ \textbullet\ \DOCNIVEAU\ \textbullet\ Leçon \DOCLECON}{white}{popink}}
\lfoot{\popxbold\fontsize{7.6}{7.6}\selectfont\color{popmuted}\MakeUppercase{Cours signé OnBuch+}\ \textbullet\ {\popsemi\DOCTITRE}}
\rfoot{\tikz[baseline=-0.6ex]\node[rounded corners=8.5pt,fill=popink,minimum height=17pt,minimum width=17pt,inner xsep=5pt,inner sep=0]{\popxbold\fontsize{8}{8}\selectfont\color{popcream}\thepage\,/\,\pageref*{LastPage}};}

% ============================================================
% Titres
% ============================================================
\newcommand{\chaptitre}[2]{%
  \begin{tcolorbox}[enhanced,colback=poporange,colframe=popink,boxrule=2.6pt,arc=11pt,
      drop shadow=popink,left=15pt,right=15pt,top=12pt,bottom=11pt,before skip=3pt,after skip=18pt]
    \poppill{#2}{popink}{popcream}\par\vspace{9pt}
    {\raggedright\hyphenpenalty=10000\exhyphenpenalty=10000\popdisplay\fontsize{22}{24}\selectfont\color{popcream}\MakeUppercase{#1}\par}
  \end{tcolorbox}
}
% Section numérotée : pastille orange avec le numéro + titre Archivo
\newcounter{popsec}
\newcommand{\coursec}[1]{%
  \stepcounter{popsec}\setcounter{popsub}{0}%
  \par\vspace{16pt}\noindent
  \begin{minipage}{\linewidth}
  \tikz[baseline=-0.7ex]\node[draw=popink,line width=1.6pt,fill=poporange,rounded corners=4pt,minimum size=22pt,inner sep=0]{\popdisplay\fontsize{12}{12}\selectfont\color{popcream}\thepopsec};%
  \hspace{8pt}{\popdisplay\fontsize{15}{17}\selectfont\color{popink}\MakeUppercase{#1}}\par
  \vspace{4pt}\noindent{\color{poporange}\rule{36pt}{3.6pt}}
  \end{minipage}\par\nopagebreak\vspace{8pt}
}
\newcounter{popsub}
\newcommand{\courssub}[1]{%
  \stepcounter{popsub}%
  \par\vspace{9pt}\noindent{\popxbold\fontsize{11.5}{13}\selectfont\color{popink}{\color{poporange}\thepopsec.\thepopsub}\hspace{6pt}#1}\par\nopagebreak\vspace{4pt}
}

% ============================================================
% Boîtes pédagogiques — même recette : fond blanc, bordure épaisse colorée,
% ombre dure encre, badge plein. Titre optionnel [#1].
% ============================================================
\tcbset{popbase/.style={breakable,enhanced,colback=white,boxrule=2.3pt,arc=9pt,
  shadow={3pt}{-3pt}{0pt}{popink},left=14pt,right=14pt,top=11pt,bottom=11pt,before skip=11pt,after skip=15pt,
  fontupper=\popsemi\fontsize{10.6}{14.5}\selectfont\color{popink},
  fontlower=\popsemi\fontsize{10.6}{14.5}\selectfont\color{popink}}}
% Coche dessinée (le glyphe ✓ n'existe pas dans les polices du document)
\newcommand{\popcheck}[1]{\tikz[baseline=-0.5ex]\draw[#1,line width=1.6pt,line cap=round,line join=round](-0.13,0.0)--(-0.04,-0.09)--(0.13,0.1);}
\providecommand{\coche}{\popcheck{popgreen}}
\providecommand{\resultat}[1]{\par\smallskip\noindent\fcolorbox{popgreen}{popgreenL}{\parbox{\dimexpr\linewidth-2\fboxsep-2\fboxrule\relax}{\textbf{Résultat.} #1}}\par\smallskip}
\newcommand{\popboxhead}[3]{\poppill{#1}{#2}{white}\ifx\relax#3\relax\else\hspace{6pt}{\popxbold\fontsize{10.6}{12}\selectfont\color{popink}#3}\fi\par\vspace{7pt}}

\newtcolorbox{aretenir}[1][]{popbase,colframe=popgreen,before upper={\popboxhead{À retenir}{popgreen}{#1}}}
% Texte en chinois (document de lecture, dialogue, extrait) : cadre
% d'étude. Un titre optionnel (source, auteur) s'affiche dans l'en-tête.
\newtcolorbox{textebox}[1][]{popbase,colframe=popblue,before upper={\popboxhead{文本 · Texte}{popblue}{#1}},after upper={}}
% Marques ✓ / ✗ (forme correcte / fautive) souvent utilisées par les modèles
\providecommand{\cross}{\textcolor{poppink}{\ensuremath{\times}}}
\providecommand{\xmark}{\cross}\providecommand{\crossmark}{\cross}\providecommand{\cmark}{\ensuremath{\checkmark}}
% Ligne de réponse / justification à compléter (inventée par les modèles)
\providecommand{\justification}[1]{\par\noindent\textit{Justification~:} \dotfill\par}
% \textipa (package tipa non chargé) : la police ne couvre pas l'API -> simple passage
\providecommand{\textipa}[1]{#1}
% Soulignement (ulem non chargé) : \uline, \ul
\providecommand{\uline}[1]{\underline{#1}}\providecommand{\ul}[1]{\underline{#1}}
% \glossaire{\item ...} inventée par les modèles : liste de glossaire
\providecommand{\glossaire}[1]{\begin{itemize}#1\end{itemize}}
% \alert (beamer) inventée par les modèles : texte mis en évidence
\providecommand{\alert}[1]{\textbf{\textcolor{poppink}{#1}}}
% variantes ASCII d'ordinaux écrites en macro
\providecommand{\iere}{\textsuperscript{re}}\providecommand{\ieme}{\textsuperscript{e}}\providecommand{\ere}{\textsuperscript{re}}\providecommand{\eme}{\textsuperscript{e}}\providecommand{\er}{\textsuperscript{er}}\providecommand{\ie}{\textsuperscript{e}}
% Ordinaux français écrits en macro par les modèles : 1\ière, 3\ième
\newcommand{\ière}{\textsuperscript{re}}\newcommand{\ième}{\textsuperscript{e}}
% \exemple (inventée par les modèles) : étiquette « Exemple : »
\providecommand{\exemple}{\textbf{Exemple~:}\ }
% \square utilisable aussi hors mode math (cases à cocher)
\AtBeginDocument{\let\squareMath\square\renewcommand{\square}{\ensuremath{\squareMath}}}
% Mot ou phrase en chinois à l'intérieur d'un paragraphe français
\newcommand{\zh}[1]{\ifmmode\text{#1}\else{#1}\fi}
\let\eng\zh \let\esp\zh \let\all\zh % alias tolérés
% flèches d'intonation inventées par les modèles
\providecommand{\textsearrow}{}\renewcommand{\textsearrow}{\ensuremath{\searrow}}\providecommand{\textnearrow}{}\renewcommand{\textnearrow}{\ensuremath{\nearrow}}
% \fr{..} : mot français (inventé par les modèles)
\providecommand{\fr}[1]{\textit{#1}}
% zhbox : encadré de texte chinois inventé par les modèles
\newenvironment{zhbox}{\begin{quote}}{\end{quote}}
% \courssubsub : titre de niveau 3 inventé par les modèles
\providecommand{\courssubsub}[1]{\par\medskip\noindent\textbf{#1}\par\smallskip}
% \zhbf : texte chinois en gras inventé par les modèles
\providecommand{\zhbf}[1]{\textbf{#1}}
% \vs : « versus » inventé par les modèles
\providecommand{\vs}{\textit{vs}\ }
% \blank : case à compléter inventée par les modèles
\providecommand{\blank}[1][]{\underline{\hspace{1.6em}}}
\newtcolorbox{exemplebox}[1][]{popbase,colframe=poporange,before upper={\popboxhead{Exemple}{poporange}{#1}}}
\newtcolorbox{definition}[1][]{popbase,colframe=popblue,before upper={\popboxhead{Définition}{popblue}{#1}}}
\newtcolorbox{propriete}[1][]{popbase,colframe=poppurple,before upper={\popboxhead{Propriété}{poppurple}{#1}}}
\newtcolorbox{methode}[1][]{popbase,colframe=popgold,before upper={\popboxhead{Méthode}{popgold}{#1}}}
\newtcolorbox{attention}[1][]{popbase,colframe=poppink,before upper={\popboxhead{Attention}{poppink}{#1}}}
\newtcolorbox{experience}[1][]{popbase,colframe=popblue,colback=popblueL,before upper={\popboxhead{Expérience}{popblue}{#1}}}
% « Le savais-tu ? » : carte violette pleine (contexte, culture, Cameroun)
\newtcolorbox{savaistu}[1][]{popbase,colback=poppurple,colframe=popink,
  fontupper=\popsemi\fontsize{10.4}{14.2}\selectfont\color{white},
  before upper={\poppill{Le savais-tu\,?}{white}{poppurple}\ifx\relax#1\relax\else\hspace{6pt}{\popxbold\color{white}#1}\fi\par\vspace{7pt}}}
% Exercice résolu : énoncé en haut, \tcblower puis la solution sur fond crème
\newcounter{popexo}\newcounter{popexr}
\newtcolorbox{exoresolu}[1][]{popbase,colframe=poporange,colbacklower=popcream,
  segmentation style={popink,line width=1.2pt,dashed},
  before upper={\stepcounter{popexr}\popboxhead{Exercice résolu \thepopexr}{poporange}{#1}},
  before lower={\poppill{Solution}{popink}{popcream}\par\vspace{6pt}}}
% Exercice (non résolu en place) — la correction est dans la partie Corrigés
\newtcolorbox{exercice}[1][]{popbase,colframe=popink,
  before upper={\stepcounter{popexo}\popboxhead{Exercice \thepopexo}{popink}{#1}}}
% Corrigé
\newtcolorbox{corrige}[1][]{popbase,colframe=popgreen,colback=popgreenL,
  before upper={\popboxhead{Corrigé}{popgreen}{#1}}}
% Objectifs / prérequis / bilan
\newtcolorbox{objectifs}{popbase,colframe=popink,colback=poporangeL,
  before upper={\popboxhead{Objectifs de la leçon}{popink}{}}}
\newtcolorbox{prerequis}{popbase,colframe=popink,colback=popblueL,
  before upper={\popboxhead{Prérequis}{popblue}{}}}
\newtcolorbox{bilan}{popbase,colframe=popink,colback=white,boxrule=2.8pt,
  before upper={\popboxhead{Fiche bilan}{poporange}{L'essentiel à connaître pour l'examen}}}
% Figure encadrée avec légende
\newtcolorbox{popfigure}[1][]{enhanced,breakable=false,colback=white,colframe=popline,boxrule=1.4pt,arc=8pt,
  left=10pt,right=10pt,top=10pt,bottom=8pt,before skip=10pt,after skip=12pt,halign=center,
  lower separated=false,halign lower=center,
  fontlower=\popsemi\fontsize{9}{11.5}\selectfont\color{popmuted},
  #1}
\newcommand{\legende}[1]{\par\vspace{6pt}{\popsemi\fontsize{9}{11.5}\selectfont\color{popmuted}#1}}

% ============================================================
% Signature OnBuch+
% ============================================================
\newcommand{\poplogoinline}[1]{{\poplogo\fontsize{#1}{#1}\selectfont\color{popink}OnBuch\color{poporange}+}}
\newcommand{\signatureonbuch}{%
  \begin{tcolorbox}[enhanced,colback=popink,colframe=popink,boxrule=2.2pt,arc=10pt,
      drop shadow=poporange,left=14pt,right=14pt,top=10pt,bottom=10pt,before skip=0pt,after skip=0pt]
    \tikz[baseline=-0.7ex]\node[circle,fill=popgreen,draw=popcream,line width=1.3pt,minimum size=22pt,inner sep=0]{\popcheck{white}};%
    \hspace{9pt}\begin{minipage}[c]{0.62\linewidth}
      {\popxbold\fontsize{12}{13}\selectfont\color{popcream}Signé\ \textbullet\ {\poplogo OnBuch\color{poporange}+}}\par\vspace{2pt}
      {\raggedright\popsemi\fontsize{8}{10.5}\selectfont\color{popline}\MakeUppercase{Équipe pédagogique OnBuch+}\\\MakeUppercase{Conforme au programme officiel MINESEC}\par}
    \end{minipage}\hfill
    {\poplogo\fontsize{22}{22}\selectfont\color{popcream}OnBuch\color{poporange}+}
  \end{tcolorbox}%
}

% ============================================================
% Page de garde
% #1 titre · #2 matière · #3 niveau · #4 module · #5 n° leçon · #6 durée ·
% #7 description / objectifs (texte ou itemize)
% ============================================================
\newcommand{\pagegarde}[7]{%
  \thispagestyle{empty}
  \newgeometry{a4paper,margin=1.7cm,top=1.6cm,bottom=1.4cm}%
  \begin{tikzpicture}[remember picture,overlay]
    \fill[poporange!18] ([xshift=-3.2cm,yshift=-3.6cm]current page.north east) circle (4.6cm);
    \fill[poppurple!14] ([xshift=2.4cm,yshift=7.6cm]current page.south west) circle (3.8cm);
  \end{tikzpicture}%
  {\poplogo\fontsize{30}{30}\selectfont\color{popink}OnBuch\color{poporange}+}\hfill
  \raisebox{0.6ex}{\poppill{Cours signé \textbullet\ Premium}{popink}{popcream}}\par
  \vspace{1pt}\popsquiggle{poppurple}\par
  \vspace{8pt}
  \begin{tcolorbox}[enhanced,colback=poporange,colframe=popink,boxrule=2.8pt,arc=13pt,
      drop shadow=popink,left=18pt,right=18pt,top=12pt,bottom=13pt,after skip=10pt]
    \poppill{#2 \textbullet\ #3}{popink}{popcream}\hspace{5pt}\poppill{#4}{white}{popink}\par\vspace{8pt}
    {\popdisplay\fontsize{14}{14}\selectfont\color{popink}LEÇON #5}\par\vspace{4pt}
    {\raggedright\hyphenpenalty=10000\exhyphenpenalty=10000\popdisplay\fontsize{25}{27}\selectfont\color{popcream}\MakeUppercase{#1}\par}
    \vspace{8pt}
    {\popxbold\fontsize{9.5}{9.5}\selectfont\color{popink}\MakeUppercase{Durée conseillée : #6}}
  \end{tcolorbox}
  \begin{tcolorbox}[enhanced,colback=white,colframe=popink,boxrule=2.2pt,arc=10pt,
      drop shadow=popink,left=16pt,right=16pt,top=10pt,bottom=10pt,after skip=8pt]
    \poppill{Dans cette leçon}{poppurple}{white}\par\vspace{5pt}
    \popsemi\fontsize{9.3}{12.6}\selectfont\color{popink}#7
  \end{tcolorbox}
  \noindent\begin{minipage}[t]{0.487\linewidth}
    \begin{tcolorbox}[enhanced,colback=popgreen,colframe=popink,boxrule=2.2pt,arc=10pt,
        drop shadow=popink,left=11pt,right=11pt,top=10pt,bottom=10pt,height=3.1cm,valign=center]
      \tcbox[colback=white,colframe=popink,boxrule=1.3pt,arc=5pt,boxsep=0pt,nobeforeafter,
        left=3pt,right=3pt,top=3pt,bottom=3pt,valign=center]{\qrcode[height=1.9cm]{https://play.google.com/store/apps/details?id=com.luvvix.onbuch&hl=fr}}%
      \hspace{8pt}\begin{minipage}[c]{0.5\linewidth}
        {\popdisplay\fontsize{11}{12.5}\selectfont\color{white}\MakeUppercase{Télécharge OnBuch}}\par\vspace{4pt}
        {\raggedright\popsemi\fontsize{8.4}{11}\selectfont\color{white}Gratuit \textbullet\ Google Play\\Tuteur IA, annales, concours, résultats\par}
      \end{minipage}
    \end{tcolorbox}
  \end{minipage}\hfill
  \begin{minipage}[t]{0.487\linewidth}
    \begin{tcolorbox}[enhanced,colback=poppurple,colframe=popink,boxrule=2.2pt,arc=10pt,
        drop shadow=popink,left=11pt,right=11pt,top=10pt,bottom=10pt,height=3.1cm,valign=center]
      \tikz[baseline=-0.7ex]\node[circle,fill=white,draw=popink,line width=1.3pt,minimum size=1.6cm,inner sep=0]{\popdisplay\fontsize{24}{24}\selectfont\color{poppurple}+};%
      \hspace{8pt}\begin{minipage}[c]{0.6\linewidth}
        {\popdisplay\fontsize{11}{12.5}\selectfont\color{white}\MakeUppercase{Passe à OnBuch+}}\par\vspace{4pt}
        {\raggedright\popsemi\fontsize{8.4}{11}\selectfont\color{white}Tous les cours signés, Tuteur IA illimité, fiches PDF — onglet OnBuch+ de l'app\par}
      \end{minipage}
    \end{tcolorbox}
  \end{minipage}
  \vspace{8pt}
  \popcontenu
  \vfill
  \signatureonbuch
  \restoregeometry
  \newpage
}

% Rangée « Ce cours contient » (page de garde)
\newcommand{\poptile}[3]{% #1 couleur #2 gros texte #3 légende
  \begin{tcolorbox}[enhanced,colback=#1,colframe=popink,boxrule=2pt,arc=9pt,shadow={2.5pt}{-2.5pt}{0pt}{popink},
      left=6pt,right=6pt,top=9pt,bottom=9pt,halign=center,valign=center,height=1.75cm,width=\linewidth,nobeforeafter]
    {\popdisplay\fontsize{12.5}{13}\selectfont\color{white}\MakeUppercase{#2}}\par\vspace{3pt}
    {\centering\popsemi\fontsize{7.8}{9.5}\selectfont\color{white}#3\par}
  \end{tcolorbox}}
\newcommand{\popcontenu}{%
  \par\noindent{\popxboldls\fontsize{7.5}{7.5}\selectfont\color{popmuted}CE COURS SIGNÉ CONTIENT}\par\vspace{7pt}
  \noindent\begin{minipage}[t]{0.235\linewidth}\poptile{popblue}{Cours}{complet, illustré, pas à pas}\end{minipage}\hfill
  \begin{minipage}[t]{0.235\linewidth}\poptile{poporange}{Méthodes}{et exercices résolus}\end{minipage}\hfill
  \begin{minipage}[t]{0.235\linewidth}\poptile{poppink}{Exercices}{type examen + corrigés}\end{minipage}\hfill
  \begin{minipage}[t]{0.235\linewidth}\poptile{popgreen}{Fiche bilan}{l'essentiel à retenir}\end{minipage}\par}

% Sommaire
\newtcolorbox{sommaire}{enhanced,colback=white,colframe=popink,boxrule=2.2pt,arc=10pt,
  drop shadow=popink,left=18pt,right=18pt,top=13pt,bottom=13pt,before skip=6pt,after skip=20pt,
  before upper={\poppill{Sommaire}{poppurple}{white}\par\vspace{9pt}\popsemi\fontsize{10.6}{14.5}\selectfont\color{popink}}}

% Signature de fin de cours — poussée tout en bas de la dernière page
\newcommand{\findecours}{%
  \par\vfill\nopagebreak
  \begin{center}\popsquiggle{poporange}\end{center}\vspace{4pt}
  \signatureonbuch
}
\AtBeginDocument{\def\llbracket{\mathopen{[\mkern-3.5mu[}}\def\rrbracket{\mathclose{]\mkern-3.5mu]}}} % intervalles d'entiers (glyphe ⟦ absent de la police)
% Glyphes absents de Fira Math : redéfinis à partir de glyphes disponibles
\AtBeginDocument{%
  \def\setminus{\mathbin{\backslash}}\let\smallsetminus\setminus
  \def\vdots{\mathord{\vbox{\baselineskip3.2pt\lineskiplimit0pt\kern4pt\hbox{.}\hbox{.}\hbox{.}}}}%
  \def\ddots{\mathinner{\mkern1mu\raise6.5pt\hbox{.}\mkern2mu\raise3.6pt\hbox{.}\mkern2mu\raise0.7pt\hbox{.}\mkern1mu}}%
  \def\triangle{\mathord{\scalebox{0.9}{$\symup{\Delta}$}}}%
  \def\nabla{\mathord{\raisebox{\depth}{\scalebox{1}[-1]{$\symup{\Delta}$}}}}%
  \def\checkmark{\popcheck{popdark}}%
  \def\star{\mathbin{\ast}}\def\bigstar{\mathord{\ast}}%
  \def\diamond{\mathbin{\scalebox{0.6}{$\Diamond$}}}%
  \def\bigcup{\mathop{\mathchoice{\raisebox{-0.3em}{\scalebox{1.7}{$\cup$}}}{\scalebox{1.25}{$\cup$}}{\cup}{\cup}}\displaylimits}%
  \def\bigcap{\mathop{\mathchoice{\raisebox{-0.3em}{\scalebox{1.7}{$\cap$}}}{\scalebox{1.25}{$\cap$}}{\cap}{\cap}}\displaylimits}%
  \def\lesseqqgtr{\mathrel{\lessgtr}}\def\bigoplus{\mathop{\oplus}\displaylimits}%
}
\providecommand{\ssi}{si et seulement si} % abréviation fréquente
\AtBeginDocument{\providecommand{\wideparen}{\overparen}} % arc de cercle

%% FIGFIX-BEGIN v4 — correctifs globaux des figures (docs/onbuch-generation/tools/figfix)
\usepackage{adjustbox}\usepackage{etoolbox}
% toute figure TikZ/pgfplots plus large que la ligne est réduite à la largeur disponible
\BeforeBeginEnvironment{tikzpicture}{\begin{adjustbox}{max width=\linewidth}}
\AfterEndEnvironment{tikzpicture}{\end{adjustbox}}
% légendes pgfplots lisibles par-dessus les courbes
\pgfplotsset{every axis/.append style={legend style={fill=white,fill opacity=0.92,text opacity=1,draw=black!15,inner sep=3pt}}}
% étiquettes d'arêtes (syntaxe edge["..."]) avec fond blanc
\tikzset{every edge quotes/.append style={fill=white,inner sep=1.2pt}}
%% FIGFIX-END
\begin{document}

\pagegarde{Leçon 33 — Expression orale : économie}{Chinois}{1ère A}{Module 4 : Activités économiques et monde du travail}{ZHA33}{2 h}{Leçon d'expression orale consacrée au thème de l'économie : les élèves apprennent à présenter un exposé simple sur une activité économique, à donner leur avis et à participer à un débat guidé en chinois, en travaillant la prononciation des tons et les formules fonctionnelles de communication.
\vspace{4pt}
\begin{multicols}{2}\small
\begin{itemize}
\item À la fin de cette leçon, je sais utiliser le vocabulaire oral de base de l'économie (marché, prix, commerce, travail, argent)
\item À la fin de cette leçon, je sais présenter oralement un exposé simple de 5 à 8 phrases sur une activité économique du Cameroun ou de la Chine
\item À la fin de cette leçon, je sais donner mon avis avec les structures 我觉得…, 我认为…, 对我来说…
\item À la fin de cette leçon, je sais exprimer l'accord et le désaccord poliment (我同意 / 我不同意 / 你说得对，但是…)
\item À la fin de cette leçon, je sais relancer un échange et poser des questions (你呢？你怎么看？为什么？)
\item À la fin de cette leçon, je sais prononcer correctement les tons des mots économiques clés, notamment les suites de 3e tons
\end{itemize}\end{multicols}}

\begin{sommaire}
\begin{multicols}{2}
\begin{enumerate}
\item Vocabulaire oral de l'économie
\item Dialogues modèles avec pinyin
\item Formules fonctionnelles : présenter, opiner, relancer
\item Les tons à travailler
\item Jeux de rôle : au marché et à l'entreprise
\item Débat guidé : le commerce Cameroun-Chine
\item Activité d'intégration
\item Méthodes et exercices résolus
\item Exercices
\item Corrigés des exercices
\item Fiche bilan
\end{enumerate}
\end{multicols}
\end{sommaire}

\begin{objectifs}
\begin{itemize}
\item À la fin de cette leçon, je sais utiliser le vocabulaire oral de base de l'économie (marché, prix, commerce, travail, argent)
\item À la fin de cette leçon, je sais présenter oralement un exposé simple de 5 à 8 phrases sur une activité économique du Cameroun ou de la Chine
\item À la fin de cette leçon, je sais donner mon avis avec les structures 我觉得…, 我认为…, 对我来说…
\item À la fin de cette leçon, je sais exprimer l'accord et le désaccord poliment (我同意 / 我不同意 / 你说得对，但是…)
\item À la fin de cette leçon, je sais relancer un échange et poser des questions (你呢？你怎么看？为什么？)
\item À la fin de cette leçon, je sais prononcer correctement les tons des mots économiques clés, notamment les suites de 3e tons
\item À la fin de cette leçon, je sais participer à un jeu de rôle de négociation au marché
\item À la fin de cette leçon, je sais participer à un petit débat guidé sur le commerce Cameroun-Chine
\end{itemize}
\end{objectifs}

\begin{prerequis}
\begin{itemize}
\item Les quatre tons et le pinyin (Seconde)
\item Les structures de phrase de base : sujet + verbe + complément ; 是 / 有 / 在
\item Le vocabulaire des métiers et des achats (leçons antérieures du Module 4)
\item Les nombres, les prix et les classificateurs courants (块, 个, 斤)
\item Les questions simples avec 吗, 呢, 怎么
\end{itemize}
\end{prerequis}

\newpage

\coursec{Vocabulaire oral de l'économie}

\courssub{Situation d'accroche : au marché Mokolo}

Imaginez la scène : il est dix heures du matin au marché Mokolo de Yaoundé. Entre les étals de ndolé et de bananes plantains, un commerçant chinois négocie le prix d'un sac de riz avec une cliente camerounaise. La cliente fronce les sourcils et lance :

\zh{太贵了！便宜一点儿吧！} \emph{Tài guì le ! Piányi yìdiǎnr ba !} « C'est trop cher ! Un peu moins cher, s'il vous plaît ! »

Le commerçant sourit, réfléchit, puis répond : \zh{好吧，给你便宜五块钱。} \emph{Hǎo ba, gěi nǐ piányi wǔ kuài qián.} « D'accord, je vous fais cinq yuans de réduction. »

Cette petite scène, des milliers de Camerounais la vivent chaque jour. De plus en plus de Camerounais travaillent avec des partenaires chinois, dans le commerce, le BTP ou l'agriculture. Savoir parler d'économie en chinois — présenter un produit, donner son avis sur les prix, débattre du commerce Cameroun-Chine — est devenu un atout professionnel réel. Cette leçon vous donne les outils oraux pour le faire.

\textbf{Problématique :} quels mots faut-il maîtriser, et avec quelle prononciation, pour parler d'économie à l'oral en chinois ?

\courssub{Observation : un mini-texte pour découvrir le lexique}

Avant d'apprendre les mots un par un, lisons ce court texte qui les réunit presque tous. Soulignez mentalement les mots que vous devinez.

\begin{textebox}[L'économie du Cameroun en quelques mots]
喀麦隆的经济很有意思。在杜阿拉和雅温得，有很多市场和商店。\\
\emph{Kāmàilóng de jīngjì hěn yǒu yìsi. Zài Dùālā hé Yǎwēndé, yǒu hěn duō shìchǎng hé shāngdiàn.}\\
« L'économie du Cameroun est très intéressante. À Douala et à Yaoundé, il y a beaucoup de marchés et de magasins. »\\[4pt]
农民种咖啡、可可和香蕉，工人在工厂工作，商人在市场做生意。\\
\emph{Nóngmín zhòng kāfēi, kěkě hé xiāngjiāo, gōngrén zài gōngchǎng gōngzuò, shāngrén zài shìchǎng zuò shēngyi.}\\
« Les paysans cultivent le café, le cacao et les bananes, les ouvriers travaillent dans les usines, les commerçants font des affaires au marché. »\\[4pt]
喀麦隆出口咖啡、可可、石油和木材，进口很多中国的东西。\\
\emph{Kāmàilóng chūkǒu kāfēi, kěkě, shíyóu hé mùcái, jìnkǒu hěn duō Zhōngguó de dōngxi.}\\
« Le Cameroun exporte du café, du cacao, du pétrole et du bois, et importe beaucoup de produits chinois. »\\[4pt]
顾客喜欢便宜的东西，但是贵的东西质量常常更好。\\
\emph{Gùkè xǐhuan piányi de dōngxi, dànshì guì de dōngxi zhìliang chángcháng gèng hǎo.}\\
« Les clients aiment les produits bon marché, mais les produits chers sont souvent de meilleure qualité. »
\end{textebox}

\textbf{Questions de compréhension (à l'oral) :}
\begin{enumerate}
\item \zh{喀麦隆出口什么？} \emph{Kāmàilóng chūkǒu shénme ?} — Qu'est-ce que le Cameroun exporte ?
\item \zh{商人在哪儿做生意？} \emph{Shāngrén zài nǎr zuò shēngyi ?} — Où les commerçants font-ils des affaires ?
\item \zh{顾客喜欢什么样的东西？} \emph{Gùkè xǐhuan shénme yàng de dōngxi ?} — Quels produits les clients aiment-ils ?
\end{enumerate}

\courssub{Les lieux et les acteurs de l'économie}

\begin{definition}[Les mots-clés de l'économie]
\zh{经济} \emph{jīngjì} (économie) — \zh{市场} \emph{shìchǎng} (marché) — \zh{价格} \emph{jiàgé} (prix) — \zh{贸易} \emph{màoyì} (commerce) — \zh{生意} \emph{shēngyi} (affaires).
\end{definition}

Observons la formation des mots : le chinois construit son vocabulaire comme des briques. \zh{市场} « marché » contient \zh{市} « ville, marché » ; \zh{商店} « magasin » contient \zh{商} « commerce », que l'on retrouve dans \zh{商人} « commerçant » (littéralement « homme du commerce »). De même, \zh{工人} « ouvrier » = \zh{工} « travail » + \zh{人} « personne », et \zh{工厂} « usine » = \zh{工} + \zh{厂} « fabrique ». Apprendre les briques, c'est apprendre plusieurs mots à la fois.

\begin{center}
\begin{tabularx}{\linewidth}{|X|X|X|X|}
\hline
\rowcolor{popblueL} Caractères & Pinyin & Nature & Traduction \\
\hline
经济 & jīngjì & nom & économie \\
\hline
市场 & shìchǎng & nom & marché \\
\hline
商店 & shāngdiàn & nom & magasin, boutique \\
\hline
公司 & gōngsī & nom & entreprise, société \\
\hline
工厂 & gōngchǎng & nom & usine \\
\hline
商人 & shāngrén & nom & commerçant, homme d'affaires \\
\hline
工人 & gōngrén & nom & ouvrier \\
\hline
农民 & nóngmín & nom & paysan, agriculteur \\
\hline
顾客 & gùkè & nom & client \\
\hline
\end{tabularx}
\end{center}

\begin{exemplebox}[Les acteurs en action]
\zh{这位商人在杜阿拉有一家公司。} \emph{Zhè wèi shāngrén zài Dùālā yǒu yì jiā gōngsī.} « Ce commerçant a une entreprise à Douala. »\\
\zh{工人们在工厂里工作。} \emph{Gōngrénmen zài gōngchǎng li gōngzuò.} « Les ouvriers travaillent dans l'usine. »\\
\zh{这个农民种咖啡和可可。} \emph{Zhège nóngmín zhòng kāfēi hé kěkě.} « Ce paysan cultive du café et du cacao. »\\
\zh{顾客在商店里买东西。} \emph{Gùkè zài shāngdiàn li mǎi dōngxi.} « Les clients achètent des choses dans le magasin. »
\end{exemplebox}

\begin{attention}[Classificateurs : 家 pour les entreprises, 位 et 个 pour les personnes]
En français on dit « une entreprise », « une usine » sans réfléchir. En chinois, le nom ne s'emploie jamais seul après un nombre ou un démonstratif : il faut un \cle{classificateur}. Pour \zh{公司}, \zh{工厂} et \zh{商店}, le classificateur est \zh{家} : \zh{一家公司} \emph{yì jiā gōngsī} « une entreprise », \zh{这家商店} \emph{zhè jiā shāngdiàn} « ce magasin ». Ne dites jamais \zh{一公司} !\\
Pour les personnes, le classificateur général est \zh{个} : \zh{一个工人} \emph{yí ge gōngrén} « un ouvrier ». Le classificateur \zh{位} est sa version polie, réservée aux personnes que l'on respecte : \zh{这位商人} \emph{zhè wèi shāngrén} « ce commerçant (que je respecte) ». À l'oral, en cas de doute, utilisez \zh{个}.
\end{attention}

\courssub{L'argent et les prix}

\begin{center}
\begin{tabularx}{\linewidth}{|X|X|X|X|}
\hline
\rowcolor{popblueL} Caractères & Pinyin & Nature & Traduction \\
\hline
钱 & qián & nom & argent \\
\hline
价格 & jiàgé & nom & prix \\
\hline
贵 & guì & adjectif & cher \\
\hline
便宜 & piányi & adjectif & bon marché \\
\hline
块钱 & kuài qián & expr. & yuans (unité monétaire) \\
\hline
买 & mǎi & verbe & acheter \\
\hline
卖 & mài & verbe & vendre \\
\hline
\end{tabularx}
\end{center}

\begin{propriete}[Dire le prix : 块 / 毛 / 分]
La monnaie chinoise est le yuan, appelé familièrement \zh{块} \emph{kuài} à l'oral. Structure : \textbf{nombre + 块 (+ nombre + 毛) + 钱}.\\
\zh{五块钱} \emph{wǔ kuài qián} « cinq yuans » ; \zh{三块五毛钱} \emph{sān kuài wǔ máo qián} « trois yuans cinquante ».\\
Pour demander le prix : \zh{多少钱？} \emph{duōshao qián ?} « Combien ça coûte ? » ou \zh{这个多少钱？} \emph{Zhège duōshao qián ?} « Combien coûte ceci ? »
\end{propriete}

\begin{exemplebox}[Négocier les prix]
\zh{这个多少钱？} \emph{Zhège duōshao qián ?} « Combien coûte ceci ? »\\
\zh{太贵了！} \emph{Tài guì le !} « C'est trop cher ! »\\
\zh{便宜一点儿吧！} \emph{Piányi yìdiǎnr ba !} « Un peu moins cher, s'il vous plaît ! »\\
\zh{这个市场的东西很便宜。} \emph{Zhège shìchǎng de dōngxi hěn piányi.} « Les produits de ce marché sont bon marché. »
\end{exemplebox}

Notez la structure \zh{太 + adjectif + 了} (\zh{太贵了}) : elle exprime « trop... » ou une exclamation. C'est la phrase numéro un de toute négociation au marché !

\courssub{Les activités économiques et les produits du Cameroun}

\begin{center}
\begin{tabularx}{\linewidth}{|X|X|X|X|}
\hline
\rowcolor{popblueL} Caractères & Pinyin & Nature & Traduction \\
\hline
贸易 & màoyì & nom & commerce (international) \\
\hline
工作 & gōngzuò & nom / verbe & travail ; travailler \\
\hline
生意 & shēngyi & nom & affaires \\
\hline
出口 & chūkǒu & verbe / nom & exporter ; exportation \\
\hline
进口 & jìnkǒu & verbe / nom & importer ; importation \\
\hline
咖啡 & kāfēi & nom & café \\
\hline
可可 & kěkě & nom & cacao \\
\hline
香蕉 & xiāngjiāo & nom & banane \\
\hline
石油 & shíyóu & nom & pétrole \\
\hline
木材 & mùcái & nom & bois \\
\hline
\end{tabularx}
\end{center}

\begin{propriete}[出口 et 进口 : la logique des caractères]
\zh{出口} \emph{chūkǒu} = \zh{出} « sortir » + \zh{口} « bouche, port » : littéralement « faire sortir par le port », donc \emph{exporter}. \zh{进口} \emph{jìnkǒu} = \zh{进} « entrer » + \zh{口} : \emph{importer}. Structure : \textbf{Sujet (pays) + 出口/进口 + produit}.\\
\zh{喀麦隆出口咖啡和可可。} \emph{Kāmàilóng chūkǒu kāfēi hé kěkě.} « Le Cameroun exporte du café et du cacao. »\\
\zh{中国进口石油和木材。} \emph{Zhōngguó jìnkǒu shíyóu hé mùcái.} « La Chine importe du pétrole et du bois. »
\end{propriete}

Remarquez aussi les mots transparents : \zh{石油} « pétrole » = \zh{石} « pierre » + \zh{油} « huile » ; \zh{木材} « bois » = \zh{木} « arbre » + \zh{材} « matériau ». Et \zh{咖啡} \emph{kāfēi} est un emprunt phonétique au mot « café » — facile à retenir !

\begin{popfigure}
\begin{tikzpicture}[mindmap, grow cyclic, every node/.style={concept, align=center},
root concept/.append style={concept color=popblue, font=\Large},
level 1/.append style={level distance=3.2cm, sibling angle=90},
level 1 concept/.append style={concept color=poporange, font=\small},
level 2/.append style={level distance=2.2cm},
level 2 concept/.append style={concept color=popgreenL, font=\footnotesize}]
\node[root concept, align=center]{经济\\ jīngjì}
  child { node[align=center]{Lieux\\ 市场 商店} [clockwise from=135]
    child { node[align=center]{公司\\ gōngsī} }
    child { node[align=center]{工厂\\ gōngchǎng} } }
  child { node[align=center]{Argent\\ 钱 价格} [clockwise from=45]
    child { node {贵 guì} }
    child { node[align=center]{便宜\\ piányi} } }
  child { node[align=center]{Actions\\ 买 卖} [clockwise from=-45]
    child { node[align=center]{出口\\ chūkǒu} }
    child { node[align=center]{进口\\ jìnkǒu} } }
  child { node[align=center]{Produits\\ 咖啡 可可} [clockwise from=-135]
    child { node[align=center]{石油\\ shíyóu} }
    child { node[align=center]{木材\\ mùcái} } };
\end{tikzpicture}
\legende{Carte mentale du vocabulaire de l'économie : quatre familles autour de 经济.}
\end{popfigure}

\courssub{Prononciation ciblée : 买 mǎi contre 卖 mài}

Voici l'opposition la plus importante de cette leçon. \zh{买} \emph{mǎi} « acheter » se prononce avec le \textbf{3\ieme{} ton} : la voix descend puis remonte, comme un « v » mélodique. \zh{卖} \emph{mài} « vendre » se prononce avec le \textbf{4\ieme{} ton} : la voix tombe brutalement, comme un ordre sec.

\begin{popfigure}
\begin{tikzpicture}[scale=1.0]
% Grille 3e ton
\draw[popline] (0,0) -- (3,0) -- (3,2) -- (0,2) -- cycle;
\draw[popline, dashed] (0,1) -- (3,1);
\draw[very thick, popblue] (0.3,1.4) .. controls (1.0,0.3) and (1.8,0.3) .. (2.7,1.9);
\node[align=center] at (1.5,-0.7) {买 mǎi\\ 3e ton : descend puis monte};
% Grille 4e ton
\draw[popline] (5,0) -- (8,0) -- (8,2) -- (5,2) -- cycle;
\draw[popline, dashed] (5,1) -- (8,1);
\draw[very thick, poporange] (5.3,1.9) -- (7.7,0.2);
\node[align=center] at (6.5,-0.7) {卖 mài\\ 4e ton : chute brutale};
\end{tikzpicture}
\legende{Courbes mélodiques : mǎi (3e ton, en « v ») contre mài (4e ton, en chute).}
\end{popfigure}

\begin{attention}[L'erreur qui inverse le sens !]
Confondre \zh{买} \emph{mǎi} (acheter, 3\ieme{} ton) et \zh{卖} \emph{mài} (vendre, 4\ieme{} ton) est l'erreur classique des francophones. Un mauvais ton inverse complètement le sens de la phrase !\\
\zh{我买咖啡。} \emph{Wǒ mǎi kāfēi.} « J'achète du café. »\\
\zh{我卖咖啡。} \emph{Wǒ mài kāfēi.} « Je vends du café. »\\
Au marché, si vous dites \zh{我卖} au lieu de \zh{我买}, le commerçant croira que vous voulez lui vendre quelque chose ! Astuce : le 3\ieme{} ton « hésite » (descend-monte) comme l'acheteur qui réfléchit ; le 4\ieme{} ton « tranche » comme le vendeur qui conclut la vente.
\end{attention}

\begin{savaistu}[买 et 卖 : des caractères frères]
Les caractères \zh{买} (acheter) et \zh{卖} (vendre) se ressemblent beaucoup : \zh{卖} = \zh{买} + \zh{十} en haut. Astuce de mémorisation : celui qui vend a « dix » (\zh{十}) articles de plus à écouler ! Dans les caractères traditionnels, la famille était encore plus visible : la clé « coquillage », ancienne monnaie chinoise, apparaissait dans ces caractères — la preuve que l'écriture chinoise raconte l'histoire du commerce.
\end{savaistu}

\courssub{Entraînement choral : la gym des tons}

\begin{experience}[Répétition chorale avec gestes des tons]
En classe entière, répétez chaque mot après le professeur, en dessinant le ton dans l'air avec la main :
\begin{enumerate}
\item \textbf{3\ieme{} ton} (main en « v ») : \zh{买} mǎi, \zh{我} wǒ, \zh{市场} shìchǎng, \zh{便宜} piányi.
\item \textbf{4\ieme{} ton} (main qui coupe vers le bas) : \zh{卖} mài, \zh{贵} guì, \zh{贸易} màoyì, \zh{价格} jiàgé.
\item \textbf{Opposition} : mǎi — mài, mǎi — mài, cinq fois de suite, de plus en plus vite.
\item \textbf{Phrase complète} : \zh{我买，他卖。} \emph{Wǒ mǎi, tā mài.} « Moi j'achète, lui il vend. »
\end{enumerate}
Puis, par binômes : l'un dit un mot, l'autre doit dire s'il a entendu « acheter » ou « vendre ».
\end{experience}

\begin{exoresolu}[Exercice : complétez avec le bon mot]
Énoncé — Complétez chaque phrase avec \zh{买}, \zh{卖}, \zh{出口} ou \zh{进口}, puis traduisez.\\
1. 喀麦隆 \_\_\_\_ 咖啡和可可。\\
2. 中国在喀麦隆 \_\_\_\_ 很多机器。\\
3. 顾客在市场 \_\_\_\_ 东西。\\
4. 商人在商店 \_\_\_\_ 东西。
\tcblower
Solution détaillée :\\
1. \zh{喀麦隆出口咖啡和可可。} \emph{Kāmàilóng chūkǒu kāfēi hé kěkě.} « Le Cameroun exporte du café et du cacao. » — Le café et le cacao \emph{sortent} du pays : \zh{出口}.\\
2. \zh{中国在喀麦隆卖很多机器。} \emph{Zhōngguó zài Kāmàilóng mài hěn duō jīqì.} « La Chine vend beaucoup de machines au Cameroun. » — La Chine \emph{vend} : \zh{卖} (4\ieme{} ton).\\
3. \zh{顾客在市场买东西。} \emph{Gùkè zài shìchǎng mǎi dōngxi.} « Les clients achètent des choses au marché. » — Le client \emph{achète} : \zh{买} (3\ieme{} ton).\\
4. \zh{商人在商店卖东西。} \emph{Shāngrén zài shāngdiàn mài dōngxi.} « Les commerçants vendent des choses dans le magasin. » — Le commerçant \emph{vend} : \zh{卖}.\\
Vérifiez toujours le sens avant de choisir le ton : qui donne l'argent (\zh{买}) ? Qui reçoit l'argent (\zh{卖}) ?
\end{exoresolu}

\begin{aretenir}[L'essentiel de la section]
\begin{itemize}
\item Quatre familles de mots : les \textbf{lieux et acteurs} (\zh{市场}, \zh{公司}, \zh{商人}, \zh{顾客}...), l'\textbf{argent} (\zh{钱}, \zh{价格}, \zh{贵}, \zh{便宜}), les \textbf{actions} (\zh{买}, \zh{卖}, \zh{出口}, \zh{进口}) et les \textbf{produits} (\zh{咖啡}, \zh{可可}, \zh{石油}, \zh{木材}).
\item Le prix se dit avec \zh{块} : \zh{五块钱} « cinq yuans » ; on le demande avec \zh{多少钱？}.
\item \zh{买} mǎi (3\ieme{} ton) = acheter ; \zh{卖} mài (4\ieme{} ton) = vendre. Le ton change le sens : entraînez-vous chaque jour avec les gestes des tons.
\item Les caractères se construisent par briques (\zh{商} → \zh{商店}, \zh{商人} ; \zh{口} → \zh{出口}, \zh{进口}) : apprenez les familles, pas les mots isolés.
\end{itemize}
\end{aretenir}

\coursec{Dialogues modèles avec pinyin}

Dans la section précédente, nous avons découvert le vocabulaire oral de l'économie : les mots du marché, des prix, du travail et des affaires. Mais un mot isolé ne parle pas tout seul ! Pour vraiment communiquer, il faut placer ces mots dans des échanges réels. C'est exactement l'objectif de cette section : observer deux dialogues modèles, comprendre leur structure, repérer les formules fonctionnelles, puis les faire vôtres en remplaçant les produits et les prix par des éléments de votre quotidien camerounais.

\courssub{Dialogue 1 — Au marché : la négociation}

\begin{textebox}[Dialogue 1 — 在市场买东西 (Au marché, on achète)]
A : 你好！这个多少钱？\\
\emph{Nǐ hǎo ! Zhège duōshao qián ?}\\
« Bonjour ! Combien coûte ceci ? »\\
B : 五十块。\\
\emph{Wǔshí kuài.}\\
« Cinquante yuans. »\\
A : 太贵了！便宜一点儿吧。\\
\emph{Tài guì le ! Piányi yìdiǎnr ba.}\\
« C'est trop cher ! Un peu moins cher, s'il vous plaît. »\\
B : 好吧，四十五块。\\
\emph{Hǎo ba, sìshíwǔ kuài.}\\
« D'accord, quarante-cinq yuans. »\\
A : 好，我买了。\\
\emph{Hǎo, wǒ mǎi le.}\\
« Bien, je le prends. »
\end{textebox}

\textbf{Observation.} Lisez d'abord le dialogue en silence, puis repérez les quatre moments de l'échange : on salue, on demande le prix, on discute le prix, on conclut. C'est exactement la même logique qu'au marché Mokolo à Yaoundé ou au marché central de Douala : la négociation est un art universel !

\begin{propriete}[Les quatre étapes d'une négociation en chinois]
\begin{enumerate}
\item \cle{问好} \emph{wènhǎo} : saluer — \zh{你好！} \emph{Nǐ hǎo !}
\item \cle{问价格} \emph{wèn jiàgé} : demander le prix — \zh{这个多少钱？} \emph{Zhège duōshao qián ?}
\item \cle{讨价还价} \emph{tǎojià huánjià} : marchander — \zh{太贵了！便宜一点儿吧。}
\item \cle{成交} \emph{chéngjiāo} : conclure — \zh{好，我买了。}
\end{enumerate}
L'expression \zh{讨价还价} est une expression figée de quatre caractères qui signifie littéralement « discuter le prix, rendre le prix » : c'est le mot exact pour « marchander ».
\end{propriete}

\begin{popfigure}
\begin{tikzpicture}[
  etape/.style={rectangle, rounded corners=4pt, draw=popblue, fill=popblueL, thick, minimum height=1.1cm, align=center, font=\small, text width=3.2cm},
  fleche/.style={-{Stealth[length=3mm]}, very thick, poporange},
  num/.style={circle, draw=poporange, fill=poporangeL, thick, inner sep=1.5pt, font=\small\bfseries}
]
\node[etape, align=center] (e1) at (0,0){问好 \emph{wènhǎo}\\ Saluer};
\node[etape, align=center] (e2) at (4.4,0){问价格 \emph{wèn jiàgé}\\ Demander le prix};
\node[etape, align=center] (e3) at (8.8,0){讨价还价 \emph{tǎojià huánjià}\\ Marchander};
\node[etape, align=center] (e4) at (13.2,0){成交 \emph{chéngjiāo}\\ Conclure};
\node[num] at (0,1.05) {1};
\node[num] at (4.4,1.05) {2};
\node[num] at (8.8,1.05) {3};
\node[num] at (13.2,1.05) {4};
\draw[fleche] (e1) -- (e2);
\draw[fleche] (e2) -- (e3);
\draw[fleche] (e3) -- (e4);
\end{tikzpicture}
\legende{Organigramme du dialogue de négociation : les quatre étapes à mémoriser.}
\end{popfigure}

\begin{definition}[Les mots-clés de la négociation]
\begin{itemize}
\item \zh{多少} \emph{duōshao} : combien (pour les nombres supérieurs à dix, les prix) ; \zh{多少钱？} = « combien d'argent ? » = « combien ça coûte ? »
\item \zh{块} \emph{kuài} : unité de monnaie (langue parlée, équivalent de \zh{元} \emph{yuán} à l'écrit) ; \zh{五十块} = « 50 yuans »
\item \zh{贵} \emph{guì} : cher ; \zh{便宜} \emph{piányi} : bon marché
\item \zh{太……了} \emph{tài... le} : « trop … » ; \zh{太贵了！} = « trop cher ! »
\item \zh{吧} \emph{ba} : particule finale de suggestion ou de demande polie
\end{itemize}
\end{definition}

\begin{attention}[一点儿 se place APRÈS l'adjectif]
En français, on dit « un peu \textbf{moins cher} » : « un peu » vient \emph{avant} l'adjectif. En chinois, c'est l'inverse : \zh{一点儿} \emph{yìdiǎnr} se place \emph{après} l'adjectif.\\
\zh{便宜一点儿。} \emph{Piányi yìdiǎnr.} « Un peu moins cher. »\\
\zh{慢一点儿说。} \emph{Màn yìdiǎnr shuō.} « Parlez un peu plus lentement. »\\
Ne dites jamais \zh{一点儿便宜} : c'est un calque du français, faute typique des francophones !
\end{attention}

\courssub{Dialogue 2 — Parler du travail}

\begin{textebox}[Dialogue 2 — 谈工作 (Parler du travail)]
A : 你爸爸做什么工作？\\
\emph{Nǐ bàba zuò shénme gōngzuò ?}\\
« Que fait ton père comme travail ? »\\
B : 他是商人，在杜阿拉做生意。\\
\emph{Tā shì shāngrén, zài Dù'ālā zuò shēngyi.}\\
« Il est commerçant, il fait des affaires à Douala. »\\
A : 生意怎么样？\\
\emph{Shēngyi zěnmeyàng ?}\\
« Comment vont les affaires ? »\\
B : 很不错！\\
\emph{Hěn búcuò !}\\
« Très bien ! »
\end{textebox}

\textbf{Observation.} Ce deuxième dialogue utilise deux structures déjà connues mais réunies ici dans un contexte professionnel : la question avec \zh{什么} \emph{shénme} (« quoi, quel ») et la question d'appréciation avec \zh{怎么样} \emph{zěnmeyàng} (« comment »).

\begin{propriete}[Présenter le métier de quelqu'un]
Structure : \cle{Personne + 是 + métier} puis \cle{Personne + 在 + lieu + verbe d'activité}.\\
\zh{他是商人。} \emph{Tā shì shāngrén.} « Il est commerçant. »\\
\zh{她在雅温得工作。} \emph{Tā zài Yǎwēndé gōngzuò.} « Elle travaille à Yaoundé. »\\
\zh{我妈妈在医院工作。} \emph{Wǒ māma zài yīyuàn gōngzuò.} « Ma mère travaille à l'hôpital. »\\
Pour demander le métier : \zh{……做什么工作？} \emph{... zuò shénme gōngzuò ?} « … fait quel travail ? »\\
Attention à la place du complément de lieu : \zh{在 + lieu} se place \emph{avant} le verbe, jamais après : on dit \zh{在杜阿拉做生意}, pas \zh{做生意在杜阿拉}.
\end{propriete}

\begin{center}
\begin{tabularx}{\linewidth}{|X|X|X|X|}
\hline
\rowcolor{popblueL} Caractères & Pinyin & Nature & Traduction \\
\hline
商人 & shāngrén & nom & commerçant, homme d'affaires \\
\hline
做生意 & zuò shēngyi & verbe + compl. & faire des affaires \\
\hline
生意 & shēngyi & nom & affaires, commerce \\
\hline
工作 & gōngzuò & nom / verbe & travail / travailler \\
\hline
怎么样 & zěnmeyàng & pronom interrog. & comment \\
\hline
不错 & búcuò & adjectif & pas mal, bien \\
\hline
多少钱 & duōshao qián & expression & combien (prix) \\
\hline
便宜 & piányi & adjectif & bon marché \\
\hline
贵 & guì & adjectif & cher \\
\hline
买 & mǎi & verbe & acheter \\
\hline
\end{tabularx}
\end{center}

\courssub{Lecture guidée : de l'écoute à la prise de parole}

Un dialogue ne s'apprend pas seulement avec les yeux, mais avec les oreilles et la bouche. Voici la progression en trois temps utilisée en classe :

\begin{enumerate}
\item \textbf{Lecture par l'enseignant.} Écoutez sans lire, livre fermé si possible. Concentrez-vous sur les tons : \zh{wǔshí} (3\textsuperscript{e} ton + 2\textsuperscript{e} ton), \zh{piányi} (2\textsuperscript{e} ton + 1\textsuperscript{er} ton), \zh{mǎi} (3\textsuperscript{e} ton).
\item \textbf{Lecture en chœur.} Toute la classe répète ensemble, phrase par phrase, en imitant le modèle. Le chœur rassure : personne n'est exposé seul.
\item \textbf{Lecture à deux.} En binôme, l'un joue le client (A), l'autre le vendeur (B), puis on échange les rôles. Regardez votre partenaire, pas seulement le texte !
\end{enumerate}

\begin{methode}[Mémoriser un dialogue en cinq étapes]
\begin{enumerate}
\item \textbf{Écouter} le dialogue plusieurs fois, sans texte, pour saisir le rythme et les tons.
\item \textbf{Répéter phrase par phrase} après le modèle, en veillant aux tons avant la vitesse.
\item \textbf{Jouer avec le texte} sous les yeux, en binôme, en soignant la prononciation.
\item \textbf{Jouer sans le texte} : ne gardez que l'organigramme des étapes (saluer → prix → marchander → conclure).
\item \textbf{Remplacer} les mots par vos propres idées : d'autres produits, d'autres prix, d'autres métiers. C'est l'étape où le dialogue devient vraiment le vôtre.
\end{enumerate}
\end{methode}

\courssub{Repérage des formules fonctionnelles}

Chaque dialogue cache des « outils de communication » réutilisables dans des centaines de situations. Apprenez à les repérer :

\begin{center}
\begin{tabularx}{\linewidth}{|X|X|X|}
\hline
\rowcolor{popblueL} Fonction & Formule (caractères + pinyin) & Traduction \\
\hline
Saluer & 你好！Nǐ hǎo ! & Bonjour ! \\
\hline
Demander le prix & 这个多少钱？Zhège duōshao qián ? & Combien coûte ceci ? \\
\hline
Réagir à un prix & 太贵了！Tài guì le ! & Trop cher ! \\
\hline
Négocier poliment & 便宜一点儿吧。Piányi yìdiǎnr ba. & Un peu moins cher, svp. \\
\hline
Accepter & 好吧。Hǎo ba. & D'accord. \\
\hline
Conclure un achat & 我买了。Wǒ mǎi le. & Je le prends. \\
\hline
Demander le métier & 做什么工作？zuò shénme gōngzuò ? & fait quel travail ? \\
\hline
Demander une appréciation & 怎么样？zěnmeyàng ? & comment ? \\
\hline
Donner un avis positif & 很不错！Hěn búcuò ! & Très bien ! \\
\hline
\end{tabularx}
\end{center}

\begin{aretenir}[La particule 吧, votre alliée de la politesse]
\zh{吧} \emph{ba} (ton neutre) adoucit une demande ou une suggestion. Comparez :\\
\zh{便宜一点儿。} \emph{Piányi yìdiǎnr.} « Moins cher. » (sec, presque impoli)\\
\zh{便宜一点儿吧。} \emph{Piányi yìdiǎnr ba.} « Un peu moins cher, s'il vous plaît. » (poli)\\
Au marché, cette petite particule fait toute la différence !
\end{aretenir}

\courssub{Substitution : adapter les dialogues au contexte camerounais}

Le vrai test de compréhension, c'est de remplacer les éléments du dialogue par votre réalité. Voici le matériel de substitution :

\begin{center}
\begin{tabularx}{\linewidth}{|X|X|X|X|}
\hline
\rowcolor{popblueL} Caractères & Pinyin & Nature & Traduction \\
\hline
山药 & shānyao & nom & igname \\
\hline
花生 & huāshēng & nom & arachide, cacahuète \\
\hline
布 & bù & nom & tissu, pagne \\
\hline
芒果 & mángguǒ & nom & mangue \\
\hline
鸡 & jī & nom & poulet, poule \\
\hline
公斤 & gōngjīn & nom de mesure & kilogramme \\
\hline
\end{tabularx}
\end{center}

\begin{exemplebox}[Dialogue 1 transformé : au marché de Yaoundé]
A : 你好！这个山药多少钱？\\
\emph{Nǐ hǎo ! Zhège shānyao duōshao qián ?}\\
« Bonjour ! Combien coûte cet igname ? »\\
B : 十块。\\
\emph{Shí kuài.}\\
« Dix yuans. »\\
A : 太贵了！便宜一点儿吧。\\
\emph{Tài guì le ! Piányi yìdiǎnr ba.}\\
« Trop cher ! Un peu moins cher, s'il vous plaît. »\\
B : 好吧，八块。\\
\emph{Hǎo ba, bā kuài.}\\
« D'accord, huit yuans. »\\
A : 好，我买了。\\
\emph{Hǎo, wǒ mǎi le.}\\
« Bien, je le prends. »
\end{exemplebox}

Même principe pour le dialogue 2 : remplacez \zh{商人} par \zh{老师} \emph{lǎoshī} (enseignant), \zh{医生} \emph{yīshēng} (médecin) ou \zh{司机} \emph{sījī} (chauffeur), et \zh{杜阿拉} par \zh{雅温得} \emph{Yǎwēndé} (Yaoundé) ou \zh{布埃亚} \emph{Bù'āiyà} (Buea).

\begin{exoresolu}[Reconstituer et transformer un dialogue]
Énoncé : remettez les phrases suivantes dans l'ordre logique d'une négociation, puis remplacez le produit par « l'arachide » (\zh{花生}).\\
1. 好，我买了。 2. 这个多少钱？ 3. 好吧，三十块。 4. 太贵了！便宜一点儿吧。 5. 你好！ 6. 三十五块。
\tcblower
Solution détaillée :\\
Ordre logique (selon l'organigramme : saluer → prix → marchander → conclure) : 5 → 2 → 6 → 4 → 3 → 1.\\
Dialogue reconstitué et transformé :\\
A : 你好！这个花生多少钱？\emph{Nǐ hǎo ! Zhège huāshēng duōshao qián ?} « Bonjour ! Combien coûte cette arachide ? »\\
B : 三十五块。\emph{Sānshíwǔ kuài.} « Trente-cinq yuans. »\\
A : 太贵了！便宜一点儿吧。\emph{Tài guì le ! Piányi yìdiǎnr ba.} « Trop cher ! Un peu moins cher, s'il vous plaît. »\\
B : 好吧，三十块。\emph{Hǎo ba, sānshí kuài.} « D'accord, trente yuans. »\\
A : 好，我买了。\emph{Hǎo, wǒ mǎi le.} « Bien, je le prends. »\\
Astuce : le produit se place toujours entre \zh{这个} et \zh{多少钱} ; le prix suit simplement le schéma nombre + \zh{块}.
\end{exoresolu}

\begin{attention}[Pièges fréquents des francophones dans ces dialogues]
\begin{itemize}
\item \textbf{Ton de 买 et 卖} : \zh{买} \emph{mǎi} (3\textsuperscript{e} ton, acheter) et \zh{卖} \emph{mài} (4\textsuperscript{e} ton, vendre) ne se distinguent que par le ton. Une erreur de ton inverse complètement le sens !
\item \textbf{Ordre de la question de prix} : on dit \zh{这个多少钱？} (ceci + combien + argent), jamais \zh{多少钱这个？} — le sujet reste en tête de phrase.
\item \textbf{块 et non 元 à l'oral} : dans la langue parlée, utilisez \zh{块} \emph{kuài} ; \zh{元} \emph{yuán} est réservé à l'écrit et aux contextes formels.
\item \textbf{Ne pas traduire « s'il vous plaît » mot à mot} : \zh{吧} suffit à rendre la demande polie ; n'ajoutez pas \zh{请} partout.
\item \textbf{了 dans 我买了} : \zh{了} marque ici la décision accomplie (« je le prends »), pas un passé composé français à traduire littéralement.
\end{itemize}
\end{attention}

\begin{exercice}[À vous de jouer]
En binôme, construisez un dialogue de négociation complet (au moins six répliques) sur l'un des produits suivants : \zh{布} (tissu), \zh{芒果} (mangue), \zh{鸡} (poulet). Respectez les quatre étapes de l'organigramme, utilisez au moins une fois \zh{一点儿} et la particule \zh{吧}, puis jouez-le devant la classe sans regarder vos notes.
\end{exercice}

Ces deux dialogues modèles sont désormais votre boîte à outils : vous savez saluer, demander un prix, marchander, conclure, et parler du travail de vos proches. La section suivante approfondira les formules fonctionnelles pour présenter, donner votre avis et relancer la conversation dans un débat.

\coursec{Formules fonctionnelles : présenter, opiner, relancer}

Dans la section précédente, nous avons écouté et lu des dialogues modèles : vous avez entendu des élèves qui parlaient du marché de Mokolo et du commerce entre le Cameroun et la Chine. Vous avez remarqué que certaines phrases revenaient tout le temps : \zh{我觉得…}, \zh{首先…}, \zh{你呢？} Ces phrases ne sont pas du vocabulaire nouveau : ce sont des \cle{formules fonctionnelles}, c'est-à-dire des outils tout prêts qui permettent de prendre la parole, de la garder et de la passer à un camarade. Un bon orateur chinois, comme un bon orateur français, ne cherche pas ses transitions : il les connaît par cœur. Cette section vous donne la « boîte à outils » complète du débatteur, en trois compartiments : \cle{présenter}, \cle{opiner}, \cle{réagir et relancer}.

\courssub{Observer d'abord : un mini-exposé modèle}

Avant toute règle, lisons un petit exposé d'élève. Soulignez mentalement les formules que vous reconnaissez.

\begin{textebox}[L'exposé d'Aminatou, élève de Première A à Yaoundé]
大家好，我今天谈谈喀麦隆的经济。\\
Dàjiā hǎo, wǒ jīntiān tántan Kāmàilóng de jīngjì.\\
« Bonjour à tous, aujourd'hui je vais parler de l'économie du Cameroun. »\\[4pt]
首先，我介绍农业：喀麦隆有很多农民，农业非常重要。\\
Shǒuxiān, wǒ jièshào nóngyè : Kāmàilóng yǒu hěn duō nóngmín, nóngyè fēicháng zhòngyào.\\
« D'abord, je présente l'agriculture : le Cameroun compte beaucoup de paysans, l'agriculture est très importante. »\\[4pt]
然后，我谈谈商业。杜阿拉是一个很大的商业城市。\\
Ránhòu, wǒ tántan shāngyè. Dù'ālā shì yí ge hěn dà de shāngyè chéngshì.\\
« Ensuite, je parle du commerce. Douala est une très grande ville commerciale. »\\[4pt]
最后，我说说我的想法。我觉得年轻人应该学习经济。\\
Zuìhòu, wǒ shuōshuo wǒ de xiǎngfǎ. Wǒ juéde niánqīngrén yīnggāi xuéxí jīngjì.\\
« Enfin, je donne mon idée. Je pense que les jeunes devraient étudier l'économie. »\\[4pt]
我的讲话完了，谢谢大家！\\
Wǒ de jiǎnghuà wán le, xièxie dàjiā !\\
« Mon exposé est terminé, merci à tous ! »
\end{textebox}

Observation : l'exposé fait cinq phrases, et chaque phrase commence par une formule-cadre : \zh{我今天谈谈…}, \zh{首先}, \zh{然后}, \zh{最后}, \zh{谢谢大家}. Le contenu (agriculture, commerce, opinion) est simple ; ce qui rend l'exposé solide, c'est l'armature. Retenez ce principe : \textbf{en expression orale, la structure compte autant que le contenu}.

\courssub{Présenter un exposé : ouvrir et ordonner}

\begin{propriete}[Ouvrir un exposé : 大家好，我今天谈谈…]
Structure : \textbf{大家好，我今天谈谈 + thème + 。}\\
\zh{大家} (dàjiā) signifie « tout le monde » ; \zh{谈谈} (tántan) est le verbe \zh{谈} « parler de » redoublé : le redoublement adoucit le verbe et signifie « parler un peu de ». C'est la formule d'ouverture standard d'un exposé scolaire en Chine.\\
Variantes : \zh{我给大家介绍…} (wǒ gěi dàjiā jièshào…, « je vous présente… »), \zh{今天的题目是…} (jīntiān de tímù shì…, « le sujet d'aujourd'hui est… »).
\end{propriete}

\begin{exemplebox}[Ouvrir et ordonner ses idées]
\zh{大家好，我今天谈谈中国的经济。} \emph{Dàjiā hǎo, wǒ jīntiān tántan Zhōngguó de jīngjì.} « Bonjour à tous, aujourd'hui je parle de l'économie chinoise. »\\[4pt]
\zh{首先，我介绍喀麦隆的经济。} \emph{Shǒuxiān, wǒ jièshào Kāmàilóng de jīngjì.} « D'abord, je présente l'économie du Cameroun. »\\[4pt]
\zh{然后，我谈谈商业。} \emph{Ránhòu, wǒ tántan shāngyè.} « Ensuite, je parle du commerce. »\\[4pt]
\zh{最后，我说说我的想法。} \emph{Zuìhòu, wǒ shuōshuo wǒ de xiǎngfǎ.} « Enfin, je donne mon avis. »
\end{exemplebox}

\begin{aretenir}[La chaîne d'ordre : 首先… 然后… 最后…]
\zh{首先} (shǒuxiān) = « d'abord, en premier » ; \zh{然后} (ránhòu) = « ensuite » ; \zh{最后} (zuìhòu) = « enfin ». Ces trois mots se placent \textbf{en tête de phrase}, suivis d'une virgule chinoise \zh{，}. Pour plus de trois idées, on peut enchaîner plusieurs \zh{然后}. Contrairement au français (« premièrement, deuxièmement »), le chinois n'utilise pas de numérotation lourde dans un exposé oral simple.
\end{aretenir}

\courssub{Donner son avis : 我觉得，我认为，对我来说，在我看来}

\begin{propriete}[Règle de 我觉得 / 我认为 : pas de « que » en chinois]
\zh{我觉得} (wǒ juéde) et \zh{我认为} (wǒ rènwéi) signifient tous deux « je pense que », mais ils introduisent \textbf{directement une phrase complète}, sans aucun mot de liaison :\\
\textbf{我觉得 / 我认为 + Sujet + Verbe (phrase complète)}\\
\zh{我觉得中国经济很强。} \emph{Wǒ juéde Zhōngguó jīngjì hěn qiáng.} « Je pense que l'économie chinoise est forte. »\\
Il n'y a \textbf{aucun mot} entre \zh{觉得} et le sujet de la subordonnée (\zh{中国经济}) : le chinois n'a pas d'équivalent du « que » français.\\
Nuance de registre : \zh{觉得} est oral et courant ; \zh{认为} est plus formel, plus écrit, et convient bien aux débats et aux exposés.
\end{propriete}

\begin{exemplebox}[Quatre façons de donner son avis]
\zh{我觉得农业很重要。} \emph{Wǒ juéde nóngyè hěn zhòngyào.} « Je pense que l'agriculture est très importante. »\\[4pt]
\zh{我认为年轻人应该学习汉语。} \emph{Wǒ rènwéi niánqīngrén yīnggāi xuéxí Hànyǔ.} « Je considère que les jeunes devraient apprendre le chinois. »\\[4pt]
\zh{对我来说，农业最重要。} \emph{Duì wǒ láishuō, nóngyè zuì zhòngyào.} « Pour moi, l'agriculture est ce qu'il y a de plus important. »\\[4pt]
\zh{在我看来，中国是一个很大的市场。} \emph{Zài wǒ kànlái, Zhōngguó shì yí ge hěn dà de shìchǎng.} « À mon avis, la Chine est un très grand marché. »
\end{exemplebox}

\begin{attention}[对我来说 et 在我看来 : attention à la place]
\zh{对我来说} (duì wǒ láishuō, « pour moi ») et \zh{在我看来} (zài wǒ kànlái, « à mon avis ») se placent \textbf{en tête de phrase}, suivis d'une virgule \zh{，}. Ne les mettez pas à la fin comme en français (« c'est important, pour moi ») : dire \zh{农业很重要，对我来说} est incorrect dans un exposé organisé. Dites : \zh{对我来说，农业很重要。}
\end{attention}

\courssub{Réagir : l'accord et le désaccord poli}

Dans un débat, il ne suffit pas de donner son avis : il faut réagir à celui des autres. Le chinois dispose de formules courtes et très fréquentes.

\begin{center}
\begin{tabularx}{\linewidth}{|X|X|X|X|}
\hline
\rowcolor{popblueL} Caractères & Pinyin & Fonction & Traduction \\
\hline
我同意 & wǒ tóngyì & accord & je suis d'accord \\
\hline
你说得对 & nǐ shuō de duì & accord + compliment & tu as raison \\
\hline
我也是这么想的 & wǒ yě shì zhème xiǎng de & accord chaleureux & je pense la même chose \\
\hline
我不同意 & wǒ bù tóngyì & désaccord direct & je ne suis pas d'accord \\
\hline
你说得对，但是… & nǐ shuō de duì, dànshì… & désaccord poli & tu as raison, mais… \\
\hline
对不起，我的看法不一样 & duìbuqǐ, wǒ de kànfǎ bù yíyàng & désaccord très poli & pardon, mon avis est différent \\
\hline
\end{tabularx}
\end{center}

\begin{attention}[Erreur fréquente des francophones : 我是同意]
Le français dit « je \emph{suis} d'accord », et beaucoup d'élèves traduisent mot à mot : \zh{我是同意}. C'est une \textbf{erreur} : en chinois, on ne met jamais \zh{是} devant un verbe. \zh{同意} est un verbe à part entière ; on dit simplement \zh{我同意} (wǒ tóngyì), et au négatif \zh{我不同意} (wǒ bù tóngyì). De même, « j'ai raison » se dit \zh{我对} et non \zh{我是对}. Retenez : \textbf{是 ne s'emploie que devant un nom}, pas devant un verbe ni un adjectif.
\end{attention}

Notez la stratégie de politesse chinoise, très proche de la nôtre : avant de contredire, on commence par reconnaître la valeur de l'argument adverse avec \zh{你说得对，但是…} (tu as raison, mais…). Le désaccord brutal \zh{我不同意} est correct grammaticalement, mais à réserver aux débats vifs entre amis ; en classe, préférez la formule amortie.

\begin{exemplebox}[Un échange de débat]
\zh{甲：我觉得中国市场对喀麦隆很重要。} \emph{Jiǎ : Wǒ juéde Zhōngguó shìchǎng duì Kāmàilóng hěn zhòngyào.} « A : Je pense que le marché chinois est très important pour le Cameroun. »\\[4pt]
\zh{乙：我同意。你说得对。} \emph{Yǐ : Wǒ tóngyì. Nǐ shuō de duì.} « B : Je suis d'accord. Tu as raison. »\\[4pt]
\zh{丙：你说得对，但是我觉得喀麦隆也应该发展自己的工业。} \emph{Bǐng : Nǐ shuō de duì, dànshì wǒ juéde Kāmàilóng yě yīnggāi fāzhǎn zìjǐ de gōngyè.} « C : Tu as raison, mais je pense que le Cameroun devrait aussi développer sa propre industrie. »
\end{exemplebox}

\courssub{Relancer la parole et conclure}

Un bon animateur de débat fait parler les autres. Quatre formules suffisent :

\begin{center}
\begin{tabularx}{\linewidth}{|X|X|X|X|}
\hline
\rowcolor{popgreenL} Caractères & Pinyin & Fonction & Traduction \\
\hline
你呢？ & nǐ ne ? & renvoyer la question & et toi ? \\
\hline
你怎么看？ & nǐ zěnme kàn ? & demander un avis & qu'en penses-tu ? \\
\hline
为什么？ & wèishénme ? & demander une raison & pourquoi ? \\
\hline
你能举个例子吗？ & nǐ néng jǔ ge lìzi ma ? & demander un exemple & peux-tu donner un exemple ? \\
\hline
\end{tabularx}
\end{center}

\begin{propriete}[La particule 呢 dans 你呢？]
\zh{你呢？} (nǐ ne ?) signifie « et toi ? » : la particule \zh{呢} renvoie la question précédente à une nouvelle personne, sans répéter le verbe. Après \zh{我觉得农业很重要。你呢？}, le camarade doit donner son avis. C'est la relance la plus courte et la plus naturelle du chinois oral. \zh{你怎么看？} (littéralement « comment vois-tu ? ») est la forme complète pour demander explicitement une opinion.
\end{propriete}

Pour conclure un exposé, une seule formule à mémoriser, utilisée par tous les écoliers chinois :

\begin{exemplebox}[Conclure]
\zh{我的讲话完了，谢谢大家！} \emph{Wǒ de jiǎnghuà wán le, xièxie dàjiā !} « Mon exposé est terminé, merci à tous ! »\\
Variante plus simple : \zh{我说完了，谢谢！} \emph{Wǒ shuō wán le, xièxie !} « J'ai fini, merci ! »
\end{exemplebox}

\courssub{Schéma récapitulatif : la boîte à outils du débatteur}

\begin{popfigure}
\begin{tikzpicture}[
  box/.style={draw=popblue, thick, rounded corners=6pt, fill=popblueL, align=center, minimum width=3.6cm, minimum height=2.2cm, font=\small},
  arr/.style={-{Stealth[length=3mm]}, thick, poporange}
]
\node[box, align=center] (a) at (0,0){\textbf{介绍 jièshào}\\présenter\\大家好，我今天谈谈…\\首先… 然后… 最后…};
\node[box, fill=popgreenL, draw=popgreen, align=center] (b) at (5.6,0){\textbf{意见 yìjiàn}\\opiner\\我觉得… / 我认为…\\对我来说… / 在我看来…};
\node[box, fill=poppinkL, draw=poppink, align=center] (c) at (11.2,0){\textbf{反应 fǎnyìng}\\réagir et relancer\\我同意 / 你说得对，但是…\\你呢？ 你怎么看？};
\draw[arr] (a) -- (b);
\draw[arr] (b) -- (c);
\draw[arr, dashed] (c.south) .. controls (5.6,-2.2) .. (a.south);
\node[font=\footnotesize\itshape, text=popmuted] at (5.6,-2.4) {on reboucle : relancer, c'est faire reparler};
\end{tikzpicture}
\legende{La boîte à outils du débatteur : trois compartiments reliés — présenter, opiner, réagir et relancer.}
\end{popfigure}

\courssub{Méthode : structurer un mini-exposé en cinq phrases}

\begin{methode}[Le plan « 1 + 3 + 1 » pour un mini-exposé oral]
\begin{enumerate}
\item \textbf{Une phrase d'introduction} avec la formule d'ouverture : \zh{大家好，我今天谈谈…} (annoncez le thème, rien de plus).
\item \textbf{Trois idées ordonnées} avec la chaîne \zh{首先… 然后… 最后…} : une idée par phrase, phrase courte, vocabulaire connu. Exemple de squelette : \zh{首先，我介绍…。然后，我谈谈…。最后，我说说我的想法。}
\item \textbf{Une phrase de conclusion} : \zh{我的讲话完了，谢谢大家！}
\item \textbf{Préparez la relance} : après l'exposé, soyez prêt à réagir (\zh{我同意} / \zh{你说得对，但是…}) et à interroger (\zh{你怎么看？}).
\item \textbf{Soignez les tons} en parlant lentement : \zh{shǒuxiān} (3-1), \zh{ránhòu} (2-4), \zh{zuìhòu} (4-4), \zh{juéde} (2-léger).
\end{enumerate}
\end{methode}

\courssub{Exercice résolu et entraînement}

\begin{exoresolu}[Remettre un exposé dans l'ordre]
Énoncé — Voici les phrases de l'exposé de Paul, mélangées. Remettez-les dans l'ordre logique et traduisez la phrase d'ouverture.\\
a) \zh{最后，我觉得商业也很重要。}\\
b) \zh{首先，我介绍农业。}\\
c) \zh{我的讲话完了，谢谢大家！}\\
d) \zh{大家好，我今天谈谈喀麦隆的经济。}\\
e) \zh{然后，我谈谈工业。}
\tcblower
Solution détaillée — On repère les marqueurs d'ordre : \zh{首先} (d'abord), \zh{然后} (ensuite), \zh{最后} (enfin), plus la formule d'ouverture \zh{大家好，我今天谈谈…} et la formule de clôture \zh{谢谢大家}. L'ordre correct est donc : \textbf{d → b → e → a → c}.\\
Traduction de l'ouverture : \zh{大家好，我今天谈谈喀麦隆的经济。} \emph{Dàjiā hǎo, wǒ jīntiān tántan Kāmàilóng de jīngjì.} « Bonjour à tous, aujourd'hui je vais parler de l'économie du Cameroun. »\\
L'exposé reconstitué : introduction (d), idée 1 : l'agriculture (b), idée 2 : l'industrie (e), idée 3 : avis personnel sur le commerce (a), conclusion (c). C'est exactement le plan « 1 + 3 + 1 ».
\end{exoresolu}

\begin{exercice}[À votre tour : réagir et relancer]
\begin{enumerate}
\item Votre camarade dit : \zh{我觉得足球对喀麦隆的经济很重要。} Réagissez d'abord en étant d'accord (deux formules différentes), puis une troisième fois en exprimant un désaccord poli avec \zh{但是}.
\item Complétez les relances : a) \zh{我觉得汉语很有用。……？} (demandez l'avis de votre voisin) ; b) \zh{……？} (demandez un exemple) ; c) \zh{……？} (demandez pourquoi).
\item Rédigez le squelette d'un mini-exposé sur « le marché de ma ville » avec le plan « 1 + 3 + 1 » : une phrase d'ouverture, trois phrases avec \zh{首先 / 然后 / 最后}, une phrase de conclusion.
\end{enumerate}
\end{exercice}

\begin{corrige}[Exercice « Réagir et relancer »]
\begin{enumerate}
\item Accord : \zh{我同意。} (wǒ tóngyì) ; \zh{你说得对。} (nǐ shuō de duì) ou \zh{我也是这么想的。} (wǒ yě shì zhème xiǎng de). Désaccord poli : \zh{你说得对，但是我觉得农业更重要。} \emph{Nǐ shuō de duì, dànshì wǒ juéde nóngyè gèng zhòngyào.} « Tu as raison, mais je pense que l'agriculture est plus importante. »
\item a) \zh{你呢？} (nǐ ne ?) ou \zh{你怎么看？} (nǐ zěnme kàn ?) ; b) \zh{你能举个例子吗？} (nǐ néng jǔ ge lìzi ma ?) ; c) \zh{为什么？} (wèishénme ?)
\item Modèle : \zh{大家好，我今天谈谈我们城市的市场。首先，我介绍市场的商品。然后，我谈谈商人的生活。最后，我说说我的想法：我觉得市场对经济很重要。我的讲话完了，谢谢大家！} \emph{Dàjiā hǎo, wǒ jīntiān tántan wǒmen chéngshì de shìchǎng. Shǒuxiān, wǒ jièshào shìchǎng de shāngpǐn. Ránhòu, wǒ tántan shāngrén de shēnghuó. Zuìhòu, wǒ shuōshuo wǒ de xiǎngfǎ : wǒ juéde shìchǎng duì jīngjì hěn zhòngyào. Wǒ de jiǎnghuà wán le, xièxie dàjiā !} « Bonjour à tous, aujourd'hui je parle du marché de notre ville. D'abord, je présente les marchandises du marché. Ensuite, je parle de la vie des commerçants. Enfin, je donne mon avis : je pense que le marché est très important pour l'économie. Mon exposé est terminé, merci à tous ! »
\end{enumerate}
\end{corrige}

\begin{aretenir}[L'essentiel de la section]
\begin{itemize}
\item Ouvrir : \zh{大家好，我今天谈谈…} ; ordonner : \zh{首先… 然后… 最后…} ; conclure : \zh{我的讲话完了，谢谢大家！}
\item Opiner : \zh{我觉得 / 我认为 + phrase complète}, sans mot pour « que » ; \zh{对我来说…} et \zh{在我看来…} en tête de phrase.
\item Réagir : \zh{我同意} (jamais \zh{我是同意} !) ; désaccord poli : \zh{你说得对，但是…}
\item Relancer : \zh{你呢？} \zh{你怎么看？} \zh{为什么？} \zh{你能举个例子吗？}
\end{itemize}
\end{aretenir}

Ces formules sont maintenant vos outils. Dans la prochaine section, nous travaillerons la prononciation : les tons de ces formules, pour que votre débat soit non seulement bien structuré, mais aussi bien prononcé.

\coursec{Les tons à travailler}

Dans la section précédente, nous avons appris les formules fonctionnelles pour présenter, donner son avis et relancer un débat sur l'économie. Mais une belle formule mal prononcée peut devenir incompréhensible : en chinois, chaque syllabe porte un \cle{ton}, c'est-à-dire une mélodie qui fait partie du mot, au même titre que ses consonnes et ses voyelles. Dire \zh{买} (mǎi, acheter) au lieu de \zh{卖} (mài, vendre), c'est inverser tout le sens d'une négociation au marché ! Cette section est donc consacrée à l'entraînement systématique des tons sur le vocabulaire économique de la leçon.

\courssub{Révision rapide : les quatre tons et le ton neutre}

\begin{definition}[Les tons du chinois mandarin]
Le mandarin possède \cle{quatre tons} et un \cle{ton neutre}. Le ton est la courbe mélodique (haute, montante, descendante-montante, chutante) sur laquelle une syllabe est prononcée. Il est noté dans le pinyin par un accent sur la voyelle principale : ā (1\textsuperscript{er} ton), á (2\textsuperscript{e} ton), ǎ (3\textsuperscript{e} ton), à (4\textsuperscript{e} ton), a (ton neutre, léger et court).
\end{definition}

Reprenons les quatre tons avec des mots économiques de notre leçon :

\begin{center}
\begin{tabularx}{\linewidth}{|X|X|X|X|}
\hline
\rowcolor{popblueL} Ton & Mélodie & Mot économique & Traduction \\
\hline
1\textsuperscript{er} ton (ā) & haut et plat, comme une note tenue & \zh{工} gōng & travail, ouvrier \\
\hline
2\textsuperscript{e} ton (á) & montant, comme une question française & \zh{钱} qián & argent \\
\hline
3\textsuperscript{e} ton (ǎ) & descend puis remonte (souvent réduit à la descente) & \zh{买} mǎi & acheter \\
\hline
4\textsuperscript{e} ton (à) & chute brusque, comme un ordre sec & \zh{卖} mài & vendre \\
\hline
Ton neutre (a) & court, léger, sans accent & \zh{吗} ma & particule de question \\
\hline
\end{tabularx}
\end{center}

\begin{popfigure}
\begin{tikzpicture}[scale=1.0, >=Stealth]
% Grille de référence
\foreach \x in {0,4,8,12}{
  \draw[popline, dashed] (\x,0) -- (\x+2.6,0);
  \draw[popline, dashed] (\x,1) -- (\x+2.6,1);
  \draw[popline, dashed] (\x,2) -- (\x+2.6,2);
}
% Ton 1 : horizontal haut
\draw[line width=2.2pt, popblue] (0.2,2) -- (2.4,2);
\node[align=center, popdark] at (1.3,-0.7) {\zh{工} \emph{gōng}\\ 1\textsuperscript{er} ton : haut et plat};
% Ton 2 : montant
\draw[line width=2.2pt, poporange] (4.2,0.6) -- (6.4,2);
\node[align=center, popdark] at (5.3,-0.7) {\zh{钱} \emph{qián}\\ 2\textsuperscript{e} ton : montant};
% Ton 3 : descendant-montant
\draw[line width=2.2pt, popgreen] (8.2,1.6) .. controls (8.9,0.1) and (9.7,0.1) .. (10.4,1.8);
\node[align=center, popdark] at (9.3,-0.7) {\zh{买} \emph{mǎi}\\ 3\textsuperscript{e} ton : bas puis remonte};
% Ton 4 : chutant
\draw[line width=2.2pt, poppink] (12.2,2) -- (14.4,0.3);
\node[align=center, popdark] at (13.3,-0.7) {\zh{卖} \emph{mài}\\ 4\textsuperscript{e} ton : chute brusque};
\end{tikzpicture}
\legende{Les quatre tons du mandarin illustrés par quatre mots économiques : \zh{工} gōng (travail), \zh{钱} qián (argent), \zh{买} mǎi (acheter), \zh{卖} mài (vendre).}
\end{popfigure}

\begin{attention}[Le français n'a pas de tons]
En français, la mélodie d'une phrase exprime l'émotion ou la question (« Tu viens ? » monte à la fin), mais elle ne change jamais le sens d'un mot. En chinois, le ton est \textbf{dans} le mot : \zh{mā} (maman), \zh{má} (chanvre), \zh{mǎ} (cheval), \zh{mà} (gronder) sont quatre mots différents. Un francophone a tendance à « chanter » la phrase entière comme en français : il faut au contraire prononcer chaque syllabe avec son propre ton, de façon stable.
\end{attention}

\courssub{Paires minimales économiques : ne pas confondre acheter et vendre}

Deux mots forment une \cle{paire minimale} quand ils ne diffèrent que par un seul élément sonore — ici, uniquement le ton. Ce sont les pièges les plus dangereux à l'oral.

\begin{center}
\begin{tabularx}{\linewidth}{|X|X|X|X|}
\hline
\rowcolor{popblueL} Paire & Pinyin & Ton qui change & Traduction \\
\hline
\zh{买} / \zh{卖} & mǎi / mài & 3\textsuperscript{e} / 4\textsuperscript{e} & acheter / vendre \\
\hline
\zh{资} / \zh{子} & zī / zǐ & 1\textsuperscript{er} / 3\textsuperscript{e} & capital / fils \\
\hline
\zh{商} / \zh{上} & shāng / shàng & 1\textsuperscript{er} / 4\textsuperscript{e} & commerce / monter \\
\hline
\zh{钱} / \zh{千} & qián / qiān & 2\textsuperscript{e} / 1\textsuperscript{er} & argent / mille \\
\hline
\end{tabularx}
\end{center}

\begin{exemplebox}[Au marché de Mokolo : le ton qui change tout]
\zh{我买苹果。} \emph{Wǒ mǎi píngguǒ.} « J'achète des pommes. »\\
\zh{我卖苹果。} \emph{Wǒ mài píngguǒ.} « Je vends des pommes. »\\[4pt]
\zh{太贵了！} \emph{Tài guì le !} « C'est trop cher ! » (guì, 4\textsuperscript{e} ton sec)\\
\zh{很便宜。} \emph{Hěn piányi.} « C'est bon marché. » (pián, 2\textsuperscript{e} ton ; yi au ton neutre)
\end{exemplebox}

Remarquez que dans \zh{便宜} piányi, la deuxième syllabe perd son ton d'origine et devient \cle{neutre} : c'est fréquent dans les mots de deux syllabes très courants. À l'oral, dites \emph{pián-yi}, léger et rapide sur la fin.

\courssub{Le sandhi du 3\textsuperscript{e} ton : la règle d'or de la prononciation}

Observons d'abord ces exemples, prononcés à voix haute par l'enseignant :

\begin{exemplebox}[Écoutez la différence entre l'écrit et le prononcé]
\zh{很好} s'écrit \emph{hěn hǎo} mais se prononce \emph{hén hǎo} « très bien ».\\
\zh{我可以} s'écrit \emph{wǒ kěyǐ} mais se prononce \emph{wó kěyǐ} « je peux ».\\
\zh{我想买} s'écrit \emph{wǒ xiǎng mǎi} mais se prononce \emph{wó xiáng mǎi} « je voudrais acheter ».
\end{exemplebox}

\begin{propriete}[Règle du sandhi du 3\textsuperscript{e} ton]
Quand \textbf{deux syllabes au 3\textsuperscript{e} ton se suivent}, la première se prononce comme un \textbf{2\textsuperscript{e} ton} (montant). On continue à \emph{écrire} le pinyin avec le ton d'origine, mais on \emph{prononce} le ton modifié.\\
Structure : 3\textsuperscript{e} + 3\textsuperscript{e} → prononcé 2\textsuperscript{e} + 3\textsuperscript{e}.\\
Exemple : \zh{很好} hěn hǎo → prononcé \emph{hén hǎo}.
\end{propriete}

Cette règle a une raison simple : enchaîner deux courbes « descente-montée » est long et fatigant pour la bouche ; la langue simplifie naturellement la première en une montée. C'est un phénomène automatique : un Chinois qui parle vite applique le sandhi sans y penser, et vous devez en faire autant.

\begin{attention}[Le demi-3\textsuperscript{e} ton : l'autre visage du 3\textsuperscript{e} ton]
Le 3\textsuperscript{e} ton complet (descente puis montée) n'est prononcé ainsi qu'\textbf{isolément} ou \textbf{en fin de phrase}. Devant un 1\textsuperscript{er}, un 2\textsuperscript{e}, un 4\textsuperscript{e} ton ou un ton neutre, on ne fait que la \textbf{descente}, basse et courte : c'est le \cle{demi-3\textsuperscript{e} ton}.\\
Exemple : \zh{我买} \emph{wǒ mǎi} — ici wǒ devant mǎi suit la règle du sandhi (wó mǎi), mais dans \zh{我说} \emph{wǒ shuō} « je parle », wǒ devant shuō (1\textsuperscript{er} ton) se prononce simplement bas, sans remonter. Erreur typique du francophone : remonter exagérément chaque 3\textsuperscript{e} ton, ce qui donne un parler saccadé et artificiel.
\end{attention}

\courssub{Suites difficiles : quatre 3\textsuperscript{e} tons d'affilée}

La phrase-phare de cette leçon est un excellent entraînement :

\begin{exemplebox}[La phrase à quatre 3\textsuperscript{e} tons]
\zh{我也想买中国产品。}\\
\emph{Wǒ yě xiǎng mǎi Zhōngguó chǎnpǐn.}\\
« Moi aussi, je voudrais acheter des produits chinois. »\\[4pt]
Découpage en groupes de deux syllabes : \textbf{wǒ yě / xiǎng mǎi} → prononcé \textbf{wó yě / xiáng mǎi}. Dans chaque groupe, seul le premier 3\textsuperscript{e} ton devient un 2\textsuperscript{e} ton ; le second garde son 3\textsuperscript{e} ton (demi-3\textsuperscript{e} ton, car il est suivi d'une autre syllabe). Pour \zh{产品} chǎnpǐn (3\textsuperscript{e} + 3\textsuperscript{e}), on prononce \emph{chán pǐn}.
\end{exemplebox}

\begin{methode}[Comment prononcer une longue suite de 3\textsuperscript{e} tons]
\begin{enumerate}
\item \textbf{Repérez} toutes les syllabes au 3\textsuperscript{e} ton dans la phrase écrite.
\item \textbf{Groupez} par paires logiques (sens + rythme) : wǒ yě / xiǎng mǎi / chǎn pǐn.
\item \textbf{Appliquez} la règle dans chaque paire : le premier ton devient un 2\textsuperscript{e} ton.
\item \textbf{Lisez} lentement, puis de plus en plus vite, sans jamais remonter les demi-3\textsuperscript{e} tons.
\item \textbf{Vérifiez} avec l'enseignant ou un enregistrement.
\end{enumerate}
\end{methode}

\courssub{Le ton de \zh{不} : bù devient bú}

Le mot \zh{不} « ne... pas » a lui aussi son comportement spécial. Observez :

\begin{exemplebox}[Écoutez \zh{不} dans quatre contextes]
\zh{不买} \emph{bù mǎi} « ne pas acheter » (devant un 3\textsuperscript{e} ton : bù reste au 4\textsuperscript{e} ton)\\
\zh{不来} \emph{bù lái} « ne pas venir » (devant un 2\textsuperscript{e} ton : bù reste au 4\textsuperscript{e} ton)\\
\zh{不贵} \emph{bú guì} « pas cher » (devant un 4\textsuperscript{e} ton : bù devient bú, 2\textsuperscript{e} ton)\\
\zh{不对} \emph{bú duì} « ce n'est pas correct » (devant un 4\textsuperscript{e} ton : bú)
\end{exemplebox}

\begin{propriete}[Règle du changement de ton de \zh{不}]
\zh{不} est au \textbf{4\textsuperscript{e} ton} (bù), mais il passe au \textbf{2\textsuperscript{e} ton} (bú) devant une syllabe au \textbf{4\textsuperscript{e} ton}. Devant tous les autres tons (1\textsuperscript{er}, 2\textsuperscript{e}, 3\textsuperscript{e}, neutre), il reste bù.\\
Structure : bù + syllabe au 4\textsuperscript{e} ton → \textbf{bú} + syllabe.\\
Exemples : \zh{不是} bú shì « ne pas être », \zh{不要} bú yào « ne pas vouloir », \zh{不贵} bú guì « pas cher », \zh{不去} bú qù « ne pas aller ».\\
Exception d'écriture : dans les manuels, on écrit souvent le ton d'origine (bù) ; à l'oral, on prononce toujours bú devant un 4\textsuperscript{e} ton.
\end{propriete}

Là encore, la raison est l'euphonie : deux chutes brusques d'affilée (bù guì) sont dures à prononcer ; la langue adoucit la première en montée (bú guì). Comparez avec le français qui transforme « je ne ai » en « je n'ai » pour faciliter la prononciation : même logique de confort.

\begin{exoresolu}[Discrimination des tons en contexte économique]
Énoncé : l'enseignant prononce les cinq phrases suivantes. Pour chacune, indiquez (a) le ton de la syllabe soulignée mentalement, (b) si un changement de ton s'applique, et (c) la traduction.\\
1. \zh{这个东西很贵。} 2. \zh{不贵，很便宜。} 3. \zh{我也想买。} 4. \zh{你卖什么？} 5. \zh{我可以看看吗？}
\tcblower
Solution détaillée :\\
1. \zh{贵} guì : 4\textsuperscript{e} ton, chute brusque ; aucun changement. « Cette chose est très chère. »\\
2. \zh{不贵} : bù devant guì (4\textsuperscript{e} ton) → prononcé \emph{bú guì} ; \zh{便宜} piányi : 2\textsuperscript{e} ton + ton neutre. « Ce n'est pas cher, c'est bon marché. »\\
3. \zh{我也想买} wǒ yě xiǎng mǎi : quatre 3\textsuperscript{e} tons ; sandhi par paires → \emph{wó yě xiáng mǎi}. « Moi aussi je voudrais acheter. »\\
4. \zh{卖} mài : 4\textsuperscript{e} ton ; attention à ne pas le confondre avec \zh{买} mǎi. « Qu'est-ce que tu vends ? »\\
5. \zh{我可以} wǒ kěyǐ : sandhi du 3\textsuperscript{e} ton → \emph{wó kěyǐ}. « Je peux regarder ? »
\end{exoresolu}

\begin{experience}[Jeu des cartes de tons]
Matériel : chaque élève prépare quatre cartes numérotées 1, 2, 3, 4 (et une carte « 0 » pour le ton neutre).\\
Déroulement : l'enseignant prononce un mot économique de la leçon (\zh{工}, \zh{钱}, \zh{买}, \zh{卖}, \zh{贵}, \zh{便宜}, \zh{公司}, \zh{工作}...) ; les élèves lèvent la carte du ton entendu. Deuxième manche : l'enseignant prononce une phrase complète (\zh{我也想买}, \zh{不贵}) et les élèves doivent dire si un changement de ton a eu lieu. L'équipe la plus rapide gagne. Variante : un élève vient prononcer au tableau et la classe lève les cartes.
\end{experience}

\begin{aretenir}[Les trois réflexes de prononciation à acquérir]
\begin{itemize}
\item \textbf{3\textsuperscript{e} + 3\textsuperscript{e}} : le premier devient un 2\textsuperscript{e} ton — \zh{很好} hěn hǎo → hén hǎo.
\item \textbf{3\textsuperscript{e} devant un autre ton} : on ne fait que la descente (demi-3\textsuperscript{e} ton) — \zh{我买} wǒ mǎi.
\item \textbf{bù devant un 4\textsuperscript{e} ton} : il devient bú — \zh{不贵} bú guì, \zh{不是} bú shì.
\end{itemize}
Ces trois règles sont automatiques chez les sinophones : les maîtriser, c'est déjà parler un chinois naturel et compréhensible, condition indispensable pour réussir l'exposé et le débat de la fin de cette leçon.
\end{aretenir}

\coursec{Jeux de rôle : au marché et à l'entreprise}

Dans la section précédente, nous avons travaillé les \cle{tons} : vous savez maintenant prononcer correctement des mots comme \zh{贵} (guì, cher) et \zh{便宜} (piányi, bon marché) sans les confondre. Mais une belle prononciation ne suffit pas : il faut la mettre en \cle{situation réelle}. C'est exactement le but de cette section : parler chinois \emph{pour agir} — négocier un prix, convaincre un employeur. Nous allons jouer deux scènes de la vie économique : le marché et l'entretien d'embauche.

\courssub{Principes généraux du jeu de rôle en classe de chinois}

\begin{definition}[Qu'est-ce qu'un jeu de rôle ?]
Un \cle{jeu de rôle} (\zh{角色扮演} juésè bànyǎn) est une activité orale où chaque élève incarne un personnage défini par une \cle{carte de rôle} (\zh{角色卡} juésè kǎ) : un vendeur, un client, un employeur, un candidat. Chacun connaît des informations que l'autre n'a pas (prix minimum, budget, compétences), ce qui crée un vrai besoin de communiquer en chinois.
\end{definition}

\begin{methode}[Réussir un jeu de rôle en cinq étapes]
\begin{enumerate}
\item \textbf{Lire sa carte de rôle} et souligner ses objectifs (que dois-je obtenir ?).
\item \textbf{Préparer mentalement} ses phrases avec les formules fonctionnelles vues en section 3 : \zh{我觉得……} (Wǒ juéde…), \zh{太贵了！} (Tài guì le!), \zh{能不能便宜一点儿？} (Néng bu néng piányi yìdiǎnr?).
\item \textbf{Parler fort et articuler} les tons : le public doit entendre la différence entre \zh{买} (mǎi, acheter) et \zh{卖} (mài, vendre).
\item \textbf{Regarder son partenaire} dans les yeux : on ne lit pas ses notes, on communique.
\item \textbf{Utiliser au moins 3 formules fonctionnelles} pendant l'échange (présenter, opiner, relancer).
\end{enumerate}
\end{methode}

\begin{popfigure}
\begin{center}
\begin{tikzpicture}[node distance=1.6cm, >=Stealth]
\node[draw=popblue, fill=popblueL, rounded corners, align=center, minimum width=3.2cm, minimum height=1.2cm] (prep) {\textbf{准备 zhǔnbèi}\\Préparation (5 min)};
\node[draw=poporange, fill=poporangeL, rounded corners, align=center, minimum width=3.2cm, minimum height=1.2cm, right=of prep] (jeu) {\textbf{表演 biǎoyǎn}\\Jeu de rôle (3 min)};
\node[draw=popgreen, fill=popgreenL, rounded corners, align=center, minimum width=3.2cm, minimum height=1.2cm, right=of jeu] (eval) {\textbf{评价 píngjià}\\Évaluation par les pairs (2 min)};
\draw[->, very thick, popdark] (prep) -- (jeu);
\draw[->, very thick, popdark] (jeu) -- (eval);
\draw[->, very thick, popmuted, dashed] (eval.south) to[bend right=25] node[below, align=center, font=\small]{nouveau binôme} (prep.south);
\end{tikzpicture}
\end{center}
\legende{Déroulé d'une séance de jeu de rôle : préparation, performance, évaluation.}
\end{popfigure}

\courssub{Jeu de rôle 1 : au marché de Mokolo}

\textbf{Situation.} Nous sommes au marché de Mokolo, à Yaoundé. Un client veut acheter des bananes et du café, mais son budget est limité. Le vendeur veut gagner le plus possible. Qui obtiendra le meilleur prix ?

\begin{textebox}[Carte de rôle : le vendeur 卖家]
你卖香蕉和咖啡。香蕉十块一斤，咖啡三十块一包。最低：香蕉八块，咖啡二十五块。\\
Nǐ mài xiāngjiāo hé kāfēi. Xiāngjiāo shí kuài yì jīn, kāfēi sānshí kuài yì bāo. Zuìdī: xiāngjiāo bā kuài, kāfēi èrshíwǔ kuài.\\
« Tu vends des bananes et du café. Les bananes coûtent dix yuans le jin, le café trente yuans le paquet. Prix minimum : huit yuans pour les bananes, vingt-cinq pour le café. »
\end{textebox}

\begin{textebox}[Carte de rôle : le client 买家]
你想买两斤香蕉和一包咖啡。你只有四十块。\\
Nǐ xiǎng mǎi liǎng jīn xiāngjiāo hé yì bāo kāfēi. Nǐ zhǐyǒu sìshí kuài.\\
« Tu veux acheter deux jin de bananes et un paquet de café. Tu n'as que quarante yuans. »
\end{textebox}

\textbf{Analyse de la situation.} Faisons le calcul : au prix affiché, deux jin de bananes coûtent vingt yuans et un paquet de café trente yuans, soit \cle{cinquante yuans} au total. Or le client n'a que quarante yuans. Il \emph{doit} donc négocier. Au prix minimum du vendeur, le total est de 16 + 25 = 41 yuans : encore un yuan de trop ! La négociation sera serrée, c'est ce qui rend le jeu intéressant.

\begin{propriete}[La structure de la négociation]
Pour marchander, on enchaîne trois mouvements :\\
\textbf{1. Demander le prix} : Sujet + \zh{多少钱} + (quantité) ?\\
\zh{香蕉多少钱一斤？} \emph{Xiāngjiāo duōshao qián yì jīn?} « Combien coûte le jin de bananes ? »\\
\textbf{2. Réagir au prix} : \zh{太 + adjectif + 了！}\\
\zh{太贵了！} \emph{Tài guì le!} « C'est trop cher ! »\\
\textbf{3. Proposer un prix} : \zh{能不能 + verbe + 一点儿？} ou \zh{……块，怎么样？}\\
\zh{能不能便宜一点儿？} \emph{Néng bu néng piányi yìdiǎnr?} « Peux-tu faire un peu moins cher ? »\\
\zh{三十五块，怎么样？} \emph{Sānshíwǔ kuài, zěnmeyàng?} « Trente-cinq yuans, ça te va ? »
\end{propriete}

\begin{exemplebox}[Dialogue modèle à imiter]
—— 老板，香蕉多少钱一斤？\emph{Lǎobǎn, xiāngjiāo duōshao qián yì jīn?} « Patron, combien le jin de bananes ? »\\
—— 十块一斤。\emph{Shí kuài yì jīn.} « Dix yuans le jin. »\\
—— 太贵了！八块行不行？\emph{Tài guì le! Bā kuài xíng bu xíng?} « Trop cher ! Huit yuans, ça marche ? »\\
—— 好吧，八块。咖啡呢？\emph{Hǎo ba, bā kuài. Kāfēi ne?} « D'accord, huit yuans. Et le café ? »\\
—— 咖啡三十块一包。\emph{Kāfēi sānshí kuài yì bāo.} « Le café, trente yuans le paquet. »\\
—— 我只有四十块。两斤香蕉加一包咖啡，四十块，怎么样？\emph{Wǒ zhǐyǒu sìshí kuài. Liǎng jīn xiāngjiāo jiā yì bāo kāfēi, sìshí kuài, zěnmeyàng?} « Je n'ai que quarante yuans. Deux jin de bananes plus un paquet de café, quarante yuans, d'accord ? »\\
—— 好，卖给你！\emph{Hǎo, mài gěi nǐ!} « Bon, je te les vends ! »
\end{exemplebox}

\begin{attention}[Pièges des francophones au marché]
\begin{itemize}
\item Ne confondez pas \zh{买} (mǎi, 3\up{e} ton, acheter) et \zh{卖} (mài, 4\up{e} ton, vendre) : un mauvais ton transforme le client en vendeur !
\item On dit \zh{两斤} (liǎng jīn) et non \zh{二斤} : devant un classificateur, « deux » se dit \zh{两}.
\item Le prix se dit sans verbe : \zh{香蕉十块一斤} (litt. « bananes dix yuans un jin »), pas besoin de traduire « coûter ».
\item \zh{一点儿} signifie « un peu » : \zh{便宜一点儿} = « un peu moins cher », formule indispensable pour négocier poliment.
\end{itemize}
\end{attention}

\begin{savaistu}[Marchander : Cameroun et Chine]
En Chine comme au Cameroun, marchander (\zh{讨价还价} tǎojià huánjià, litt. « discuter le prix, rendre le prix ») est une pratique tout à fait normale au marché : le premier prix annoncé est rarement le prix final. Attention cependant : dans les \cle{supermarchés} (\zh{超市} chāoshì), les prix sont \emph{fixes} dans les deux pays. Marchander au supermarché serait très mal vu ! La règle est simple : au marché on négocie, au magasin on paie le prix affiché.
\end{savaistu}

\courssub{Jeu de rôle 2 : l'entretien d'embauche}

\textbf{Situation.} Une entreprise chinoise installée à Douala cherche un employé. Le directeur chinois (\zh{经理} jīnglǐ) reçoit un candidat camerounais. L'entretien porte sur trois thèmes : le métier, l'expérience et les langues.

\begin{textebox}[Carte de rôle : le directeur 经理]
你是一家中国公司的经理。你要问三个问题：1. 你会做什么工作？2. 你以前做过这个工作吗？3. 你会说什么语言？最后说：我们明天给你打电话。\\
Nǐ shì yì jiā Zhōngguó gōngsī de jīnglǐ. Nǐ yào wèn sān ge wèntí: 1. Nǐ huì zuò shénme gōngzuò? 2. Nǐ yǐqián zuòguo zhège gōngzuò ma? 3. Nǐ huì shuō shénme yǔyán? Zuìhòu shuō: Wǒmen míngtiān gěi nǐ dǎ diànhuà.\\
« Tu es le directeur d'une entreprise chinoise. Pose trois questions : 1. Quel travail sais-tu faire ? 2. As-tu déjà fait ce travail ? 3. Quelles langues parles-tu ? Termine par : nous t'appellerons demain. »
\end{textebox}

\begin{textebox}[Carte de rôle : le candidat 应聘者]
你是喀麦隆人，你想在这家公司工作。你会说法语、英语和一点儿汉语。你以前在一个商店工作过两年。你要介绍自己，回答问题，最后说：谢谢您，我等您的电话。\\
Nǐ shì Kāmàilóng rén, nǐ xiǎng zài zhè jiā gōngsī gōngzuò. Nǐ huì shuō Fǎyǔ, Yīngyǔ hé yìdiǎnr Hànyǔ. Nǐ yǐqián zài yí ge shāngdiàn gōngzuòguo liǎng nián. Nǐ yào jièshào zìjǐ, huídá wèntí, zuìhòu shuō: Xièxie nín, wǒ děng nín de diànhuà.\\
« Tu es Camerounais et tu veux travailler dans cette entreprise. Tu parles français, anglais et un peu chinois. Tu as travaillé deux ans dans un magasin. Présente-toi, réponds aux questions, et termine par : merci, j'attends votre appel. »
\end{textebox}

\begin{propriete}[Parler de l'expérience avec 过]
Pour dire qu'on a \emph{déjà fait} quelque chose (expérience vécue), on place \cle{过} (guo) après le verbe :\\
\textbf{Sujet + Verbe + 过 + (durée / complément)}\\
\zh{我在一个商店工作过两年。} \emph{Wǒ zài yí ge shāngdiàn gōngzuòguo liǎng nián.} « J'ai travaillé deux ans dans un magasin. »\\
\zh{你以前做过这个工作吗？} \emph{Nǐ yǐqián zuòguo zhège gōngzuò ma?} « As-tu déjà fait ce travail ? »\\
À la forme négative : \zh{没(有) + verbe + 过} — \zh{我没做过这个工作。} \emph{Wǒ méi zuòguo zhège gōngzuò.} « Je n'ai jamais fait ce travail. »\\
\textbf{Comparaison avec le français} : le français utilise le passé composé (« j'ai travaillé ») ; le chinois ne conjugue pas, il ajoute simplement la particule 过 après le verbe, qui reste invariable.
\end{propriete}

\begin{center}
\begin{tabularx}{\linewidth}{|X|X|X|X|}
\hline
\rowcolor{popblueL} Caractères & Pinyin & Nature & Traduction \\
\hline
角色扮演 & juésè bànyǎn & loc. verbale & jeu de rôle \\
\hline
讨价还价 & tǎojià huánjià & loc. verbale & marchander \\
\hline
老板 & lǎobǎn & n. & patron, vendeur \\
\hline
经理 & jīnglǐ & n. & directeur, manager \\
\hline
公司 & gōngsī & n. & entreprise, société \\
\hline
工作 & gōngzuò & n. / v. & travail ; travailler \\
\hline
经验 & jīngyàn & n. & expérience \\
\hline
语言 & yǔyán & n. & langue \\
\hline
应聘者 & yìngpìnzhě & n. & candidat (à un emploi) \\
\hline
面试 & miànshì & n. / v. & entretien d'embauche ; passer un entretien \\
\hline
便宜 & piányi & adj. & bon marché \\
\hline
怎么样 & zěnmeyàng & pron. & comment ? qu'en dis-tu ? \\
\hline
\end{tabularx}
\end{center}

\begin{center}
\begin{tabularx}{\linewidth}{|X|X|X|}
\hline
\rowcolor{popblueL} Fonction & Formule chinoise + pinyin & Traduction \\
\hline
Se présenter & 我叫……，我是喀麦隆人。Wǒ jiào…, wǒ shì Kāmàilóng rén. & Je m'appelle…, je suis Camerounais. \\
\hline
Demander le prix & 多少钱？Duōshao qián? & Combien ça coûte ? \\
\hline
Réagir & 太贵了！Tài guì le! & C'est trop cher ! \\
\hline
Proposer & ……块，怎么样？…kuài, zěnmeyàng? & … yuans, ça te va ? \\
\hline
Parler d'expérience & 我工作过两年。Wǒ gōngzuòguo liǎng nián. & J'ai travaillé deux ans. \\
\hline
Conclure poliment & 谢谢您！Xièxie nín! & Merci (formel) ! \\
\hline
\end{tabularx}
\end{center}

\courssub{Grille d'observation et organisation du tour de rôle}

Pendant que deux élèves jouent, les autres ne sont pas spectateurs passifs : ils sont \cle{observateurs} et remplissent une grille d'évaluation. C'est la phase \zh{评价} (píngjià).

\begin{center}
\begin{tabularx}{\linewidth}{|X|X|X|X|}
\hline
\rowcolor{popblueL} Critère & Excellent (2 pts) & Moyen (1 pt) & À retravailler (0 pt) \\
\hline
Tons corrects (买/卖, 十/四) & presque aucune erreur & quelques erreurs & erreurs fréquentes \\
\hline
Formules fonctionnelles (3 minimum) & 3 ou plus utilisées & 1 ou 2 utilisées & aucune \\
\hline
Fluidité (sans lire ses notes) & parle librement & hésite parfois & lit ses notes \\
\hline
Objectif atteint (accord trouvé) & oui, avec négociation & accord partiel & pas d'accord \\
\hline
\end{tabularx}
\end{center}

\begin{experience}[Tour de rôle en classe]
\begin{enumerate}
\item \textbf{Préparation (5 min)} : formez des binômes, distribuez les cartes de rôle (vendeur/client ou directeur/candidat). Chacun prépare ses phrases en silence.
\item \textbf{Performance (3 min)} : deux binômes jouent devant la classe ; les autres binômes jouent en îlots simultanés, puis échangent les rôles.
\item \textbf{Évaluation (2 min)} : les observateurs remplissent la grille et donnent un point fort et un axe d'amélioration à chaque binôme.
\item \textbf{Feedback collectif} : le professeur note au tableau les formules les mieux réussies et corrige les erreurs de tons les plus fréquentes.
\end{enumerate}
\end{experience}

\begin{exoresolu}[Construire une réplique de négociation]
Énoncé : vous êtes le client au marché de Mokolo. Le vendeur annonce : \zh{咖啡三十块一包。} Vous n'avez que quarante yuans pour tout acheter. Écrivez une réplique complète en chinois (avec pinyin et traduction) qui réagit au prix et propose un prix inférieur.
\tcblower
Solution détaillée : la réplique doit contenir deux mouvements : (1) réagir avec \zh{太……了}, (2) proposer avec \zh{怎么样} ou \zh{行不行}.\\
Réponse modèle : \zh{太贵了！我只有四十块，还要买香蕉。咖啡二十五块，行不行？}\\
\emph{Tài guì le! Wǒ zhǐyǒu sìshí kuài, hái yào mǎi xiāngjiāo. Kāfēi èrshíwǔ kuài, xíng bu xíng?}\\
« C'est trop cher ! Je n'ai que quarante yuans et je dois encore acheter des bananes. Le café à vingt-cinq yuans, ça marche ? »\\
Vérification : tons corrects sur \zh{贵} (guì, 4\up{e} ton) et \zh{买} (mǎi, 3\up{e} ton) ; le prix est proposé sans verbe « coûter » ; la question finale \zh{行不行} est une formule fonctionnelle de relance.
\end{exoresolu}

\begin{exercice}[À vous de jouer]
\begin{enumerate}
\item En binôme, rejouez la scène du marché en changeant les produits : mangues (\zh{芒果} mángguǒ) à cinq yuans le jin et arachides (\zh{花生} huāshēng) à douze yuans le paquet. Le client a vingt yuans.
\item Préparez trois questions supplémentaires pour l'entretien d'embauche (par exemple sur les horaires ou le salaire).
\item Notez pendant le jeu de vos camarades : combien de formules fonctionnelles ont-ils utilisées ? Quels tons ont posé problème ?
\end{enumerate}
\end{exercice}

\begin{aretenir}[L'essentiel de la section]
\begin{itemize}
\item Un jeu de rôle se déroule en trois temps : \zh{准备} (préparation), \zh{表演} (performance), \zh{评价} (évaluation).
\item Au marché, on négocie avec : \zh{多少钱？} — \zh{太贵了！} — \zh{……块，怎么样？}
\item En entretien, on parle de son expérience avec la particule \zh{过} placée après le verbe : \zh{我工作过两年。}
\item On parle fort, on regarde son partenaire, on n'utilise pas ses notes, et on place au moins trois formules fonctionnelles.
\end{itemize}
\end{aretenir}

\coursec{Débat guidé : le commerce Cameroun-Chine}

Après les jeux de rôle au marché et à l'entreprise, où vous avez appris à négocier, à demander le prix et à présenter un produit, nous passons à l'activité orale la plus exigeante de cette leçon : le \cle{débat guidé}. Ici, il ne s'agit plus seulement de parler chinois, mais de \textbf{défendre une opinion} en chinois, devant toute la classe. C'est exactement le type d'épreuve orale que vous rencontrerez au baccalauréat : prendre position, argumenter, écouter les autres et conclure. Le sujet choisi est un vrai sujet de société, que les Camerounais discutent chaque jour au marché, au lycée ou à la maison.

\courssub{Le sujet du débat}

\begin{definition}[Le sujet]
\zh{中国产品对喀麦隆好不好？}\\
\emph{Zhōngguó chǎnpǐn duì Kāmàilóng hǎo bu hǎo ?}\\
« Les produits chinois sont-ils bons pour le Cameroun ? »
\end{definition}

Observons la structure de cette question avant de débattre. Elle utilise deux outils grammaticaux que vous connaissez déjà :

\begin{itemize}
\item \zh{对} \emph{duì} + nom : « pour, envers ». Structure : \textbf{对 + personne/pays + adjectif}. Exemple : \zh{茶对身体好。} \emph{Chá duì shēntǐ hǎo.} « Le thé est bon pour la santé. »
\item La question affirmative-négative \textbf{adjectif + 不 + adjectif} : \zh{好不好} « est-ce bon ou pas ? », comme \zh{贵不贵} « est-ce cher ? », \zh{多不多} « y en a-t-il beaucoup ? ».
\end{itemize}

La classe se divise en deux équipes :

\begin{center}
\begin{tabularx}{\linewidth}{|X|X|X|X|}
\hline
\rowcolor{popblueL} Caractères & Pinyin & Nature & Traduction \\
\hline
赞成 & zànchéng & verbe / équipe & être pour, approuver \\
\hline
反对 & fǎnduì & verbe / équipe & être contre, s'opposer \\
\hline
中国产品 & Zhōngguó chǎnpǐn & nom & les produits chinois \\
\hline
便宜 & piányi & adjectif & bon marché \\
\hline
工作机会 & gōngzuò jīhuì & nom & les emplois, les opportunités de travail \\
\hline
基础设施 & jīchǔ shèshī & nom & les infrastructures \\
\hline
质量 & zhìliàng & nom & la qualité \\
\hline
本地产品 & běndì chǎnpǐn & nom & les produits locaux \\
\hline
竞争 & jìngzhēng & nom / verbe & la concurrence, être en concurrence \\
\hline
应该 & yīnggāi & verbe modal & devoir, il faut \\
\hline
\end{tabularx}
\end{center}

\courssub{Observer d'abord : un mini-résumé de débat}

Avant de prendre la parole, lisons comment un élève de Première peut \textbf{résumer un débat en trois phrases}. C'est le modèle à imiter pour l'examen.

\begin{textebox}[Modèle de résumé de débat]
有的同学认为中国产品很便宜，对喀麦隆人好。另一些同学认为我们应该买本地产品。我觉得两个都有道理。\\
\emph{Yǒu de tóngxué rènwéi Zhōngguó chǎnpǐn hěn piányi, duì Kāmàilóng rén hǎo. Lìng yìxiē tóngxué rènwéi wǒmen yīnggāi mǎi běndì chǎnpǐn. Wǒ juéde liǎng gè dōu yǒu dàolǐ.}\\
« Certains élèves pensent que les produits chinois sont bon marché et bons pour les Camerounais. D'autres élèves pensent que nous devrions acheter des produits locaux. Je pense que les deux ont raison. »
\end{textebox}

Ce petit texte contient la \textbf{structure d'or} du résumé de débat, en trois temps :

\begin{enumerate}
\item \zh{有的同学认为…} \emph{Yǒu de tóngxué rènwéi…} « Certains élèves pensent que… » (premier camp) ;
\item \zh{另一些同学认为…} \emph{Lìng yìxiē tóngxué rènwéi…} « D'autres élèves pensent que… » (deuxième camp) ;
\item \zh{我觉得…} \emph{Wǒ juéde…} « Je pense que… » (ton avis personnel, obligatoire à la fin).
\end{enumerate}

\begin{attention}[有的 … 另一些 …]
\zh{有的} \emph{yǒu de} signifie « certains » et se place \textbf{devant le nom} : \zh{有的同学} « certains élèves ». Ne dites jamais \zh{同学有的认为}. Pour le camp opposé, on utilise \zh{另一些} \emph{lìng yìxiē} « d'autres ». Cette paire \zh{有的…另一些…} est la manière la plus naturelle de présenter deux opinions opposées en chinois.
\end{attention}

\courssub{Grammaire du débat : le modal 应该}

Dans le modèle ci-dessus apparaît \zh{我们应该买本地产品}. Observons : où se place \zh{应该} ? Devant le verbe \zh{买}, directement, sans mot entre les deux.

\begin{propriete}[Le verbe modal 应该 yīnggāi « devoir, il faut »]
\textbf{Structure : Sujet + 应该 + verbe (+ complément).}\\
\zh{应该} exprime le devoir moral ou le conseil (« il faudrait »), comme les modaux \zh{会}, \zh{能}, \zh{想} : il se place \textbf{avant le verbe principal} et n'est jamais suivi de \zh{的}.\\
\zh{我们应该买本地产品。} \emph{Wǒmen yīnggāi mǎi běndì chǎnpǐn.} « Nous devrions acheter des produits locaux. »\\
\textbf{Négation :} Sujet + 不应该 + verbe. \zh{我们不应该只买便宜的东西。} \emph{Wǒmen bù yīnggāi zhǐ mǎi piányi de dōngxi.} « Nous ne devrions pas acheter seulement des choses bon marché. »\\
\textbf{Question :} \zh{我们应该怎么办？} \emph{Wǒmen yīnggāi zěnme bàn ?} « Que devrions-nous faire ? »
\end{propriete}

Comparaison avec le français : le français conjugue « devoir » (je dois, nous devrions) ; le chinois ne conjugue rien : \zh{应该} reste invariable pour toutes les personnes.

\begin{center}
\begin{tabularx}{\linewidth}{|X|X|X|}
\hline
\rowcolor{popblueL} Structure & Exemple (caractères + pinyin) & Traduction \\
\hline
Sujet + 应该 + verbe & 我应该学习汉语。Wǒ yīnggāi xuéxí Hànyǔ. & Je devrais étudier le chinois. \\
\hline
Sujet + 应该 + verbe + complément & 政府应该修更多的路。Zhèngfǔ yīnggāi xiū gèng duō de lù. & Le gouvernement devrait construire plus de routes. \\
\hline
Sujet + 不应该 + verbe & 我们不应该忘记本地产品。Wǒmen bù yīnggāi wàngjì běndì chǎnpǐn. & Nous ne devrions pas oublier les produits locaux. \\
\hline
应该 + verbe + 吗 ? & 我们应该买中国产品吗？Wǒmen yīnggāi mǎi Zhōngguó chǎnpǐn ma ? & Devrions-nous acheter des produits chinois ? \\
\hline
\end{tabularx}
\end{center}

\courssub{Les arguments des deux équipes}

Chaque équipe prépare ses arguments avec le vocabulaire de la leçon. Voici la carte mentale du débat :

\begin{popfigure}
\begin{tikzpicture}[
  central/.style={rectangle, rounded corners, draw=popdark, fill=popblueL, align=center, font=\bfseries, minimum width=3.2cm, minimum height=1.1cm},
  camp/.style={rectangle, rounded corners, draw=popdark, align=center, font=\bfseries, minimum width=2.6cm},
  arg/.style={rectangle, rounded corners, draw=popline, fill=popcream, align=center, minimum width=3.4cm, font=\small},
  lien/.style={-{Stealth[length=2.5mm]}, thick, popmuted}
]
\node[central, align=center] (c) at (0,0){中国产品\\ \emph{\small Zhōngguó chǎnpǐn}};
\node[camp, fill=popgreenL, align=center] (pour) at (-5.2,0){赞成\\ \emph{\small zànchéng (pour)}};
\node[camp, fill=poppinkL, align=center] (contre) at (5.2,0){反对\\ \emph{\small fǎnduì (contre)}};
\node[arg, align=center] (a1) at (-5.2,1.8){便宜 piányi\\ bon marché};
\node[arg, align=center] (a2) at (-5.2,3.3){工作机会 gōngzuò jīhuì\\ les emplois};
\node[arg, align=center] (a3) at (-5.2,4.8){基础设施 jīchǔ shèshī\\ routes, ponts, écoles};
\node[arg, align=center] (b1) at (5.2,1.8){质量 zhìliàng\\ la qualité};
\node[arg, align=center] (b2) at (5.2,3.3){本地产品 běndì chǎnpǐn\\ les produits locaux};
\node[arg, align=center] (b3) at (5.2,4.8){竞争 jìngzhēng\\ la concurrence};
\draw[lien] (pour) -- (c);
\draw[lien] (contre) -- (c);
\draw[lien, popgreen] (a1) -- (pour);
\draw[lien, popgreen] (a2) -- (pour);
\draw[lien, popgreen] (a3) -- (pour);
\draw[lien, poppink] (b1) -- (contre);
\draw[lien, poppink] (b2) -- (contre);
\draw[lien, poppink] (b3) -- (contre);
\end{tikzpicture}
\legende{Carte mentale du débat : les arguments des deux équipes autour du thème 中国产品.}
\end{popfigure}

\begin{exemplebox}[Arguments « pour » (赞成)]
\zh{我觉得中国产品很便宜，大家都可以买。} \emph{Wǒ juéde Zhōngguó chǎnpǐn hěn piányi, dàjiā dōu kěyǐ mǎi.} « Je pense que les produits chinois sont bon marché, tout le monde peut les acheter. »\\
\zh{中国公司给喀麦隆人很多工作机会。} \emph{Zhōngguó gōngsī gěi Kāmàilóng rén hěn duō gōngzuò jīhuì.} « Les entreprises chinoises donnent beaucoup d'emplois aux Camerounais. »\\
\zh{他们修了很多路、桥和学校。} \emph{Tāmen xiū le hěn duō lù, qiáo hé xuéxiào.} « Ils ont construit beaucoup de routes, de ponts et d'écoles. »
\end{exemplebox}

\begin{exemplebox}[Arguments « contre » (反对)]
\zh{我认为有些中国产品质量不太好。} \emph{Wǒ rènwéi yǒuxiē Zhōngguó chǎnpǐn zhìliàng bú tài hǎo.} « Je pense que la qualité de certains produits chinois n'est pas très bonne. »\\
\zh{我们应该买本地产品，帮助喀麦隆的工人。} \emph{Wǒmen yīnggāi mǎi běndì chǎnpǐn, bāngzhù Kāmàilóng de gōngrén.} « Nous devrions acheter des produits locaux pour aider les travailleurs camerounais. »\\
\zh{中国产品和本地产品有竞争。} \emph{Zhōngguó chǎnpǐn hé běndì chǎnpǐn yǒu jìngzhēng.} « Les produits chinois et les produits locaux sont en concurrence. »
\end{exemplebox}

\begin{attention}[质量 zhìliàng : deux quatrièmes tons]
\zh{质量} se prononce \emph{zhìliàng} : \textbf{zhì} (4\textsuperscript{ème} ton) + \textbf{liàng} (4\textsuperscript{ème} ton), deux tons descendants nets et secs. Les francophones ont tendance à adoucir le deuxième ton, ce qui donne \emph{zhìliáng} — incompréhensible. Entraînez-vous en frappant les deux syllabes comme deux coups : \emph{zhì ! liàng !} Même vigilance pour \zh{反对} \emph{fǎnduì} (3\textsuperscript{ème} + 4\textsuperscript{ème}) et \zh{竞争} \emph{jìngzhēng} (4\textsuperscript{ème} + 1\textsuperscript{er}).
\end{attention}

\courssub{Déroulé du débat et prise de parole}

\begin{methode}[Comment prendre la parole dans un débat en chinois]
\begin{enumerate}
\item \textbf{Demander la parole} : lever la main et dire \zh{我想说一点。} \emph{Wǒ xiǎng shuō yìdiǎn.} « Je voudrais dire une chose. »
\item \textbf{Annoncer son avis} : commencer obligatoirement par \zh{我觉得…} \emph{Wǒ juéde…} ou \zh{我认为…} \emph{Wǒ rènwéi…} « Je pense que… ».
\item \textbf{Donner un argument} : une phrase claire avec le vocabulaire du débat (\zh{因为…所以…} si possible : \zh{因为中国产品很便宜，所以对穷人好。}).
\item \textbf{Parler au moins 30 secondes} : deux ou trois phrases suffisent, à condition d'être complètes et bien prononcées.
\item \textbf{Terminer par une relance} : \zh{你呢？} \emph{Nǐ ne ?} « Et toi ? » ou \zh{你同意吗？} \emph{Nǐ tóngyì ma ?} « Es-tu d'accord ? »
\end{enumerate}
\end{methode}

L'enseignant joue le rôle d'\cle{arbitre} : il distribue la parole et relance le débat avec les formules vues dans la boîte « Méthode » ci-dessus : \zh{谁想说？} « Qui veut parler ? », \zh{你同意他的看法吗？} « Es-tu d'accord avec son avis ? », \zh{反对的同学，你们怎么想？} « Équipe contre, qu'en pensez-vous ? ».

À la fin du débat, la classe procède au \textbf{vote final en chinois} : chaque élève se lève et déclare \zh{我赞成。} \emph{Wǒ zànchéng.} « Je suis pour. » ou \zh{我反对。} \emph{Wǒ fǎnduì.} « Je suis contre. » L'enseignant compte les voix au tableau : \zh{十五个同学赞成，十个同学反对。} « Quinze élèves sont pour, dix sont contre. »

\begin{experience}[Le débat en classe]
\begin{enumerate}
\item Former deux équipes : \zh{赞成队} et \zh{反对队}. Chaque équipe prépare cinq arguments en dix minutes.
\item Chaque élève parle au moins 30 secondes, en commençant par \zh{我觉得} ou \zh{我认为}.
\item L'équipe adverse peut répondre : \zh{我不同意，因为…} \emph{Wǒ bù tóngyì, yīnwèi…} « Je ne suis pas d'accord, parce que… ».
\item Vote final de la classe, puis résumé collectif en trois phrases au tableau.
\end{enumerate}
\end{experience}

\begin{exoresolu}[Résumer le débat en trois phrases]
Énoncé : après le débat, rédigez un résumé en trois phrases avec \zh{有的同学认为…}, \zh{另一些同学认为…} et \zh{我觉得…}.
\tcblower
Solution détaillée :\\
1. Premier camp : \zh{有的同学认为中国产品给喀麦隆带来很多工作机会。} \emph{Yǒu de tóngxué rènwéi Zhōngguó chǎnpǐn gěi Kāmàilóng dàilái hěn duō gōngzuò jīhuì.} « Certains élèves pensent que les produits chinois apportent beaucoup d'emplois au Cameroun. »\\
2. Deuxième camp : \zh{另一些同学认为我们应该先买本地产品。} \emph{Lìng yìxiē tóngxué rènwéi wǒmen yīnggāi xiān mǎi běndì chǎnpǐn.} « D'autres élèves pensent que nous devrions d'abord acheter des produits locaux. »\\
3. Avis personnel : \zh{我觉得我们应该买质量好的产品，不管是中国的还是喀麦隆的。} \emph{Wǒ juéde wǒmen yīnggāi mǎi zhìliàng hǎo de chǎnpǐn, bùguǎn shì Zhōngguó de háishì Kāmàilóng de.} « Je pense que nous devrions acheter des produits de bonne qualité, qu'ils soient chinois ou camerounais. »\\
Méthode : une phrase par camp, puis ton avis personnel avec une nuance — c'est exactement ce que l'examinateur attend à l'oral.
\end{exoresolu}

\begin{attention}[Pièges fréquents des francophones au débat]
\begin{itemize}
\item Ne traduisez pas « je suis d'accord » par \zh{我是同意} : dites \zh{我同意} \emph{Wǒ tóngyì.}, sans \zh{是}.
\item \zh{觉得} et \zh{认为} se placent \textbf{avant} l'opinion : \zh{我觉得中国产品便宜}, jamais l'inverse.
\item Après \zh{应该}, mettez un verbe directement : pas de \zh{的}, pas de \zh{了}.
\end{itemize}
\end{attention}

\begin{savaistu}[Le Cameroun et la Chine : un vrai débat économique]
Le port en eau profonde de Kribi, le stade de Japoma à Douala et de nombreuses routes du Cameroun ont été construits avec des entreprises chinoises. Au marché de Douala comme à celui de Mokolo à Yaoundé, on trouve partout des produits chinois : téléphones, chaussures, tissus. C'est pourquoi le sujet \zh{中国产品对喀麦隆好不好？} n'est pas un exercice artificiel : c'est une question que les Camerounais se posent réellement — et un excellent sujet de réflexion pour l'épreuve orale du baccalauréat, où l'on te demandera de donner ton avis avec des arguments.
\end{savaistu}

\begin{exercice}[À vous de débattre]
\begin{enumerate}
\item Écris deux arguments « pour » et deux arguments « contre » en chinois, avec pinyin.
\item Prépare ta prise de parole de 30 secondes : \zh{我想说一点。我觉得…你呢？}
\item Résume le débat de ta classe en trois phrases avec \zh{有的同学认为…另一些同学认为…我觉得…}.
\end{enumerate}
\end{exercice}

\begin{corrige}[Exercice]
1. Pour : \zh{中国产品很便宜。} / \zh{中国公司修了很多路。} Contre : \zh{有些产品质量不好。} / \zh{我们应该帮助本地工人。}\\
2. Exemple : \zh{我想说一点。我觉得中国产品对喀麦隆好，因为他们给我们工作机会，也修了很多学校和桥。你呢？}\\
3. \zh{有的同学认为中国产品便宜又好。另一些同学认为本地产品更重要。我觉得我们应该两个都买。} — toute phrase correcte avec les trois structures est acceptée.
\end{corrige}

\newpage

\coursec{Activité d'intégration}

\coursec{Leçon 33 — Expression orale : économie}

\courssub{Objectifs de l'activité}

À l'issue de cette activité d'intégration, tu seras capable de :
\begin{itemize}
\item préparer et présenter un \cle{exposé oral simple} de 5 à 8 phrases sur une activité économique du Cameroun, en suivant le plan \zh{首先}…\zh{然后}…\zh{最后}… ;
\item utiliser les \cle{formules fonctionnelles} pour donner ton avis (\zh{我觉得…}, \zh{我同意}, \zh{你说得对，但是…}) et pour relancer la parole (\zh{你怎么看？}, \zh{为什么？}) ;
\item écouter un camarade, poser une question pertinente et réagir à son propos ;
\item prononcer correctement les \cle{tons}, en particulier les suites de troisièmes tons (règle du \emph{sandhi}) et le ton neutre ;
\item participer à un débat guidé et voter en chinois.
\end{itemize}

\courssub{Situation}

Dans le cadre de la semaine des langues vivantes, ton lycée de Yaoundé organise un \cle{forum économique des classes}. Ta classe de Première A a choisi de présenter, en chinois, les grandes activités économiques du Cameroun : le cacao, le marché, le port de Kribi, les métiers. Des élèves de Terminale et un professeur d'économie viendront écouter vos exposés. Le forum se terminera par un grand débat : \zh{中国产品对喀麦隆好不好？} (\emph{Zhōngguó chǎnpǐn duì Kāmàilóng hǎo bu hǎo ?} — « Les produits chinois sont-ils bons ou non pour le Cameroun ? »).

Pour préparer le forum, le professeur distribue le programme officiel :

\begin{textebox}[Programme du forum — 论坛节目单]
经济论坛：喀麦隆的经济\\
\emph{Jīngjì lùntán : Kāmàilóng de jīngjì}\\
« Forum économique : l'économie du Cameroun »\\[4pt]
一、每个同学介绍一个经济活动（五到八句话）。\\
\emph{Měi ge tóngxué jièshào yí ge jīngjì huódòng (wǔ dào bā jù huà).}\\
« Chaque élève présente une activité économique (cinq à huit phrases). »\\
二、听众提问题，说自己的意见。\\
\emph{Tīngzhòng tí wèntí, shuō zìjǐ de yìjian.}\\
« Les auditeurs posent des questions et donnent leur avis. »\\
三、大辩论：中国产品对喀麦隆好不好？最后投票。\\
\emph{Dà biànlùn : Zhōngguó chǎnpǐn duì Kāmàilóng hǎo bu hǎo ? Zuìhòu tóupiào.}\\
« Grand débat : les produits chinois sont-ils bons pour le Cameroun ? Vote final. »
\end{textebox}

\courssub{Consignes}

\begin{enumerate}
\item \textbf{Compréhension.} Lis le programme du forum ci-dessus. Réponds en français : combien de phrases doit contenir chaque exposé ? Que font les auditeurs ? Sur quoi porte le débat final ?
\item \textbf{Choix du sujet.} Choisis (seul ou en binôme) une activité économique du Cameroun dans la liste : le cacao, le marché de Mokolo, le port de Kribi, le football, un métier (chauffeur de taxi, vendeuse, agriculteur…). Écris ton exposé de 5 à 8 phrases en suivant le plan : \zh{首先} (d'abord : présentation), \zh{然后} (ensuite : description), \zh{最后} (enfin : ton avis).
\item \textbf{Préparation orale.} Relis ton exposé à voix haute trois fois. Souligne les troisièmes tons et vérifie le ton neutre de \zh{我们}, \zh{什么}, \zh{怎么样}. Demande au professeur de corriger ta prononciation.
\item \textbf{Présentation.} Présente ton exposé devant la classe, sans lire mot à mot : commence par \zh{大家好！今天我介绍…} et termine par \zh{谢谢大家！}.
\item \textbf{Écoute active.} Pendant les exposés de tes camarades, note une question (\zh{为什么…？}, \zh{…多少钱？}) et un avis (\zh{我同意。} ou \zh{你说得对，但是…}). Chaque auditeur doit intervenir au moins une fois.
\item \textbf{Grand débat.} La classe se divise en deux camps : \zh{好} (pour) et \zh{不好} (contre). Chaque élève donne un argument avec \zh{我觉得…，因为…}. Le débat se termine par un vote : chacun dit \zh{我投好。} ou \zh{我投不好。} (\emph{Wǒ tóu hǎo / bù hǎo} — « je vote pour / contre »).
\item \textbf{Évaluation par les pairs.} Pour chaque intervenant, remplis la grille : tons (2 points), formules fonctionnelles (2 points), fluidité (2 points), participation au débat (2 points). Justifie oralement une note en chinois : \zh{他的声调很好。} (\emph{Tā de shēngdiào hěn hǎo.} — « Ses tons sont très bons. »).
\end{enumerate}

\courssub{Ressources à mobiliser}

\begin{aretenir}[Plan de l'exposé et formules fonctionnelles]
\begin{itemize}
\item \textbf{Présenter :} \zh{大家好！今天我介绍喀麦隆的可可。} (\emph{Dàjiā hǎo ! Jīntiān wǒ jièshào Kāmàilóng de kěkě.} — « Bonjour à tous ! Aujourd'hui je présente le cacao du Cameroun. »)
\item \textbf{Ordonner :} \zh{首先}…，\zh{然后}…，\zh{最后}… (\emph{shǒuxiān}…, \emph{ránhòu}…, \emph{zuìhòu}…).
\item \textbf{Donner son avis :} \zh{我觉得…} (\emph{Wǒ juéde…} — « je pense que… ») ; \zh{我同意。} (\emph{Wǒ tóngyì.} — « je suis d'accord ») ; \zh{你说得对，但是…} (\emph{Nǐ shuō de duì, dànshì…} — « tu as raison, mais… »).
\item \textbf{Relancer :} \zh{你怎么看？} (\emph{Nǐ zěnme kàn ?} — « qu'en penses-tu ? ») ; \zh{为什么？} (\emph{Wèishénme ?} — « pourquoi ? »).
\item \textbf{Conclure :} \zh{谢谢大家！} (\emph{Xièxie dàjiā !} — « merci à tous ! »).
\end{itemize}
\end{aretenir}

\begin{center}
\begin{tabularx}{\linewidth}{|X|X|X|X|}
\hline
\rowcolor{popblueL} Caractères & Pinyin & Nature & Traduction \\
\hline
经济 & jīngjì & nom & économie \\
\hline
论坛 & lùntán & nom & forum \\
\hline
介绍 & jièshào & verbe & présenter, introduire \\
\hline
可可 & kěkě & nom & cacao \\
\hline
市场 & shìchǎng & nom & marché \\
\hline
港口 & gǎngkǒu & nom & port \\
\hline
产品 & chǎnpǐn & nom & produit \\
\hline
意见 & yìjian & nom & avis, opinion \\
\hline
同意 & tóngyì & verbe & être d'accord \\
\hline
辩论 & biànlùn & nom / verbe & débat, débattre \\
\hline
投票 & tóupiào & verbe & voter \\
\hline
因为 & yīnwei & conjonction & parce que \\
\hline
\end{tabularx}
\end{center}

\courssub{Dialogues modèles : situations de communication}

\begin{exemplebox}[Dialogue 1 : Au marché Mokolo — 莫科洛市场买东西]
\textbf{A :} 这个多少钱一斤？\\
\emph{Zhège duōshao qián yī jīn ?}\\
« Combien coûte ceci le jin (500 g) ? »\\[4pt]
\textbf{B :} 十块钱一斤。\\
\emph{Shí kuài qián yī jīn.}\\
« Dix kuai le jin. »\\[4pt]
\textbf{A :} 太贵了，便宜一点儿可以吗？\\
\emph{Tài guì le, piányi yīdiǎnr kěyǐ ma ?}\\
« C'est trop cher, peut-on baisser un peu ? »\\[4pt]
\textbf{B :} 八块钱，不能再便宜了。\\
\emph{Bā kuài qián, bù néng zài piányi le.}\\
« Huit kuai, impossible de baisser plus. »\\[4pt]
\textbf{A :} 好，买两斤。\\
\emph{Hǎo, mǎi liǎng jīn.}\\
« D'accord, j'en achète deux jin. »
\end{exemplebox}

\begin{exemplebox}[Dialogue 2 : À l'entreprise (Port de Kribi) — 在港口公司洽谈]
\textbf{A :} 你好，我想了解一下贵公司的产品。\\
\emph{Nǐ hǎo, wǒ xiǎng liǎojiě yíxià guì gōngsī de chǎnpǐn.}\\
« Bonjour, je voudrais me renseigner sur les produits de votre entreprise. »\\[4pt]
\textbf{B :} 好的，我们主要出口木材和可可。\\
\emph{Hǎo de, wǒmen zhǔyào chūkǒu mùcái hé kěkě.}\\
« D'accord, nous exportons principalement du bois et du cacao. »\\[4pt]
\textbf{A :} 你们有合作伙伴在中国吗？\\
\emph{Nǐmen yǒu hézuǒ huǒbàn zài Zhōngguó ma ?}\\
« Avez-vous des partenaires en Chine ? »\\[4pt]
\textbf{B :} 有，我们和中国的港口合作很多。\\
\emph{Yǒu, wǒmen hé Zhōngguó de gǎngkǒu hézuò hěn duō.}\\
« Oui, nous collaborons beaucoup avec les ports chinois. »
\end{exemplebox}

\courssub{Les tons à travailler}

\begin{attention}[Règles de sandhi et tons neutres — points de vigilance pour le forum]
\begin{itemize}
\item \zh{我同意} (\emph{wǒ tóngyì}) : \zh{我} (3e ton) devant \zh{同} (2e ton) se prononce en \cle{demi-troisième ton} (bas, sans remontée) : \emph{wǒ} $\to$ \emph{wǒ̆} (notation linguistique) / perception : ton bas court.
\item \zh{很好} (\emph{hěn hǎo}) : deux 3e tons consécutifs $\to$ le premier devient 2e ton : \emph{hén hǎo}. Règle générale : 3e + 3e = 2e + 3e.
\item \zh{我们} (\emph{wǒmen}), \zh{什么} (\emph{shénme}), \zh{怎么} (\emph{zěnme}) : la deuxième syllabe (\zh{们}, \zh{么}, \zh{么}) est au \cle{ton neutre} (léger, court, sans marque de ton), la première garde son ton propre.
\item \zh{为什么} (\emph{wèishénme}) : \zh{为} 4e ton ferme (\emph{wèi}), \zh{什} 2e ton (\emph{shén}), \zh{么} ton neutre (\emph{me}).
\item \zh{不} (\emph{bù}) : devient 2e ton (\emph{bú}) devant un 4e ton (\zh{太} \emph{tài}, \zh{是} \emph{shì}, \zh{好} \emph{hào} sens « être bon pour »). Ex : \zh{不太多} \emph{bú tài duō}, \zh{不好} \emph{bú hǎo}.
\end{itemize}
\end{attention}

\courssub{Jeux de rôle guidés}

\begin{experience}[Jeu de rôle 1 : Acheteur / Vendeur au marché (Mokolo / Buea)]
\begin{enumerate}
\item \textbf{Rôles :} Élève A = client (achète cacao, bananes, tissu) ; Élève B = vendeur (donne prix, négocie).
\item \textbf{Contraintes linguistiques :} Utiliser \zh{多少钱}, \zh{太贵了}, \zh{便宜一点}, \zh{可以 / 不可以}, \zh{买}, \zh{斤} (classificateur de poids).
\item \textbf{Déroulement :} 2 min de préparation (notes autorisées en pinyin), 2 min de jeu devant la classe. La classe note le vocabulaire économique entendu.
\end{enumerate}
\end{experience}

\begin{experience}[Jeu de rôle 2 : Responsable portuaire / Investisseur chinois (Port de Kribi)]
\begin{enumerate}
\item \textbf{Rôles :} Élève A = directeur du port de Kribi (présente infrastructures, volume) ; Élève B = représentant d'une entreprise chinoise (pose questions, donne avis).
\item \textbf{Contraintes linguistiques :} Utiliser \zh{港口}, \zh{出口}, \zh{进口}, \zh{合作}, \zh{重要}, \zh{因为…所以…}, \zh{我觉得…}.
\item \textbf{Déroulement :} 3 min de préparation, 3 min de jeu. L'auditoire vote pour la meilleure argumentation en chinois (\zh{我投A / 我投B}).
\end{enumerate}
\end{experience}

\courssub{Éléments socioculturels : Cameroun — Chine}

\begin{savaistu}[Relations économiques Cameroun-Chine : quelques repères factuels]
\begin{itemize}
\item La Chine est le \textbf{premier partenaire commercial} du Cameroun depuis le début des années 2010 (source : Douanes camerounaises, rapports annuels).
\item Le \textbf{port en eau profonde de Kribi} (phase 1) a été construit par l'entreprise chinoise \emph{China Harbour Engineering Company} (CHEC) ; il est géré en partie par \emph{Kribi Conteneur Terminal} (KCT), filiale de \emph{China Merchants Port Holdings}.
\item Principales \textbf{exportations camerounaises vers la Chine} : bois (grumes, sciages), coton, caoutchouc, pétrole brut, minerais (fer, bauxite).
\item Principales \textbf{importations camerounaises depuis la Chine} : machines, équipements électriques, textiles, chaussures, produits manufacturés, véhicules (motos, camions).
\item \textbf{Coopération agricole :} centres de démonstration de riz hybride (variétés chinoises adaptées au Sahel), formation d'agronomes camerounais en Chine (boursiers CSC).
\item \textbf{Débat public au Cameroun :} on entend souvent \zh{中国产品便宜，质量怎么样？} (« Produits chinois bon marché, quelle qualité ? ») et \zh{中国投资带来路桥，但本地就业够不够？} (« Investissements chinois = routes/ponts, mais emploi local suffisant ? »). Le forum vise à structurer ce débat en chinois.
\end{itemize}
\end{savaistu}

\courssub{Proposition de production et corrigé commenté}

\begin{exoresolu}[Exposé modèle : le cacao du Cameroun]
Présente un exposé de 5 à 8 phrases sur une activité économique du Cameroun, puis réponds à une question d'un auditeur.
\tcblower
\textbf{Exposé modèle :}\\[4pt]
大家好！今天我介绍喀麦隆的可可。\\
\emph{Dàjiā hǎo ! Jīntiān wǒ jièshào Kāmàilóng de kěkě.}\\
« Bonjour à tous ! Aujourd'hui, je présente le cacao du Cameroun. »\\
首先，喀麦隆有很多农民种可可。\\
\emph{Shǒuxiān, Kāmàilóng yǒu hěn duō nóngmín zhòng kěkě.}\\
« D'abord, au Cameroun, beaucoup de paysans cultivent le cacao. »\\
然后，他们把可可卖给公司。\\
\emph{Ránhòu, tāmen bǎ kěkě mài gěi gōngsī.}\\
« Ensuite, ils vendent le cacao aux entreprises. »\\
可可可以换很多钱，所以它对喀麦隆的经济很重要。\\
\emph{Kěkě kěyǐ huàn hěn duō qián, suǒyǐ tā duì Kāmàilóng de jīngjì hěn zhòngyào.}\\
« Le cacao rapporte beaucoup d'argent, donc il est très important pour l'économie du Cameroun. »\\
最后，我觉得农民的工作很辛苦，但是很有用。\\
\emph{Zuìhòu, wǒ juéde nóngmín de gōngzuò hěn xīnkǔ, dànshì hěn yǒuyòng.}\\
« Enfin, je pense que le travail des paysans est dur, mais très utile. »\\
谢谢大家！\\
\emph{Xièxie dàjiā !}\\
« Merci à tous ! »\\[6pt]
\textbf{Échange avec un auditeur :}\\[4pt]
— 你怎么看？可可农民有钱吗？\\
\emph{Nǐ zěnme kàn ? Kěkě nóngmín yǒu qián ma ?}\\
« Qu'en penses-tu ? Les planteurs de cacao sont-ils riches ? »\\
— 你说得对，这是一个问题。我觉得他们工作很多，但是钱不太多。\\
\emph{Nǐ shuō de duì, zhè shì yí ge wèntí. Wǒ juéde tāmen gōngzuò hěn duō, dànshì qián bú tài duō.}\\
« Tu as raison, c'est un vrai problème. Je pense qu'ils travaillent beaucoup, mais qu'ils ne gagnent pas beaucoup. »\\[6pt]
\textbf{Commentaire des choix de langue :}
\begin{itemize}
\item Le plan \zh{首先}…\zh{然后}…\zh{最后}… structure clairement l'exposé : c'est la colonne vertébrale exigée par la consigne.
\item \zh{把可可卖给公司} : la construction en \zh{把} met en avant l'objet (\zh{可可}) ; au niveau Première, on peut aussi dire simplement \zh{他们卖可可给公司} (SVO + 给).
\item \zh{所以} (« donc ») relie la cause à la conséquence, comme la corrélation \zh{因为}…\zh{所以}… vue en grammaire.
\item \zh{对…很重要} : structure \cle{A 对 B 很 + adjectif}, indispensable pour parler d'impact économique (\zh{可可对喀麦隆很重要}).
\item Dans la réponse, \zh{你说得对，但是…} montre qu'on écoute avant de nuancer : formule de réaction polie à réutiliser dans le débat.
\item Attention à \zh{不太多} : \zh{不} devient deuxième ton (\emph{bú}) devant le quatrième ton de \zh{太} (règle du \emph{sandhi} sur \zh{不}).
\end{itemize}
\end{exoresolu}

\courssub{Bilan de l'activité}

Ce forum économique t'a permis de mobiliser en situation réelle tous les savoir-faire de la leçon : tu as \cle{présenté} un exposé structuré avec \zh{首先}、\zh{然后}、\zh{最后}, \cle{donné ton avis} avec \zh{我觉得} et \zh{我同意}, \cle{réagi et relancé} avec \zh{你怎么看？} et \zh{你说得对，但是…}, en soignant tes tons (troisièmes tons consécutifs, ton neutre, \emph{sandhi} de \zh{不}). L'évaluation par les pairs t'a aidé à identifier tes points forts et tes progrès à faire en prononciation. Pour aller plus loin, réécoute-toi (enregistre ton exposé au téléphone) et repère les tons à retravailler ; tu pourras réutiliser toutes ces formules dans le prochain débat du module, sur le monde du travail.

\newpage

\coursec{Méthodes et exercices résolus}

\courssub{Vocabulaire de base : l'économie}

\begin{center}
\begin{tabularx}{\linewidth}{|X|X|X|X|}
\hline
\rowcolor{popblueL} \textbf{Caractères} & \textbf{Pinyin} & \textbf{Nature} & \textbf{Traduction} \\ \hline
经济 & jīngjì & n. & économie \\ \hline
农业 & nóngyè & n. & agriculture \\ \hline
工业 & gōngyè & n. & industrie \\ \hline
贸易 & màoyì & n. & commerce, échanges commerciaux \\ \hline
市场 & shìchǎng & n. & marché \\ \hline
价格 & jiàgé & n. & prix \\ \hline
出口 & chūkǒu & v./n. & exporter / exportation \\ \hline
进口 & jìnkǒu & v./n. & importer / importation \\ \hline
发展 & fāzhǎn & v. & développer, se développer \\ \hline
投资 & tóuzī & v./n. & investir / investissement \\ \hline
企业 & qǐyè & n. & entreprise \\ \hline
产品 & chǎnpǐn & n. & produit \\ \hline
本地 & běndì & adj. & local \\ \hline
外资 & wàizī & n. & capitaux étrangers \\ \hline
基础设施 & jīchǔ shèshī & n. & infrastructures \\ \hline
法郎 & fǎláng & n. & franc (CFA) \\ \hline
公斤 & gōngjīn & n. & kilogramme \\ \hline
买 & mǎi & v. & acheter \\ \hline
卖 & mài & v. & vendre \\ \hline
贵 & guì & adj. & cher \\ \hline
便宜 & piányi & adj. & bon marché \\ \hline
\end{tabularx}
\end{center}

\begin{popfigure}
\begin{tikzpicture}[
  mindmap, grow cyclic, every node/.style=concept, concept color=popblue,
  level 1/.append style={level distance=4.5cm, sibling angle=90, concept color=popblue!70},
  level 2/.append style={level distance=3cm, sibling angle=45, concept color=poporange!70},
  text=white, font=\small
]
\node {Exposé oral}
  child { node[align=center]{1. Saluer \& Annoncer\\ \zh{大家好，我介绍一下……}} }
  child { node[align=center]{2. Structurer\\ \zh{第一……第二……第三……}} }
  child { node[align=center]{3. Illustrer\\ \zh{比如……比如说……}} }
  child { node[align=center]{4. Conclure\\ \zh{介绍完了，谢谢大家！}} };
\end{tikzpicture}
\legende{Structure type d'un exposé oral simple (4 étapes)}
\end{popfigure}

\courssub{Fiche méthode 1 — Présenter un exposé simple en chinois}

\begin{methode}[Réussir un exposé oral simple sur l'économie]
Un exposé oral en chinois se construit en quatre étapes claires :
\begin{enumerate}
\item \textbf{Saluer et annoncer le sujet.} Commence par une formule de politesse puis annonce ton thème avec \zh{我给大家介绍一下……} (\emph{Wǒ gěi dàjiā jièshào yíxià…}, « je vais vous présenter… ») ou \zh{今天我要谈的题目是……} (\emph{Jīntiān wǒ yào tán de tímù shì…}, « le sujet dont je vais parler aujourd'hui est… »).
\item \textbf{Développer en trois points.} Annonce ta structure : \zh{第一……第二……第三……} (\emph{dì-yī… dì-èr… dì-sān…}, « premièrement… deuxièmement… troisièmement… »). Une idée = une phrase courte : Sujet + Verbe + Complément.
\item \textbf{Illustrer avec un exemple concret.} Utilise \zh{比如……} (\emph{bǐrú…}, « par exemple… ») ou \zh{比如说……} : \zh{比如说，杜阿拉是一个很重要的港口城市。} (\emph{Bǐrú shuō, Dù'ālā shì yí ge hěn zhòngyào de gǎngkǒu chéngshì.}, « Par exemple, Douala est une ville portuaire très importante. »)
\item \textbf{Conclure et remercier.} Termine par \zh{我的介绍完了，谢谢大家！} (\emph{Wǒ de jièshào wán le, xièxie dàjiā!}, « Ma présentation est terminée, merci à tous ! »)
\end{enumerate}
\textbf{Réflexes :} phrases courtes (8 à 12 caractères), débit lent, tons bien marqués, regard vers le public. Prépare une fiche avec 5 mots-clés seulement, jamais un texte lu mot à mot.
\end{methode}

\begin{exoresolu}[Exposé : l'économie du Cameroun]
\textbf{Énoncé.} Tu dois présenter un exposé de 5 phrases minimum sur le thème « L'économie du Cameroun ». Utilise au moins trois des mots suivants : \zh{经济} (économie), \zh{农业} (agriculture), \zh{市场} (marché), \zh{发展} (développement), \zh{出口} (exportation). Rédige ton exposé complet avec pinyin et traduction.
\tcblower
\textbf{Raisonnement.} On suit la démarche en 4 étapes : salutation + annonce du sujet, développement en points numérotés, exemple concret (cacao, pétrole), conclusion. On vérifie que chaque phrase respecte l'ordre Sujet + Verbe + Complément et que les mots imposés sont bien placés.

\textbf{Réponse rédigée :}

\zh{大家好！今天我给大家介绍一下喀麦隆的经济。}\\
\emph{Dàjiā hǎo! Jīntiān wǒ gěi dàjiā jièshào yíxià Kāmàilóng de jīngjì.}\\
« Bonjour à tous ! Aujourd'hui, je vais vous présenter l'économie du Cameroun. »

\zh{第一，农业非常重要：很多喀麦隆人在田里工作。}\\
\emph{Dì-yī, nóngyè fēicháng zhòngyào: hěn duō Kāmàilóng rén zài tián lǐ gōngzuò.}\\
« Premièrement, l'agriculture est très importante : beaucoup de Camerounais travaillent dans les champs. »

\zh{第二，喀麦隆出口可可、咖啡和石油。}\\
\emph{Dì-èr, Kāmàilóng chūkǒu kěkě, kāfēi hé shíyóu.}\\
« Deuxièmement, le Cameroun exporte le cacao, le café et le pétrole. »

\zh{第三，城市里的市场越来越大，经济正在发展。}\\
\emph{Dì-sān, chéngshì lǐ de shìchǎng yuè lái yuè dà, jīngjì zhèngzài fāzhǎn.}\\
« Troisièmement, les marchés des villes deviennent de plus en plus grands, l'économie est en train de se développer. »

\zh{我的介绍完了，谢谢大家！}\\
\emph{Wǒ de jièshào wán le, xièxie dàjiā!}\\
« Ma présentation est terminée, merci à tous ! »

\textbf{Conclusion.} L'exposé contient les 4 étapes, 3 mots imposés (\zh{经济}, \zh{农业}, \zh{出口}, \zh{市场}, \zh{发展}), des phrases courtes et la structure \zh{越来越} (\emph{yuè lái yuè}, « de plus en plus ») qui valorise la copie.
\end{exoresolu}

\courssub{Fiche méthode 2 — Donner son avis et réagir dans un débat}

\begin{methode}[Exprimer son opinion et réagir à celle des autres]
\begin{enumerate}
\item \textbf{Donner son avis} : \zh{我觉得……} (\emph{Wǒ juéde…}, « je pense que… »), \zh{我认为……} (\emph{Wǒ rènwéi…}, « je considère que… », plus soutenu), \zh{在我看来，……} (\emph{Zài wǒ kànlái,…}, « à mon avis… »).
\item \textbf{Être d'accord} : \zh{我同意你的看法。} (\emph{Wǒ tóngyì nǐ de kànfǎ.}, « je suis d'accord avec ton point de vue. »), \zh{你说得对！} (\emph{Nǐ shuō de duì!}, « tu as raison ! »).
\item \textbf{N'être pas d'accord (avec politesse)} : \zh{我不太同意，因为……} (\emph{Wǒ bú tài tóngyì, yīnwèi…}, « je ne suis pas tout à fait d'accord, parce que… »). Toujours justifier avec \zh{因为……所以……} (\emph{yīnwèi… suǒyǐ…}, « parce que… donc… »).
\item \textbf{Relancer le débat} : \zh{你呢？你怎么看？} (\emph{Nǐ ne? Nǐ zěnme kàn?}, « et toi ? qu'en penses-tu ? »), \zh{谁还有别的意见？} (\emph{Shéi hái yǒu bié de yìjiàn?}, « qui a encore un autre avis ? »).
\end{enumerate}
\textbf{Réflexe :} en chinois, on adoucit toujours le désaccord avec \zh{不太} (« pas tout à fait ») plutôt qu'un refus sec, pour rester poli.
\end{methode}

\begin{exoresolu}[Mini-débat : faut-il acheter des produits locaux ?]
\textbf{Énoncé.} Complète ce dialogue avec les formules fonctionnelles adaptées, puis traduis les répliques 2 et 4 en français.\\
A : \zh{我觉得我们应该多买喀麦隆的产品。你呢？}\\
B : (1) \zh{…………} (exprimer l'accord + une raison : les produits locaux sont frais et moins chers)\\
A : \zh{可是中国产品也很便宜，质量也不错。}\\
B : (2) \zh{…………} (exprimer un désaccord poli + relancer la discussion vers un troisième camarade)
\tcblower
\textbf{Raisonnement.} Réplique (1) : accord → \zh{我同意} ou \zh{你说得对}, puis justification avec \zh{因为}. Réplique (2) : désaccord poli → \zh{我不太同意} + raison, puis relance → \zh{你怎么看？}

\textbf{Réponses rédigées :}

(1) \zh{你说得对！因为本地产品又新鲜又便宜，所以我同意你的看法。}\\
\emph{Nǐ shuō de duì! Yīnwèi běndì chǎnpǐn yòu xīnxiān yòu piányi, suǒyǐ wǒ tóngyì nǐ de kànfǎ.}\\
Traduction : « Tu as raison ! Parce que les produits locaux sont à la fois frais et bon marché, je suis donc d'accord avec ton point de vue. »

(2) \zh{我不太同意。虽然中国产品便宜，但是买本地产品可以帮助我们的农民。小王，你怎么看？}\\
\emph{Wǒ bú tài tóngyì. Suīrán Zhōngguó chǎnpǐn piányi, dànshì mǎi běndì chǎnpǐn kěyǐ bāngzhù wǒmen de nóngmín. Xiǎo Wáng, nǐ zěnme kàn?}\\
Traduction : « Je ne suis pas tout à fait d'accord. Bien que les produits chinois soient bon marché, acheter des produits locaux peut aider nos paysans. Petit Wang, qu'en penses-tu ? »

\textbf{Conclusion.} La structure \zh{虽然……但是……} (\emph{suīrán… dànshì…}, « bien que… mais… ») permet de nuancer : c'est la marque d'un débat de bon niveau.
\end{exoresolu}

\courssub{Fiche méthode 3 — Maîtriser les tons à l'oral}

\begin{methode}[Travailler les tons pour être compris]
\begin{enumerate}
\item \textbf{Identifier les tons critiques du lexique économique.} Attention aux paires proches : \zh{mǎi} (\zh{买}, 3e ton, « acheter ») / \zh{mài} (\zh{卖}, 4e ton, « vendre ») ; \zh{jià} (\zh{价}, 4e ton, « prix ») / \zh{jiā} (\zh{家}, 1er ton, « famille »).
\item \textbf{Appliquer la règle du 3e ton (sandhi).} Deux 3e tons qui se suivent : le premier devient un 2e ton à l'oral : \zh{nǐ hǎo} se prononce \emph{ní hǎo}. Devant un autre ton (1er, 2e, 4e), le 3e ton est prononcé « à moitié » (bas, sans remonter).
\item \textbf{Attention à \zh{不} (bù) et \zh{一} (yī).} \zh{不} devient \emph{bú} (2e ton) devant un 4e ton : \zh{不对} se dit \emph{bú duì}. \zh{一} devient \emph{yí} (2e ton) devant un 4e ton (\zh{一个} → \emph{yí ge}) et \emph{yì} (4e ton) devant un 1er, 2e ou 3e ton (\zh{一公斤} → \emph{yì gōngjīn}).
\item \textbf{S'entraîner par paires minimales.} Répète à voix haute : \zh{mā} (妈, mère), \zh{má} (麻, chanvre), \zh{mǎ} (马, cheval), \zh{mà} (骂, gronder), \zh{ma} (吗, particule). Puis enregistre-toi et compare avec le modèle du professeur.
\end{enumerate}
\textbf{Réflexe :} un mauvais ton peut changer le sens : dire \zh{wǒ mài} au lieu de \zh{wǒ mǎi} transforme « j'achète » en « je vends » !
\end{methode}

\begin{exoresolu}[Dictée de tons et paires minimales]
\textbf{Énoncé.} 1) Indique le ton correct de chaque syllabe : \zh{买} (mǎi/maí), \zh{卖} (mái/mài), \zh{贵} (gùi/guì), \zh{便宜} (piányi/piānyí). 2) Comment se prononce \zh{很好} à l'oral et pourquoi ? 3) Comment se prononce \zh{不要} et pourquoi ?
\tcblower
\textbf{Raisonnement et réponses.}

1) \zh{买} = \emph{mǎi} (3e ton, « acheter ») ; \zh{卖} = \emph{mài} (4e ton, « vendre ») ; \zh{贵} = \emph{guì} (4e ton, « cher ») ; \zh{便宜} = \emph{piányi} (2e ton + ton neutre, « bon marché »).

2) \zh{很好} s'écrit \emph{hěn hǎo} mais se prononce \emph{hén hǎo} : deux 3e tons consécutifs, le premier passe au 2e ton (règle de sandhi du 3e ton).

3) \zh{不要} s'écrit \emph{bù yào} mais se prononce \emph{bú yào} : \zh{不} passe au 2e ton devant un 4e ton (\zh{要}, yào).

\textbf{Conclusion.} À l'écrit, on garde le pinyin original (\emph{hěn hǎo}, \emph{bù yào}) ; c'est seulement la prononciation qui change. À l'examen oral, appliquer ces règles montre une vraie maîtrise.
\end{exoresolu}

\courssub{Fiche méthode 4 — Jeu de rôle : négocier au marché}

\begin{methode}[Conduire une négociation commerciale simple]
\begin{enumerate}
\item \textbf{Le client demande le prix} : \zh{这个多少钱？} (\emph{Zhège duōshao qián?}, « combien coûte ceci ? »), \zh{一公斤多少钱？} (\emph{Yì gōngjīn duōshao qián?}, « combien le kilo ? »).
\item \textbf{Le vendeur annonce le prix} : \zh{一公斤五百法郎。} (\emph{Yì gōngjīn wǔbǎi fǎláng.}, « 500 francs le kilo. »). Structure : nombre + classificateur + \zh{块/法郎}.
\item \textbf{Le client négocie} : \zh{太贵了！便宜一点儿吧。} (\emph{Tài guì le! Piányi yìdiǎnr ba.}, « C'est trop cher ! Un peu moins cher, s'il vous plaît. »)
\item \textbf{Le vendeur propose} : \zh{那四百五吧。} (\emph{Nà sìbǎi wǔ ba.}, « Alors 450, d'accord. »)
\item \textbf{Conclure l'achat} : \zh{好，我要两公斤。} (\emph{Hǎo, wǒ yào liǎng gōngjīn.}, « D'accord, j'en prends deux kilos. ») — attention : \zh{两} (liǎng) et non \zh{二} (èr) devant un classificateur !
\end{enumerate}
\textbf{Réflexe :} au marché camerounais comme en Chine, négocier est naturel et attendu ; la particule finale \zh{吧} adoucit la demande.
\end{methode}

\begin{exoresolu}[Jeu de rôle : au marché Mokolo de Yaoundé]
\textbf{Énoncé.} Rédige un dialogue de 6 répliques entre un client et une vendeuse de tomates au marché Mokolo : demande du prix, annonce du prix (600 FCFA le kilo), négociation, contre-proposition (500 FCFA), accord, achat de 3 kilos. Donne le pinyin de chaque réplique.
\tcblower
\textbf{Raisonnement.} On suit les 5 étapes de la négociation. On vérifie : \zh{三公斤} avec \zh{三} (3, pas de problème), \zh{太贵了} avec \zh{了} de changement de situation, et la question de prix avec \zh{多少钱}.

\textbf{Dialogue rédigé :}

A (client) : \zh{你好！西红柿一公斤多少钱？}\\
\emph{Nǐ hǎo! Xīhóngshì yì gōngjīn duōshao qián?}\\
« Bonjour ! Combien coûte le kilo de tomates ? »

B (vendeuse) : \zh{一公斤六百法郎。}\\
\emph{Yì gōngjīn liùbǎi fǎláng.}\\
« 600 francs le kilo. »

A : \zh{太贵了！便宜一点儿吧。}\\
\emph{Tài guì le! Piányi yìdiǎnr ba.}\\
« C'est trop cher ! Faites un peu moins cher, s'il vous plaît. »

B : \zh{那五百吧，这是最低价。}\\
\emph{Nà wǔbǎi ba, zhè shì zuì dī jià.}\\
« Alors 500, c'est le prix le plus bas. »

A : \zh{好，我要三公斤。}\\
\emph{Hǎo, wǒ yào sān gōngjīn.}\\
« D'accord, j'en prends trois kilos. »

B : \zh{一共一千五百法郎。谢谢，欢迎再来！}\\
\emph{Yígòng yìqiān wǔbǎi fǎláng. Xièxie, huānyíng zài lái!}\\
« Cela fait mille cinq cents francs en tout. Merci, revenez bientôt ! »

\textbf{Conclusion.} Le dialogue respecte la progression naturelle d'une négociation et utilise \zh{一共} (\emph{yígòng}, « en tout ») et \zh{最} (\emph{zuì}, « le plus »), deux mots-outils appréciés à l'examen.
\end{exoresolu}

\courssub{Fiche méthode 5 — Jeu de rôle : dans une entreprise (entretien / réunion)}

\begin{methode}[Communiquer en contexte professionnel simple]
\begin{enumerate}
\item \textbf{Se présenter / Présenter un collègue} : \zh{我是……，负责……} (\emph{Wǒ shì…, fùzé…,} « Je suis…, responsable de… »), \zh{这是我们的经理，王先生。} (\emph{Zhè shì wǒmen de jīnglǐ, Wáng xiānsheng.}, « Voici notre directeur, M. Wang. »).
\item \textbf{Évoquer un projet / chiffre d'affaires} : \zh{我们的项目进展顺利。} (\emph{Wǒmen de xiàngmù jìnzhǎn shùnlì.}, « Notre projet progresse bien. »), \zh{今年的贸易额增加了。} (\emph{Jīnnián de màoyì'é zēngjiā le.}, « Le volume des échanges de cette année a augmenté. »).
\item \textbf{Exprimer une difficulté / proposer une solution} : \zh{有个问题：物流太慢。} (\emph{Yǒu ge wèntí: wùliú tài màn.}, « Il y a un problème : la logistique est trop lente. »), \zh{建议我们换一家运输公司。} (\emph{Jiànyì wǒmen huàn yì jiā yùnshū gōngsī.}, « Je suggère qu'on change de société de transport. »).
\item \textbf{Conclure la réunion} : \zh{今天的会议开完了，谢谢大家。} (\emph{Jīntiān de huìyì kāi wán le, xièxie dàjiā.}, « La réunion d'aujourd'hui est terminée, merci à tous. »).
\end{enumerate}
\textbf{Vocabulaire clé :} \zh{会议} (huìyì, réunion), \zh{合同} (hétong, contrat), \zh{合作} (hézuò, coopération), \zh{利润} (lìrùn, profit/bénéfice), \zh{亏损} (kuīsǔn, perte), \zh{签字} (qiānzì, signer).
\end{methode}

\begin{exoresolu}[Simulation : réunion d'équipe à Douala]
\textbf{Énoncé.} Rédige une courte intervention (4 phrases) pour une réunion : annonce que le projet de construction de route avance bien (grâce à l'entreprise chinoise), signale un problème de délai pour les matériaux, propose une solution, termine poliment.
\tcblower
\textbf{Raisonnement.} Structure : situation positive (\zh{进展顺利}) → problème (\zh{但是/可是}) → proposition (\zh{建议}) → conclusion polie.

\textbf{Intervention rédigée :}

\zh{我们的修路项目进展顺利，中国公司帮了大忙。}\\
\emph{Wǒmen de xiūlù xiàngmù jìnzhǎn shùnlì, Zhōngguó gōngsī bāng le dàmáng.}\\
« Notre projet de construction de route progresse bien, l'entreprise chinoise a beaucoup aidé. »

\zh{但是钢材的交货时间有点晚。}\\
\emph{Dànshì gāngcái de jiāohuò shíjiān yǒudiǎn wǎn.}\\
« Mais le délai de livraison de l'acier est un peu tardif. »

\zh{建议我们提前两周下订单。}\\
\emph{Jiànyì wǒmen tíqián liǎng zhōu xià dìngdān.}\\
« Je suggère qu'on passe commande deux semaines à l'avance. »

\zh{好的，今天就说这么多，谢谢大家。}\\
\emph{Hǎo de, jīntiān jiù shuō zhème duō, xièxie dàjiā.}\\
« Bon, on s'arrête là pour aujourd'hui, merci à tous. »

\textbf{Conclusion.} L'usage de \zh{提前} (\emph{tíqián}, « en avance ») et \zh{下订单} (\emph{xià dìngdān}, « passer commande ») montre un vocabulaire professionnel précis.
\end{exoresolu}

\courssub{Fiche méthode 6 — Animer un débat guidé}

\begin{popfigure}
\begin{tikzpicture}[
  node distance=1.2cm and 0.5cm,
  box/.style={draw=popblue, thick, rounded corners, fill=popblueL, text width=3.2cm, align=center, minimum height=1.2cm},
  arrow/.style={-Stealth, thick, popdark}
]
\node[box, align=center] (prep){Préparer\\2 arguments Pour / Contre\\\zh{首先、而且、但是}};
\node[box, right=of prep, align=center] (speak){Prendre la parole\\\zh{我想说两句}};
\node[box, right=of speak, align=center] (struct){Structurer\\\zh{我认为…因为…比如…}};
\node[box, right=of struct, align=center] (react){Réagir / Reformuler\\\zh{你的意思是…对吗？}};
\node[box, right=of react, align=center] (concl){Conclure\\\zh{总的来说…}};
\draw[arrow] (prep) -- (speak);
\draw[arrow] (speak) -- (struct);
\draw[arrow] (struct) -- (react);
\draw[arrow] (react) -- (concl);
\end{tikzpicture}
\legende{Flux d'une intervention en débat guidé}
\end{popfigure}

\begin{methode}[Participer à un débat guidé en classe]
\begin{enumerate}
\item \textbf{Préparer ses arguments} : note 2 arguments « pour » et 2 arguments « contre » avec des connecteurs : \zh{首先} (\emph{shǒuxiān}, « d'abord »), \zh{然后} (\emph{ránhòu}, « ensuite »), \zh{而且} (\emph{érqiě}, « de plus »), \zh{但是} (\emph{dànshì}, « mais »).
\item \textbf{Prendre la parole} : \zh{我想说两句。} (\emph{Wǒ xiǎng shuō liǎng jù.}, « je voudrais dire deux mots. »)
\item \textbf{Structurer son intervention} : opinion (\zh{我认为…}) + raison (\zh{因为…}) + exemple (\zh{比如…}).
\item \textbf{Réagir aux autres} : reformule avant de répondre : \zh{你的意思是……，对吗？} (\emph{Nǐ de yìsi shì…, duì ma?}, « tu veux dire que…, c'est ça ? »)
\item \textbf{Conclure le débat} : \zh{总的来说，……} (\emph{Zǒng de lái shuō,…}, « en résumé,… »)
\end{enumerate}
\end{methode}

\begin{exoresolu}[Débat : le commerce Cameroun-Chine est-il bon pour le Cameroun ?]
\textbf{Énoncé.} Rédige une intervention de débat de 4 phrases : donne ton avis sur le commerce entre le Cameroun et la Chine, avec une raison et un exemple, puis relance un camarade. Utilise \zh{我认为}, \zh{因为}, \zh{比如}, \zh{你怎么看}.
\tcblower
\textbf{Raisonnement.} Structure imposée : opinion → raison → exemple → relance. On choisit une position nuancée (positive mais avec une condition), ce qui montre l'esprit critique.

\textbf{Intervention rédigée :}

\zh{我认为喀麦隆和中国的贸易对我们国家有好处。}\\
\emph{Wǒ rènwéi Kāmàilóng hé Zhōngguó de màoyì duì wǒmen guójiā yǒu hǎochù.}\\
« Je considère que le commerce entre le Cameroun et la Chine est bénéfique pour notre pays. »

\zh{因为中国公司帮助我们修了很多路。}\\
\emph{Yīnwèi Zhōngguó gōngsī bāngzhù wǒmen xiū le hěn duō lù.}\\
« Parce que les entreprises chinoises nous ont aidés à construire beaucoup de routes. »

\zh{比如说，从雅温得到杜阿拉的路现在好走多了。}\\
\emph{Bǐrú shuō, cóng Yǎwēndé dào Dù'ālā de lù xiànzài hǎo zǒu duō le.}\\
« Par exemple, la route de Yaoundé à Douala est maintenant bien plus praticable. »

\zh{但是我们也应该发展自己的工业。小李，你怎么看？}\\
\emph{Dànshì wǒmen yě yīnggāi fāzhǎn zìjǐ de gōngyè. Xiǎo Lǐ, nǐ zěnme kàn?}\\
« Mais nous devons aussi développer notre propre industrie. Petit Li, qu'en penses-tu ? »

\textbf{Conclusion.} L'intervention est complète : avis clair, justification avec \zh{因为}, exemple local précis avec \zh{比如}, nuance avec \zh{但是}, et relance finale. Note l'usage de \zh{多了} (« bien plus ») après l'adjectif, structure typique du comparatif de degré.
\end{exoresolu}

\courssub{Activité de classe : Mini-débat et exposés}

\begin{experience}[Débat guidé : « Acheter local ou importer ? »]
\textbf{Objectif :} Produire un échange oral spontané de 5 minutes en petits groupes (4 élèves).
\textbf{Déroulement :}
\begin{enumerate}
\item \textbf{Préparation (5 min) :} Chaque groupe tire un sujet (ex: « Le café camerounais vs café importé », « Les téléphones chinois vs marques occidentales », « Investir dans l'agriculture vs l'industrie »). Chaque élève note 2 arguments (Pour/Contre) sur une fiche avec mots-clés en chinois.
\item \textbf{Débat (10 min) :} Un élève lance : \zh{我觉得……你呢？}. Les autres répondent en utilisant \zh{我同意/我不太同意，因为……}, \zh{比如……}, \zh{你怎么看？}. L'enseignant circule, note les erreurs de tons (surtout \zh{买/卖}, \zh{不/一}) et de structure (\zh{因为…所以…}).
\item \textbf{Retour collectif (5 min) :} Un groupe joue son débat devant la classe. Correction bienveillante au tableau : focus sur les tons et les connecteurs logiques.
\end{enumerate}
\textbf{Critères d'évaluation :} Fluidité (pas de lecture), justesse des tons (sandhi 3e ton, \zh{不}, \zh{一}), richesse du vocabulaire économique, politesse du désaccord (\zh{不太同意}).
\end{experience}

\courssub{Le saviez-vous ? : Relations économiques Cameroun — Chine}

\begin{savaistu}[Chine — Cameroun : un partenariat ancien et dense]
\begin{itemize}
\item \textbf{Histoire :} Les relations diplomatiques datent de 1971. La Chine est le premier partenaire commercial du Cameroun depuis le début des années 2010.
\item \textbf{Échanges :} Le Cameroun exporte principalement du \zh{木材} (mùcái, bois), \zh{棉花} (miánhuā, coton), \zh{可可} (kěkě, cacao), \zh{石油} (shíyóu, pétrole) vers la Chine. Il importe des \zh{机电产品} (jīdiàn chǎnpǐn, produits mécaniques et électroniques), \zh{纺织品} (fǎngzhīpǐn, textiles), \zh{摩托车} (mótuōchē, motos) et du \zh{riz} (mǐ).
\item \textbf{Infrastructures :} De grands chantiers camerounais sont réalisés par des entreprises chinoises : \zh{克里比深水港} (Kèlǐbǐ shēnshuǐgǎng, Port en eau profonde de Kribi), \zh{雅温得-杜阿拉高速公路} (Yǎwēndé-Dù'ālā gāosù gōnglù, Autoroute Yaoundé-Douala), barrages hydroélectriques (\zh{梅金} Mékin, \zh{洛姆潘加尔} Luòmǔpángā'ěr).
\item \textbf{Coopération agricole :} Des centres de démonstration de \zh{杂交水稻} (zájiāo shuǐdào, riz hybride) chinois existent au Cameroun pour booster le rendement rizicole.
\item \textbf{Point de vigilance :} La balance commerciale est souvent excédentaire pour la Chine. Le défi pour le Cameroun est la \zh{加工} (jiāgōng, transformation) locale des matières premières avant exportation pour créer de la \zh{ valeur ajoutée} (jiàzhí).
\end{itemize}
\end{savaistu}

\courssub{Exercices d'entraînement}

\begin{exercice}[Exercice 1 : Exposé écrit → Oral]
Prépare un exposé de 6 phrases sur le thème : \zh{我家乡的经济} (L'économie de ma ville/région d'origine). Utilise \textbf{au moins 5 mots} du tableau de vocabulaire (p. 1). Structure : Salutation → 3 points (\zh{第一/二/三}) + exemple (\zh{比如}) → Conclusion. Écris les caractères, le pinyin et la traduction.
\end{exercice}

\begin{exercice}[Exercice 2 : Compléter le dialogue (Négociation entreprise)]
Complète les répliques manquantes (B) en caractères + pinyin.
A : \zh{王经理，我们的合同准备好了吗？} (Directeur Wang, le contrat est-il prêt ?)
B : (1) \zh{…………} (Répondre : pas encore, il y a un problème de prix / \zh{价格})
A : \zh{什么问题？} (Quel problème ?)
B : (2) \zh{…………} (Expliquer : le prix des matières premières a augmenté / \zh{上涨 shàngzhǎng})
A : \zh{那我们能不能便宜一点？} (Peut-on baisser un peu ?)
B : (3) \zh{…………} (Répondre poliment : difficile, mais on peut offrir un meilleur service / \zh{服务 fúwù})
A : \zh{好，那我们再谈谈。} (Bon, on en reparle.)
\end{exercice}

\begin{exercice}[Exercice 3 : Tons et prononciation]
Pour chaque phrase, indique le pinyin \textbf{réel à l'oral} (avec sandhi) et surligne les syllabes modifiées.
1. \zh{我买三公斤苹果。} (J'achète 3 kg de pommes.)
2. \zh{不卖了，太便宜了。} (Je ne vends plus, c'est trop bon marché.)
3. \zh{一个很好的市场。} (Un très bon marché.)
4. \zh{我不想去那个工厂。} (Je ne veux pas aller dans cette usine.)
\end{exercice}

\begin{exercice}[Exercice 4 : Traduction thème (Français → Chinois)]
Traduis en chinois (caractères + pinyin).
1. L'économie du Cameroun se développe vite.
2. Je pense que les produits locaux sont moins chers que les produits importés.
3. Bien que le prix soit cher, la qualité est bonne. (Utilise \zh{虽然……但是……})
4. Combien coûte ce kilo de riz ? C'est trop cher, faites 400 francs.
5. En résumé, la coopération entre la Chine et le Cameroun est très importante.
\end{exercice}

\courssub{Corrigés des exercices}

\begin{corrige}[Exercice 1]
\textbf{Modèle de réponse (exemple pour Yaoundé) :}
\zh{大家好！今天我介绍一下雅温得的经济。}\\
\emph{Dàjiā hǎo! Jīntiān wǒ jièshào yíxià Yǎwēndé de jīngjì.}

\zh{第一，服务业很发达，很多人在公司上班。}\\
\emph{Dì-yī, fúwùyè hěn fādá, hěn duō rén zài gōngsī shàngbān.}

\zh{第二，市场很多，比如莫科洛市场很大。}\\
\emph{Dì-èr, shìchǎng hěn duō, bǐrú Mòkēluò shìchǎng hěn dà.}

\zh{第三，建筑业发展很快，到处在盖楼。}\\
\emph{Dì-sān, jiànzhúyè fāzhǎn hěn kuài, dàochù zài gài lóu.}

\zh{比如，新的商业中心马上就要开业了。}\\
\emph{Bǐrú, xīn de shāngyè zhōngxīn mǎshàng jiù yào kāiyè le.}

\zh{我的介绍完了，谢谢大家！}\\
\emph{Wǒ de jièshào wán le, xièxie dàjiā!}
\end{corrige}

\begin{corrige}[Exercice 2]
(1) B : \zh{还没准备好，价格有点问题。} (\emph{Hái méi zhǔnbèi hǎo, jiàgé yǒudiǎn wèntí.})
(2) B : \zh{原材料的价格上涨了。} (\emph{Yuáncáiliào de jiàgé shàngzhǎng le.})
(3) B : \zh{价格很难便宜，但是我们可以提供更好的服务。} (\emph{Jiàgé hěn nán piányi, dànshì wǒmen kěyǐ tígōng gèng hǎo de fúwù.})
\end{corrige}

\begin{corrige}[Exercice 3]
1. \zh{我买三公斤苹果。} → \emph{Wǒ \textbf{mǎi} sān gōngjīn píngguǒ.} (Pas de sandhi : \zh{买} 3e ton isolé ou devant 1er ton \zh{三} → 3e ton « demi » bas). \textbf{Attention :} \zh{三} (sān, 1er ton) ne déclenche pas le sandhi du 3e ton sur \zh{买}.
2. \zh{不卖了，太便宜了。} → \emph{\textbf{bú} mài le, tài piányi le.} (\zh{不} → \emph{bú} devant \zh{卖} 4e ton).
3. \zh{一个很好的市场。} → \emph{\textbf{yí} ge \textbf{hén} hǎo de shìchǎng.} (\zh{一} → \emph{yí} devant \zh{个} 4e ton ; \zh{很} \zh{好} → \emph{hén hǎo} deux 3e tons).
4. \zh{我不想去那个工厂。} → \emph{Wǒ \textbf{bú} xiǎng qù nà ge gōngchǎng.} (\zh{不} → \emph{bú} devant \zh{想} 3e ton ? \textbf{Non !} \zh{想} est 3e ton. \zh{不} garde son ton d'origine \emph{bù} (ou 2e ton « demi » bas) devant 1er, 2e, 3e ton. \textbf{Correction :} \emph{Wǒ bù xiǎng…} ou \emph{Wǒ bú xiǎng…} (usage courant variable, règle stricte : \emph{bù} devant 3e ton). On enseigne la règle stricte : \emph{bù}).
\end{corrige}

\begin{corrige}[Exercice 4]
1. \zh{喀麦隆的经济发展得很快。} (\emph{Kāmàilóng de jīngjì fāzhǎn de hěn kuài.})
2. \zh{我认为本地产品比进口产品便宜。} (\emph{Wǒ rènwéi běndì chǎnpǐn bǐ jìnkǒu chǎnpǐn piányi.})
3. \zh{虽然价格贵，但是质量好。} (\emph{Suīrán jiàgé guì, dànshì zhìliàng hǎo.})
4. \zh{大米一公斤多少钱？太贵了，四百法郎吧。} (\emph{Dàmǐ yì gōngjīn duōshao qián? Tài guì le, sìbǎi fǎláng ba.})
5. \zh{总的来说，中喀合作非常重要。} (\emph{Zǒng de lái shuō, Zhōng-Kā hézuò fēicháng zhòngyào.}) — \zh{中喀} = abréviation Chine-Cameroun.
\end{corrige}

\begin{attention}[Les 5 erreurs qui coûtent des points au Bac]
\begin{enumerate}
\item \textbf{Confondre \zh{买} (mǎi, 3e ton, acheter) et \zh{卖} (mài, 4e ton, vendre).} Erreur de ton ET de sens : à l'oral, cela inverse toute la phrase. Astuce mnémotechnique : \zh{卖} a un « chapeau » \zh{十} (shi) en plus — celui qui vend a quelque chose en plus à donner (le bien).
\item \textbf{Dire \zh{二公斤} (èr gōngjīn) au lieu de \zh{两公斤} (liǎng gōngjīn).} Devant un classificateur ou une mesure, « deux » se dit toujours \zh{两} (liǎng), jamais \zh{二} (èr). On dit \zh{两个人} (liǎng ge rén), \zh{两块钱} (liǎng kuài qián), \zh{两公斤} (liǎng gōngjīn).
\item \textbf{Calquer l'ordre des mots français (Complément de lieu).} Ne dis pas \zh{我买东西在市场} pour « j'achète des choses au marché ». En chinois, le lieu se place \textbf{AVANT} le verbe : \zh{我在市场买东西} (\emph{Wǒ zài shìchǎng mǎi dōngxi}). Règle : Sujet + Lieu + Verbe + Objet.
\item \textbf{Oublier le changement de ton de \zh{不} (bù) et \zh{一} (yī).} À l'oral : \zh{不要} se prononce \emph{bú yào} (devant 4e ton), \zh{一个} se prononce \emph{yí ge} (devant 4e ton), \zh{一公斤} se prononce \emph{yì gōngjīn} (devant 1er ton). Un examinateur oral y est très attentif.
\item \textbf{Mauvaise gestion de la particule \zh{了} (le).} \zh{太贵} (trop cher en général) vs \zh{太贵了} (constatation situationnelle : « c'est trop cher ! »). \zh{经济发展} (fait général) vs \zh{经济发展了} (changement : « l'économie s'est développée »). \zh{正在发展} (action en cours) \textbf{ne prend jamais} de \zh{了} final. Ne pas cumuler \zh{正在} et \zh{了}.
\end{enumerate}
\end{attention}

\newpage

\coursec{Exercices}

\courssub{Je vérifie mes connaissances}

\begin{exercice}[QCM : le vocabulaire de l'économie]
Pour chaque question, choisis la bonne réponse (a, b ou c).
\begin{enumerate}
\item « Le commerce » se dit : a) \zh{农业} \quad b) \zh{商业} \quad c) \zh{工厂}
\item \zh{经济} (jīngjì) signifie : a) la culture \quad b) l'économie \quad c) l'école
\item Pour demander le prix, on dit : a) \zh{太贵了！} \quad b) \zh{多少钱？} \quad c) \zh{便宜一点儿！}
\item \zh{公司} (gōngsī) désigne : a) une entreprise \quad b) un marché \quad c) un agriculteur
\end{enumerate}
\end{exercice}

\begin{exercice}[Vrai ou faux justifié]
Réponds par VRAI ou FAUX, puis justifie ta réponse en une phrase en français.
\begin{enumerate}
\item \zh{首先} (shǒuxiān) sert à conclure un exposé.
\item \zh{我觉得} (wǒ juéde) permet de donner son avis.
\item \zh{卖} (mài) signifie « acheter ».
\item \zh{最后} (zuìhòu) est un connecteur qui introduit la dernière étape d'un exposé.
\end{enumerate}
\end{exercice}

\begin{exercice}[Vocabulaire : complète les phrases]
Complète chaque phrase avec le mot qui convient, choisi dans la liste : \zh{市场}、\zh{工人}、\zh{产品}、\zh{发展}。
\begin{enumerate}
\item 中国的\zh{＿＿＿＿}发展得很快。(Zhōngguó de ...... fāzhǎn de hěn kuài.)
\item 妈妈在\zh{＿＿＿＿}买水果。(Māma zài ...... mǎi shuǐguǒ.)
\item 这家工厂有很多\zh{＿＿＿＿}。(Zhè jiā gōngchǎng yǒu hěn duō ......)
\item 喀麦隆的\zh{＿＿＿＿}很好。(Kāmàilóng de ...... hěn hǎo.)
\end{enumerate}
\end{exercice}

\begin{exercice}[Associe le pinyin aux caractères]
Relie chaque caractère ou mot à son pinyin.
\begin{center}
\begin{tabularx}{\linewidth}{|X|X|}\hline
\rowcolor{popblueL} Caractères & Pinyin \\ \hline
\zh{贵} & mǎi \\ \hline
\zh{买} & shēngyì \\ \hline
\zh{生意} & guì \\ \hline
\zh{卖} & piányi \\ \hline
\zh{便宜} & mài \\ \hline
\end{tabularx}
\end{center}
\end{exercice}

\courssub{Je m'entraîne}

\begin{exercice}[Traduction du français vers le chinois]
Traduis les phrases suivantes en chinois (caractères + pinyin).
\begin{enumerate}
\item Je pense que l'agriculture est importante.
\item Pour moi, le commerce est utile.
\item D'abord, je présente l'économie de ma région ; ensuite, je donne mon avis ; enfin, je conclus.
\item C'est trop cher ! Un peu moins cher, s'il vous plaît !
\end{enumerate}
\end{exercice}

\begin{exercice}[Dialogue à trous : au marché]
Complète ce dialogue de marchandage avec les formules fonctionnelles étudiées : \zh{太贵了}、\zh{便宜一点儿}、\zh{多少钱}、\zh{好的}。
\begin{textebox}[Au marché de Mokolo]
A：老板，这个\zh{＿＿＿＿}？(Lǎobǎn, zhège ......?)\\
B：五十块。(Wǔshí kuài.)\\
A：\zh{＿＿＿＿}！(......!) 能\zh{＿＿＿＿}吗？(Néng ...... ma?)\\
B：四十五块，怎么样？(Sìshíwǔ kuài, zěnmeyàng?)\\
A：\zh{＿＿＿＿}，我买两个。(......, wǒ mǎi liǎng ge.)
\end{textebox}
\end{exercice}

\begin{exercice}[Remise en ordre des mots]
Remets les mots dans le bon ordre pour former des phrases correctes, puis traduis-les en français.
\begin{enumerate}
\item \zh{经济 / 的 / 中国 / 很快 / 发展}
\item \zh{我 / 重要 / 觉得 / 很 / 农业}
\item \zh{喀麦隆 / 和 / 之间 / 中国 / 的 / 贸易}
\end{enumerate}
\end{exercice}

\begin{exercice}[Discrimination des tons]
Le professeur lit un mot de chaque paire minimale. Entoure le mot entendu, puis écris son pinyin avec le ton correct.
\begin{enumerate}
\item \zh{买} (mǎi) / \zh{卖} (mài)
\item \zh{贵} (guì) / \zh{归} (guī)
\item \zh{商} (shāng) / \zh{上} (shàng)
\item \zh{工} (gōng) / \zh{共} (gòng)
\end{enumerate}
\end{exercice}

\courssub{Je me prépare au Bac}

\begin{exercice}[Compréhension de l'écrit et questions de langue]
Lis le texte suivant, puis réponds aux questions en français.
\begin{textebox}[L'économie de ma région]
我住在杜阿拉。杜阿拉是喀麦隆最大的城市，也是一个很重要的港口。\\
Wǒ zhù zài Dù'ālā. Dù'ālā shì Kāmàilóng zuì dà de chéngshì, yě shì yí ge hěn zhòngyào de gǎngkǒu.\\
这里的经济很发达：有很多公司、工厂和市场。很多人在市场做生意，卖水果、鱼和衣服。\\
Zhèlǐ de jīngjì hěn fādá : yǒu hěn duō gōngsī, gōngchǎng hé shìchǎng. Hěn duō rén zài shìchǎng zuò shēngyì, mài shuǐguǒ, yú hé yīfu.\\
我觉得杜阿拉的经济会越来越发达。\\
Wǒ juéde Dù'ālā de jīngjì huì yuè lái yuè fādá.
\end{textebox}
\begin{enumerate}
\item Où habite le narrateur et pourquoi cette ville est-elle importante ?
\item Quelles activités économiques sont mentionnées dans le texte ?
\item Que vendent les gens au marché ?
\item Quelle est l'opinion du narrateur sur l'avenir de sa ville ? Relève la formule d'opinion utilisée.
\item Question de langue : relève dans le texte un superlatif avec \zh{最} et un complément de degré avec \zh{很}, puis traduis-les.
\end{enumerate}
\end{exercice}

\begin{exercice}[Thème et version]
\textbf{Thème.} Traduis en chinois (caractères + pinyin) :
\begin{enumerate}
\item Le commerce entre le Cameroun et la Chine se développe rapidement.
\item Je pense que les produits chinois ne sont pas chers.
\end{enumerate}
\textbf{Version.} Traduis en français :
\begin{enumerate}
\item \zh{首先，我介绍我们地区的农业；然后，我说说市场；最后，我谈谈我的看法。}\\
(Shǒuxiān, wǒ jièshào wǒmen dìqū de nóngyè ; ránhòu, wǒ shuōshuo shìchǎng ; zuìhòu, wǒ tántan wǒ de kànfǎ.)
\item \zh{很多年轻人想在公司工作。} (Hěn duō niánqīngrén xiǎng zài gōngsī gōngzuò.)
\end{enumerate}
\end{exercice}

\begin{exercice}[Expression orale et écrite : exposé et débat]
\textbf{Partie A — Préparation écrite.} Rédige un exposé de \textbf{6 phrases minimum} (environ 60 à 80 caractères) présentant l'activité économique d'une région du Cameroun (l'agriculture de l'Ouest, le port de Douala, le tourisme de Buea, etc.). Tu dois utiliser :
\begin{itemize}
\item les trois connecteurs \zh{首先}、\zh{然后}、\zh{最后} ;
\item au moins une formule d'opinion (\zh{我觉得}、\zh{我认为} ou \zh{对我来说}).
\end{itemize}
\textbf{Partie B — Passage à l'oral.} Présente ton exposé devant la classe en 1 minute, en soignant la prononciation des tons.\\
\textbf{Partie C — Débat.} Rédige \textbf{3 arguments} \zh{赞成} (pour) ou \zh{反对} (contre) les produits chinois au Cameroun, puis défends ta position oralement en 30 secondes. Ton évaluation portera sur la prononciation des tons et l'emploi d'au moins 4 formules fonctionnelles.
\end{exercice}

\newpage

\coursec{Corrigés des exercices}

\coursec{Leçon 33 — Expression orale : économie — Corrigés des exercices}

\begin{corrige}[Exercice 1 — QCM : le vocabulaire de l'économie]
\begin{enumerate}
\item \textbf{b) \zh{商业} (shāngyè)}. \zh{农业} (nóngyè) = l'agriculture ; \zh{工厂} (gōngchǎng) = l'usine.
\item \textbf{b) l'économie}. \zh{经济} (jīngjì) = économie.
\item \textbf{b) \zh{多少钱？} (duōshao qián?)}. \zh{太贵了！} (tài guì le) = C'est trop cher ! ; \zh{便宜一点儿！} (piányi yìdiǎnr) = Moins cher !
\item \textbf{a) une entreprise}. \zh{公司} (gōngsī) = entreprise / société ; \zh{市场} (shìchǎng) = marché ; \zh{农民} (nóngmín) = agriculteur.
\end{enumerate}
\end{corrige}

\begin{corrige}[Exercice 2 — Vrai ou faux justifié]
\begin{enumerate}
\item \textbf{FAUX}. \zh{首先} (shǒuxiān) signifie « d'abord, en premier lieu » et sert à \textbf{commencer} un exposé, non à le conclure (pour conclure : \zh{最后} zuìhòu, \zh{总之} zǒngzhī).
\item \textbf{VRAI}. \zh{我觉得} (wǒ juéde) signifie « je pense que / à mon avis / j'ai l'impression que » ; c'est la formule d'opinion la plus courante à l'oral (niveau HSK 2-3).
\item \textbf{FAUX}. \zh{卖} (mài, 4\textsuperscript{e} ton) signifie « \textbf{vendre} ». « Acheter » se dit \zh{买} (mǎi, 3\textsuperscript{e} ton). \textbf{Attention à la confusion de tons} : mǎi (acheter, ton descendant-remontant) vs mài (vendre, ton descendant).
\item \textbf{VRAI}. \zh{最后} (zuìhòu) signifie « finalement, en dernier lieu » et introduit la dernière partie ou la conclusion d'un exposé.
\end{enumerate}
\end{corrige}

\begin{corrige}[Exercice 3 — Vocabulaire : complète les phrases]
\begin{enumerate}
\item 中国的\zh{经济} (jīngjì) 发展得很快。 (L'économie de la Chine se développe très vite.)
\item 妈妈在\zh{市场} (shìchǎng) 买水果。 (Maman achète des fruits au marché.)
\item 这家工厂有很多\zh{工人} (gōngrén)。 (Cette usine a beaucoup d'ouvriers.)
\item 喀麦隆的\zh{产品} (chǎnpǐn) 很好。 (Les produits du Cameroun sont très bons.)
\end{enumerate}
\emph{Analyse grammaticale :} Dans la phrase 1, \zh{经济} est le sujet de \zh{发展} ; le complément de degré \zh{得很快} précise la manière. Dans la phrase 2, \zh{市场} est le lieu (complément de lieu placé avant le verbe grâce à \zh{在}). Dans la phrase 3, \zh{工人} est l'objet de \zh{有}. Dans la phrase 4, \zh{产品} est le sujet de la phrase adjectivale \zh{很好} (très bon).
\end{corrige}

\begin{corrige}[Exercice 4 — Associe le pinyin aux caractères]
\begin{center}
\begin{tabularx}{\linewidth}{|X|X|X|X|}\hline
\rowcolor{popblueL} Caractères & Pinyin & Nature & Traduction \\ \hline
\zh{贵} & guì & adj. & cher \\ \hline
\zh{买} & mǎi & verbe & acheter \\ \hline
\zh{生意} & shēngyì & nom & affaires, commerce \\ \hline
\zh{卖} & mài & verbe & vendre \\ \hline
\zh{便宜} & piányi & adj. & bon marché \\ \hline
\end{tabularx}
\end{center}
\end{corrige}

\begin{corrige}[Exercice 5 — Traduction du français vers le chinois]
\begin{enumerate}
\item \zh{我觉得农业很重要。} (Wǒ juéde nóngyè hěn zhòngyào.) \\ \emph{Variante :} \zh{我认为农业非常重要。} (Wǒ rènwéi nóngyè fēicháng zhòngyào.)
\item \zh{对我来说，商业很有用。} (Duì wǒ lái shuō, shāngyè hěn yǒuyòng.) \\ \emph{Variante :} \zh{我觉得商业很有用。}
\item \zh{首先，我介绍我地区的经济；然后，我谈谈我的看法；最后，我做总结。} (Shǒuxiān, wǒ jièshào wǒ dìqū de jīngjì; ránhòu, wǒ tántan wǒ de kànfǎ; zuìhòu, wǒ zuò zǒngjié.) \\ \emph{Structures imposées :} \zh{首先} (d'abord) — \zh{然后} (ensuite) — \zh{最后} (enfin). Verbes : \zh{介绍} (présenter), \zh{谈谈} (parler un peu de), \zh{做总结} (conclure / faire un résumé).
\item \zh{太贵了！便宜一点儿，谢谢！} (Tài guì le! Piányi yìdiǎnr, xièxie!) \\ \emph{Formules de marchandage :} \zh{太贵了} (c'est trop cher), \zh{便宜一点儿} (moins cher / un peu moins cher), \zh{谢谢} (merci / s'il vous plaît — poli pour clore la négociation).
\end{enumerate}
\end{corrige}

\begin{corrige}[Exercice 6 — Dialogue à trous : au marché]
\begin{textebox}[Au marché de Mokolo — Corrigé]
A：老板，这个\zh{多少钱}？(Lǎobǎn, zhège \textbf{duōshao qián}?)\\
B：五十块。(Wǔshí kuài.)\\
A：\zh{太贵了}！(\textbf{Tài guì le}!) 能\zh{便宜一点儿}吗？(Néng \textbf{piányi yìdiǎnr} ma?)\\
B：四十五块，怎么样？(Sìshíwǔ kuài, zěnmeyàng?)\\
A：\zh{好的}，我买两个。(\textbf{Hǎo de}, wǒ mǎi liǎng ge.)
\end{textebox}
\emph{Explications :} \zh{多少钱} (combien) pour demander le prix ; \zh{太贵了} (exclamation de protestation) ; \zh{能……吗？} (peut-on… ?) + \zh{便宜一点儿} (moins cher) pour négocier ; \zh{好的} (d'accord) pour accepter le prix final. Classificateur \zh{个} (gè) pour objets courants.
\end{corrige}

\begin{corrige}[Exercice 7 — Remise en ordre des mots]
\begin{enumerate}
\item \zh{中国的经济发展很快。} (Zhōngguó de jīngjì fāzhǎn hěn kuài.) — \textbf{L'économie de la Chine se développe très vite.} \\ \emph{Structure :} Sujet (\zh{中国的经济}) + Verbe (\zh{发展}) + Complément de degré (\zh{很} + \zh{快}).
\item \zh{我觉得农业很重要。} (Wǒ juéde nóngyè hěn zhòngyào.) — \textbf{Je pense que l'agriculture est très importante.} \\ \emph{Structure :} Sujet + \zh{觉得} + Proposition (\zh{农业很重要}). \zh{很} lie le sujet à l'adjectif (pas de verbe « être »).
\item \zh{喀麦隆和中国之间的贸易很发达。} (Kāmàilóng hé Zhōngguó zhījiān de màoyì hěn fādá.) — \textbf{Le commerce entre le Cameroun et la Chine est très développé.} \\ \emph{Structure :} Sujet complexe (\zh{喀麦隆和中国之间的贸易}) + \zh{很} + Adjectif (\zh{发达}). \zh{之间} (entre) structure le complément du nom \zh{贸易}.
\end{enumerate}
\end{corrige}

\begin{corrige}[Exercice 8 — Discrimination des tons]
\textbf{Corrigé du professeur (clé) :}
\begin{enumerate}
\item \zh{买} (mǎi) — 3\textsuperscript{e} ton (descendant-remontant) = « acheter ». \zh{卖} (mài) — 4\textsuperscript{e} ton (descendant) = « vendre ». \emph{Différence : seul le ton change.}
\item \zh{贵} (guì) — 4\textsuperscript{e} ton = « cher ». \zh{归} (guī) — 1\textsuperscript{e} ton (haut plat) = « retourner / rentrer ».
\item \zh{商} (shāng) — 1\textsuperscript{e} ton = « commerce » (dans \zh{商业}, \zh{商品}). \zh{上} (shàng) — 4\textsuperscript{e} ton = « haut / sur / monter / classe ».
\item \zh{工} (gōng) — 1\textsuperscript{e} ton = « travail / ouvrier / industrie ». \zh{共} (gòng) — 4\textsuperscript{e} ton = « ensemble / commun / total ».
\end{enumerate}
\emph{Conseil pédagogique :} Travailler ces paires minimales en « écho » (professeur → classe → binômes) pour fixer la distinction tonale, cruciale pour le sens.
\end{corrige}

\begin{corrige}[Exercice 9 — Compréhension de l'écrit et questions de langue]
\begin{enumerate}
\item Le narrateur habite \textbf{Douala} (\zh{杜阿拉} Dù'ālā). Cette ville est importante car c'est \textbf{la plus grande ville du Cameroun} et \textbf{un port très important}.
\item Activités mentionnées : \textbf{beaucoup d'entreprises (\zh{公司}), d'usines (\zh{工厂}) et de marchés (\zh{市场})} ; \textbf{le commerce (\zh{做生意})} au marché.
\item On y vend : \textbf{des fruits (\zh{水果}), du poisson (\zh{鱼}) et des vêtements (\zh{衣服})}.
\item L'opinion : \textbf{l'économie de Douala sera de plus en plus développée (\zh{会越来越发达})}. Formule d'opinion relevée : \zh{我觉得} (wǒ juéde) — « je pense que / à mon avis ».
\item \textbf{Superlatif avec \zh{最} :} \zh{最大} (zuì dà) — « le plus grand / la plus grande » (dans \zh{喀麦隆最大的城市}). \textbf{Complément de degré avec \zh{很} :} \zh{很重要} (hěn zhòngyào) — « très important » ; \zh{很发达} (hěn fādá) — « très développé ».
\end{enumerate}
\end{corrige}

\begin{corrige}[Exercice 10 — Thème et version]
\textbf{Thème (Français → Chinois)}
\begin{enumerate}
\item \zh{喀麦隆和中国之间的贸易发展得很快。} (Kāmàilóng hé Zhōngguó zhījiān de màoyì fāzhǎn de hěn kuài.) \\ \emph{Points clés :} \zh{之间} (entre), \zh{贸易} (commerce), \zh{发展得很快} (se développe très vite — complément de degré \zh{得} + adjectif).
\item \zh{我觉得中国产品不贵。} (Wǒ juéde Zhōngguó chǎnpǐn bù guì.) \\ \emph{Variante :} \zh{我认为中国产品不贵。} \zh{产品} (produit), \zh{不贵} (pas cher).
\end{enumerate}

\textbf{Version (Chinois → Français)}
\begin{enumerate}
\item \textbf{« D'abord, j'introduis l'agriculture de notre région ; ensuite, je parle du marché ; enfin, j'expose mon point de vue. »} \\ \emph{Analyse :} \zh{首先} (d'abord), \zh{介绍} (présenter/introduire), \zh{我们地区的农业} (l'agriculture de notre région) ; \zh{然后} (ensuite), \zh{说说} (parler un peu de / évoquer), \zh{市场} (le marché) ; \zh{最后} (enfin), \zh{谈谈} (discuter/évoquer), \zh{我的看法} (mon point de vue / mon opinion).
\item \textbf{« Beaucoup de jeunes veulent travailler dans une entreprise. »} \\ \emph{Analyse :} \zh{很多} (beaucoup de), \zh{年轻人} (jeunes), \zh{想} (vouloir), \zh{在公司} (dans une entreprise), \zh{工作} (travailler).
\end{enumerate}
\end{corrige}

\begin{corrige}[Exercice 11 — Expression orale et écrite : exposé et débat]
\textbf{Partie A — Proposition de rédaction modèle (environ 70 caractères)}
\begin{textebox}[Modèle d'exposé : L'agriculture de l'Ouest]
首先，我介绍西部地区的农业。这里土地肥沃，气候好。\\
(Shǒuxiān, wǒ jièshào Xībù dìqū de nóngyè. Zhèlǐ tǔdì féiwò, qìhòu hǎo.)\\
然后，主要产品是咖啡、可可和玉米。很多农民在田里工作。\\
(Ránhòu, zhǔyào chǎnpǐn shì kāfēi, kěkě hé yùmǐ. Hěn duō nóngmín zài tiánlǐ gōngzuò.)\\
最后，我觉得西部的农业对喀麦隆很重要，会越来越好。\\
(Zuìhòu, wǒ juéde Xībù de nóngyè duì Kāmàilóng hěn zhòngyào, huì yuè lái yuè hǎo.)
\end{textebox}
\textbf{Vérification des contraintes :}
\begin{itemize}
\item 3 phrases complexes (découpables en 6 énoncés simples pour l'oral).
\item Connecteurs : \zh{首先}, \zh{然后}, \zh{最后} ✓
\item Formule d'opinion : \zh{我觉得} ✓
\item Vocabulaire économie : \zh{农业}, \zh{土地}, \zh{产品}, \zh{农民}, \zh{重要} ✓
\end{itemize}

\textbf{Partie B — Points de vigilance pour le passage à l'oral (1 min)}
\begin{itemize}
\item \textbf{Tons :} \zh{西部} (xībù, 1-4), \zh{农业} (nóngyè, 2-4), \zh{肥沃} (féiwò, 2-4), \zh{咖啡} (kāfēi, 1-1), \zh{可可} (kěkě, 3-3 \rightarrow \textbf{kékě} par sandhi 3+3), \zh{重要} (zhòngyào, 4-4), \zh{越来越} (yuè lái yuè, 4-2-4). \textbf{Attention :} \zh{可可} se prononce \textbf{kékě} (2\textsuperscript{e} + 3\textsuperscript{e} ton) à cause de la règle du sandhi sur deux 3\textsuperscript{e} tons consécutifs (comme \zh{很好} → hén hǎo).
\item \textbf{Liaison / rythme :} Marquer une pause après les connecteurs (\zh{首先，} / \zh{然后，} / \zh{最后，}).
\item \textbf{Volume et regard :} S'adresser à la classe, ne pas lire sur la feuille.
\end{itemize}

\textbf{Partie C — Arguments types pour le débat (Produits chinois au Cameroun)}
\textbf{Position \zh{赞成} (Pour) :}
\begin{enumerate}
\item \zh{中国产品便宜，老百姓买得起。} (Produits chinois bon marché, le peuple peut se les offrir.)
\item \zh{种类多，什么都有，选择大。} (Variétés nombreuses, tout y est, grand choix.)
\item \zh{促进了喀麦隆和中国的贸易关系。} (Favorise les relations commerciales Cameroun-Chine.)
\end{enumerate}
\textbf{Position \zh{反对} (Contre) :}
\begin{enumerate}
\item \zh{有的产品质量不好，不耐用。} (Certains produits de mauvaise qualité, pas durables.)
\item \zh{冲击本地工厂和手工业，工人失业。} (Impact négatif sur usines locales et artisanat, chômage.)
\item \zh{售后服务难，坏了修不了。} (Service après-vente difficile, impossible à réparer.)
\end{enumerate}

\textbf{Formules fonctionnelles à employer (minimum 4) :}
\zh{我觉得…} / \zh{我认为…} (opinion) ; \zh{一方面……另一方面……} (argumentation équilibrée) ; \zh{因为……所以……} (cause-conséquence) ; \zh{比如说…} (exemple) ; \zh{我不同意，因为…} (désaccord) ; \zh{你说得有道理，但是…} (concession + contre-argument).

\textbf{Barème indicatif (sur 20) — Expression orale (Exposé + Débat)}
\begin{center}
\begin{tabularx}{\linewidth}{|X|c|X|X|}\hline
\rowcolor{popblueL} Critères & Points & Indicateurs \\ \hline
Prononciation (tons, intonation) & 5 & Tons justes sur vocabulaire clé, phrase fluide, sandhi 3+3 maîtrisé. \\ \hline
Vocabulaire & structures & 5 & Vocabulaire économie utilisé (\zh{经济, 贸易, 产品, 发展}…), structures \zh{首先/然后/最后}, \zh{觉得/认为} correctes. \\ \hline
Cohérence & contenu & 5 & Exposé structuré (3 parties), arguments logiques, exemples Cameroun/Chine. \\ \hline
Interaction & formules & 5 & 4+ formules fonctionnelles, réactivité dans le débat, politesse (\zh{请问}, \zh{谢谢}). \\ \hline
\end{tabularx}
\end{center}
\end{corrige}

\newpage

\coursec{Fiche bilan}

\begin{popfigure}
\begin{tikzpicture}[
    mindmap/.style={concept, execute at begin node=\hskip0pt, font=\small\bfseries},
    level 1 concept/.append style={level distance=3.5cm, sibling angle=360/7, font=\small\bfseries},
    level 2 concept/.append style={level distance=2.5cm, font=\footnotesize},
    concept color=popblue,
    branch1/.style={concept color=poporange},
    branch2/.style={concept color=popgreen},
    branch3/.style={concept color=poppurple},
    branch4/.style={concept color=poppink},
    branch5/.style={concept color=popgold},
    branch6/.style={concept color=popdark},
    branch7/.style={concept color=popblue!70!popgreen},
    every node/.style={align=center, text width=2.2cm}
]
\node[mindmap, align=center]{Économie\\经济\\ \emph{Jīngjì}}
    child[branch1] { node[align=center]{Vocabulaire\\ clé\\ 核心词汇} }
    child[branch2] { node[align=center]{Dialogues\\ modèles\\ 示范对话} }
    child[branch3] { node[align=center]{Formules\\ fonctionnelles\\ 功能语块} }
    child[branch4] { node[align=center]{Tons\\ difficiles\\ 重点声调} }
    child[branch5] { node[align=center]{Jeux de\\ rôle\\ 情景剧} }
    child[branch6] { node[align=center]{Débat\\ Cameroun-Chine\\ 中喀辩论} }
    child[branch7] { node[align=center]{Stratégies\\ orales\\ 口语技巧} };
\end{tikzpicture}
\legende{Carte mentale : Leçon 33 — Expression orale : Économie}
\end{popfigure}

\begin{bilan}

\textbf{Lexique clé (HSK 3--4) — 核心词汇}
\begin{center}
\begin{tabularx}{\linewidth}{|X|X|X|X|}
\hline
\rowcolor{popblueL} \textbf{Caractères} & \textbf{Pinyin} & \textbf{Nature} & \textbf{Traduction} \\ \hline
经济 & jīngjì & n. & économie \\ \hline
贸易 & màoyì & n. & commerce, échange commercial \\ \hline
投资 & tóuzī & v./n. & investir / investissement \\ \hline
市场 & shìchǎng & n. & marché \\ \hline
进口 / 出口 & jìnkǒu / chūkǒu & v. & importer / exporter \\ \hline
合作 & hézuò & v./n. & coopérer / coopération \\ \hline
发展 & fāzhǎn & v. & développer, se développer \\ \hline
机会 & jīhuì & n. & opportunité, occasion \\ \hline
挑战 & tiǎozhàn & n. & défi, challenge \\ \hline
优势 & yōushì & n. & avantage, atout \\ \hline
基础设施 & jīchǔ shèshī & n. & infrastructures \\ \hline
农产品 & nóngchǎnpǐn & n. & produits agricoles \\ \hline
签署 & qiānshǔ & v. & signer (accord, contrat) \\ \hline
互利共赢 & hùlì gòngyíng & id. & gagnant-gagnant, mutuellement bénéfique \\ \hline
\end{tabularx}
\end{center}

\textbf{Structures grammaticales \& Formules fonctionnelles — 语法结构与功能语块}
\begin{center}
\begin{tabularx}{\linewidth}{|X|X|X|}
\hline
\rowcolor{poporangeL} \textbf{Structure / Fonction} & \textbf{Exemple (Caractères + Pinyin)} & \textbf{Traduction} \\ \hline
\cle{Présenter un exposé} : 首先…，其次…，最后… & 首先，我介绍中喀贸易。其次，分析机会。最后，提建议。 & D'abord…, ensuite…, enfin… \\ \hline
\cle{Donner son avis} : 我认为 / 我觉得… & 我认为中国市场对喀麦隆很重要。 & Je pense que / À mon avis… \\ \hline
\cle{Exprimer accord/désaccord} : 我同意 / 我不同意… & 我同意你的看法，但是… & Je suis d'accord / Pas d'accord… \\ \hline
\cle{Relancer / Demander précision} : 请问… / 具体来说？ & 请问，具体指哪些农产品？ & Pourriez-vous préciser… ? \\ \hline
\cle{Comparer (比 bǐ)} : A 比 B + Adj. & 中国经济比喀麦隆经济发达。 & L'économie chinoise est plus développée que celle du Cameroun. \\ \hline
\cle{Indiquer but (为了 wèile)} : 为了 + Objectif, Sujet + Action & 为了加强合作，我们签署了协议。 & Pour renforcer la coopération, nous avons signé un accord. \\ \hline
\cle{Passif / Subir (被 bèi)} : Sujet + 被 + Agent + Verbe & 很多产品被出口到中国。 & Beaucoup de produits sont exportés vers la Chine. \\ \hline
\cle{Aspect accompli (了 le)} : Verbe + 了 (action terminée) & 我们已经达成协议了。 & Nous avons déjà conclu un accord. \\ \hline
\end{tabularx}
\end{center}

\textbf{Méthodes express — 口语考试技巧}
\begin{itemize}
    \item \textbf{Préparer une fiche « mots-clés »} (pas de phrases entières) : structure \zh{首先-其次-最后} + 3 arguments + conclusion.
    \item \textbf{Tons} : repérer les \cle{3e tons} consécutifs (ex: \zh{经济 jīngjì} → \emph{jíngjí}) et les \cle{bu/bù} + 4e ton (\zh{不是 bú shì}).
    \item \textbf{Connecteurs} : utilisez \zh{另外 lìngwài} (de plus), \zh{不过 bùguò} (mais), \zh{例如 lìrú} (par ex.) pour fluidifier.
    \item \textbf{Débat} : écoutez → \zh{记笔记 jì bǐjì} (mots-clés) → reformulez \zh{你的意思是… Nǐ de yìsi shì…} → donnez votre avis.
\end{itemize}

\end{bilan}

\begin{aretenir}[Checklist avant le Bac — 口语经济专项]
\begin{itemize}
    \item $\square$ Je connais le vocabulaire de base (économie, commerce, investir, importer/exporter, coopération, infrastructures).
    \item $\square$ Je sais présenter un mini-exposé structuré : \zh{首先、其次、最后} (2 min max).
    \item $\square$ Je maîtrise les formules : \zh{我认为、我同意/不同意、请问、具体来说}.
    \item $\square$ Je prononce correctement les tons des mots-clés : \zh{经济 jīngjì, 投资 tóuzī, 合作 hézuò, 互利共赢 hùlì gòngyíng}.
    \item $\square$ Je gère les enchaînements tonaux (3e ton + 3e ton → 2e ton + 3e ton ; \zh{不 bù} devant 4e ton).
    \item $\square$ Je sais utiliser \zh{比} pour comparer Cameroun/Chine (PIB, population, marché).
    \item $\square$ Je sais utiliser \zh{被} pour décrire les flux (produits exportés, accords signés).
    \item $\square$ Je peux jouer un rôle (acheteur/vendeur, entrepreneur/investisseur) avec politesse (\zh{请、谢谢、麻烦您}).
    \item $\square$ Je peux citer 2 avantages Cameroun (ressources, jeunesse) et 2 avantages Chine (marché, technologie, capital).
    \item $\square$ Je connais l'expression idiomatique \zh{互利共赢} et je peux l'employer en conclusion.
    \item $\square$ Je gère mon temps de parole : ni trop court (< 1 min), ni hors-sujet, je regarde le jury/camarades.
    \item $\square$ En cas de trou de mémoire : \zh{对不起，让我想一下… Duìbuqǐ, ràng wǒ xiǎng yíxià…} (ne pas rester silencieux).
\end{itemize}
\end{aretenir}

\findecours

\end{document}
